%% Generated by Sphinx.
\def\sphinxdocclass{report}
\documentclass[letterpaper,10pt,english]{sphinxmanual}
\ifdefined\pdfpxdimen
   \let\sphinxpxdimen\pdfpxdimen\else\newdimen\sphinxpxdimen
\fi \sphinxpxdimen=.75bp\relax
\ifdefined\pdfimageresolution
    \pdfimageresolution= \numexpr \dimexpr1in\relax/\sphinxpxdimen\relax
\fi
%% let collapsible pdf bookmarks panel have high depth per default
\PassOptionsToPackage{bookmarksdepth=5}{hyperref}

\PassOptionsToPackage{booktabs}{sphinx}
\PassOptionsToPackage{colorrows}{sphinx}

\PassOptionsToPackage{warn}{textcomp}
\usepackage[utf8]{inputenc}
\ifdefined\DeclareUnicodeCharacter
% support both utf8 and utf8x syntaxes
  \ifdefined\DeclareUnicodeCharacterAsOptional
    \def\sphinxDUC#1{\DeclareUnicodeCharacter{"#1}}
  \else
    \let\sphinxDUC\DeclareUnicodeCharacter
  \fi
  \sphinxDUC{00A0}{\nobreakspace}
  \sphinxDUC{2500}{\sphinxunichar{2500}}
  \sphinxDUC{2502}{\sphinxunichar{2502}}
  \sphinxDUC{2514}{\sphinxunichar{2514}}
  \sphinxDUC{251C}{\sphinxunichar{251C}}
  \sphinxDUC{2572}{\textbackslash}
\fi
\usepackage{cmap}
\usepackage[T1]{fontenc}
\usepackage{amsmath,amssymb,amstext}
\usepackage{babel}



\usepackage{tgtermes}
\usepackage{tgheros}
\renewcommand{\ttdefault}{txtt}



\usepackage[Bjarne]{fncychap}
\usepackage{sphinx}

\fvset{fontsize=auto}
\usepackage{geometry}


% Include hyperref last.
\usepackage{hyperref}
% Fix anchor placement for figures with captions.
\usepackage{hypcap}% it must be loaded after hyperref.
% Set up styles of URL: it should be placed after hyperref.
\urlstyle{same}

\addto\captionsenglish{\renewcommand{\contentsname}{Choose your language:}}

\usepackage{sphinxmessages}
\setcounter{tocdepth}{0}



\title{Pelagos Scientific data on KMaP}
\date{Mar 26, 2026}
\release{0.1}
\author{Pelagos Agreement}
\newcommand{\sphinxlogo}{\vbox{}}
\renewcommand{\releasename}{Release}
\makeindex
\begin{document}

\ifdefined\shorthandoff
  \ifnum\catcode`\=\string=\active\shorthandoff{=}\fi
  \ifnum\catcode`\"=\active\shorthandoff{"}\fi
\fi

\pagestyle{empty}
\sphinxmaketitle
\pagestyle{plain}
\sphinxtableofcontents
\pagestyle{normal}
\phantomsection\label{\detokenize{index::doc}}


\sphinxstepscope


\chapter{KMaP \textendash{} Knowledge Management Platform}
\label{\detokenize{en/index:kmap-knowledge-management-platform}}\label{\detokenize{en/index::doc}}
\sphinxAtStartPar
\sphinxstylestrong{User Documentation for the Pelagos Agreement Group}

\begin{sphinxadmonition}{note}{Note:}
\sphinxAtStartPar
This manual is dedicated to Kmap users that are also members of the \sphinxstylestrong{Pelagos Agreement} group
Some features described here may appear diffferently or not be visible to users outside
this group.
\end{sphinxadmonition}

\sphinxAtStartPar
Following The Permanent Secretariat of the Pelagos Agreementahs signed a Memorandun of Understanding in 2024
to publish the results of scientific calls conmsisting in
maps and geospatial data onto the UNEP\sphinxhyphen{}MAP Knowledge Management Platform (KMaP) managed by INFO/RAC;
The KMaP platform is accessible at: \sphinxurl{https://kmap.info-rac.org}

\sphinxAtStartPar
This manual describe how to access and visualize data, metadata and maps in the KMaP platform

\begin{figure}[htbp]
\centering

\noindent\sphinxincludegraphics{{call1_map}.png}
\end{figure}

\sphinxstepscope


\section{Welcome to KMaP}
\label{\detokenize{en/01_welcome:welcome-to-kmap}}\label{\detokenize{en/01_welcome:welcome-en}}\label{\detokenize{en/01_welcome::doc}}

\subsection{What is KMaP?}
\label{\detokenize{en/01_welcome:what-is-kmap}}
\sphinxAtStartPar
\sphinxstylestrong{KMaP} (Knowledge Management Platform) is the Spatial Data Infrastructure developed by
\sphinxstylestrong{INFO/RAC}, the Regional Activity Centre for Information and Communication of the
\sphinxhref{https://www.unep.org/unepmap/}{Mediterranean Action Plan} (UNEP/MAP).

\sphinxAtStartPar
KMaP is built on \sphinxhref{https://geonode.org/}{GeoNode}, an open\sphinxhyphen{}source web platform for
managing and sharing geospatial information. It acts as a unified access hub to the
environmental and spatial knowledge heritage of the Mediterranean region.

\sphinxAtStartPar
The platform is structured around three hubs:


\begin{savenotes}\sphinxattablestart
\sphinxthistablewithglobalstyle
\centering
\begin{tabular}[t]{\X{20}{100}\X{15}{100}\X{65}{100}}
\sphinxtoprule
\sphinxstyletheadfamily 
\sphinxAtStartPar
Hub
&\sphinxstyletheadfamily 
\sphinxAtStartPar
Access
&\sphinxstyletheadfamily 
\sphinxAtStartPar
Description
\\
\sphinxmidrule
\sphinxtableatstartofbodyhook
\sphinxAtStartPar
\sphinxstylestrong{Data Hub}
&
\sphinxAtStartPar
\sphinxstyleemphasis{Maps} {\color{red}\bfseries{}|mapsbut|}
&
\sphinxAtStartPar
Geospatial data: datasets, interactive maps, dashboards, and geostories.
\\
\sphinxhline
\sphinxAtStartPar
\sphinxstylestrong{Knowledge Hub}
&
\sphinxAtStartPar
\sphinxstyleemphasis{Library} {\color{red}\bfseries{}|librarybut|}
&
\sphinxAtStartPar
Documentary heritage: reports, publications, and reference documents.
\\
\sphinxhline
\sphinxAtStartPar
\sphinxstylestrong{Exchange Hub}
&
\sphinxAtStartPar
\sphinxstyleemphasis{Network} {\color{red}\bfseries{}|networkbut|}
&
\sphinxAtStartPar
Collaboration and knowledge exchange.
\\
\sphinxbottomrule
\end{tabular}
\sphinxtableafterendhook\par
\sphinxattableend\end{savenotes}

\begin{sphinxadmonition}{note}{Note:}
\sphinxAtStartPar
As a member of the \sphinxstylestrong{Pelagos Agreement} group, you have access to specific resources,
metadata editing rights, and group\sphinxhyphen{}level permissions that are not available to
general public users.
\end{sphinxadmonition}


\subsection{Resource Types}
\label{\detokenize{en/01_welcome:resource-types}}
\sphinxAtStartPar
KMaP organises all content into five resource types. Each type serves a distinct purpose.


\subsubsection{Dataset}
\label{\detokenize{en/01_welcome:dataset}}
\sphinxAtStartPar
A \sphinxstylestrong{Dataset} is a geographic data layer representing real\sphinxhyphen{}world features or measurements.
It is the foundational building block of the platform — the raw material from which maps,
dashboards, and analyses are constructed.

\sphinxAtStartPar
Datasets can be:
\begin{itemize}
\item {} 
\sphinxAtStartPar
\sphinxstylestrong{Vector} \textendash{} discrete geographic features represented as points, lines, or polygons
(e.g. coastlines, protected area boundaries, monitoring station locations).

\item {} 
\sphinxAtStartPar
\sphinxstylestrong{Raster} \textendash{} continuous surfaces represented as a grid of cells
(e.g. sea surface temperature, satellite imagery, depth models).

\end{itemize}

\sphinxAtStartPar
What you can do with a Dataset:
\begin{itemize}
\item {} 
\sphinxAtStartPar
View it interactively on a web map

\item {} 
\sphinxAtStartPar
Inspect its attribute table

\item {} 
\sphinxAtStartPar
Download it in standard formats (Shapefile, GeoTIFF, GeoJSON, KML, CSV, …)

\item {} 
\sphinxAtStartPar
Read and edit its metadata

\item {} 
\sphinxAtStartPar
Add it as a layer to a Map

\end{itemize}

\begin{sphinxadmonition}{tip}{Tip:}
\sphinxAtStartPar
Datasets published for the Pelagos Agreement are tagged with the keywords
\sphinxcode{\sphinxupquote{pelagos sanctuary}} or \sphinxcode{\sphinxupquote{pelagos agreement}}. Use these keywords to filter
the catalogue quickly (see {\hyperref[\detokenize{en/02_find_and_visualize:find-en}]{\sphinxcrossref{\DUrole{std}{\DUrole{std-ref}{How to Find and Visualise Your Data}}}}}).
\end{sphinxadmonition}


\subsubsection{Document}
\label{\detokenize{en/01_welcome:document}}
\sphinxAtStartPar
A \sphinxstylestrong{Document} is any uploaded file that complements spatial data: a PDF report,
a field photograph, a spreadsheet, a presentation, or any other reference material.

\sphinxAtStartPar
Documents are fully integrated resources in KMaP — they have their own metadata page,
can be searched and filtered, and can be linked to Datasets or Maps.

\sphinxAtStartPar
Supported file types include: \sphinxcode{\sphinxupquote{.pdf}}, \sphinxcode{\sphinxupquote{.docx}}, \sphinxcode{\sphinxupquote{.xlsx}}, \sphinxcode{\sphinxupquote{.pptx}}, \sphinxcode{\sphinxupquote{.jpg}},
\sphinxcode{\sphinxupquote{.png}}, \sphinxcode{\sphinxupquote{.tif}}, \sphinxcode{\sphinxupquote{.mp4}}, \sphinxcode{\sphinxupquote{.zip}}, and many others.


\subsubsection{Map}
\label{\detokenize{en/01_welcome:map}}
\sphinxAtStartPar
A \sphinxstylestrong{Map} is an interactive, multi\sphinxhyphen{}layer web map created by composing one or more
Datasets. Maps are built and viewed using the integrated \sphinxstylestrong{MapStore} map viewer.

\sphinxAtStartPar
Each Map:
\begin{itemize}
\item {} 
\sphinxAtStartPar
Combines multiple Dataset layers with individual styles and visibility settings

\item {} 
\sphinxAtStartPar
Includes a basemap (e.g. OpenStreetMap or satellite imagery)

\item {} 
\sphinxAtStartPar
Supports pan, zoom, feature query, measurement, and print tools

\item {} 
\sphinxAtStartPar
Can be shared via a permanent link or embedded in external websites

\end{itemize}


\subsubsection{GeoStory}
\label{\detokenize{en/01_welcome:geostory}}
\sphinxAtStartPar
A \sphinxstylestrong{GeoStory} is a scrollable, multimedia narrative that weaves together text,
interactive maps, images, and video — similar to a story map.

\sphinxAtStartPar
GeoStories are composed of sections:


\begin{savenotes}\sphinxattablestart
\sphinxthistablewithglobalstyle
\centering
\begin{tabular}[t]{\X{25}{100}\X{75}{100}}
\sphinxtoprule
\sphinxstyletheadfamily 
\sphinxAtStartPar
Section type
&\sphinxstyletheadfamily 
\sphinxAtStartPar
Description
\\
\sphinxmidrule
\sphinxtableatstartofbodyhook
\sphinxAtStartPar
\sphinxstylestrong{Title}
&
\sphinxAtStartPar
Full\sphinxhyphen{}screen opening panel with background image or video.
\\
\sphinxhline
\sphinxAtStartPar
\sphinxstylestrong{Banner}
&
\sphinxAtStartPar
Horizontal divider or chapter header with text and/or image.
\\
\sphinxhline
\sphinxAtStartPar
\sphinxstylestrong{Paragraph}
&
\sphinxAtStartPar
Free\sphinxhyphen{}form text and media (image, video, map).
\\
\sphinxhline
\sphinxAtStartPar
\sphinxstylestrong{Immersive}
&
\sphinxAtStartPar
Full\sphinxhyphen{}screen interactive map that the text scrolls over.
\\
\sphinxhline
\sphinxAtStartPar
\sphinxstylestrong{GeoCarousel}
&
\sphinxAtStartPar
Slideshow linked to map locations.
\\
\sphinxhline
\sphinxAtStartPar
\sphinxstylestrong{Web page}
&
\sphinxAtStartPar
Embedded live webpage displayed inline.
\\
\sphinxbottomrule
\end{tabular}
\sphinxtableafterendhook\par
\sphinxattableend\end{savenotes}

\sphinxAtStartPar
GeoStories are ideal for communicating environmental topics to non\sphinxhyphen{}specialist audiences.


\subsubsection{Dashboard}
\label{\detokenize{en/01_welcome:dashboard}}
\sphinxAtStartPar
A \sphinxstylestrong{Dashboard} is a configurable panel of interactive widgets — charts, tables, maps,
counters, and text blocks — linked to one or more Datasets.

\sphinxAtStartPar
Available widget types:


\begin{savenotes}\sphinxattablestart
\sphinxthistablewithglobalstyle
\centering
\begin{tabular}[t]{\X{25}{100}\X{75}{100}}
\sphinxtoprule
\sphinxstyletheadfamily 
\sphinxAtStartPar
Widget
&\sphinxstyletheadfamily 
\sphinxAtStartPar
Description
\\
\sphinxmidrule
\sphinxtableatstartofbodyhook
\sphinxAtStartPar
\sphinxstylestrong{Map}
&
\sphinxAtStartPar
Interactive mini\sphinxhyphen{}map displaying dataset layers.
\\
\sphinxhline
\sphinxAtStartPar
\sphinxstylestrong{Chart}
&
\sphinxAtStartPar
Bar, line, pie, or scatter plots derived from dataset attributes.
\\
\sphinxhline
\sphinxAtStartPar
\sphinxstylestrong{Table}
&
\sphinxAtStartPar
Sortable and filterable tabular view of dataset attributes.
\\
\sphinxhline
\sphinxAtStartPar
\sphinxstylestrong{Counter}
&
\sphinxAtStartPar
Large\sphinxhyphen{}format summary statistic (count, sum, average, …).
\\
\sphinxhline
\sphinxAtStartPar
\sphinxstylestrong{Text}
&
\sphinxAtStartPar
Free\sphinxhyphen{}text block for titles or commentary.
\\
\sphinxbottomrule
\end{tabular}
\sphinxtableafterendhook\par
\sphinxattableend\end{savenotes}

\sphinxAtStartPar
Widgets in a Dashboard are interconnected: selecting a region on the map widget
automatically filters all other widgets to display data for that region only.


\subsection{Interoperable Services}
\label{\detokenize{en/01_welcome:interoperable-services}}
\sphinxAtStartPar
KMaP exposes its data through standard \sphinxstylestrong{OGC web services}, allowing any compatible
GIS software or web application to access the data directly without downloading files.


\begin{savenotes}\sphinxattablestart
\sphinxthistablewithglobalstyle
\centering
\begin{tabular}[t]{\X{15}{100}\X{15}{100}\X{70}{100}}
\sphinxtoprule
\sphinxstyletheadfamily 
\sphinxAtStartPar
Service
&\sphinxstyletheadfamily 
\sphinxAtStartPar
Standard
&\sphinxstyletheadfamily 
\sphinxAtStartPar
Description
\\
\sphinxmidrule
\sphinxtableatstartofbodyhook
\sphinxAtStartPar
\sphinxstylestrong{WMS}
&
\sphinxAtStartPar
OGC Web Map Service
&
\sphinxAtStartPar
Returns dataset layers as rendered map images. Ideal for visualisation in
desktop GIS (QGIS, ArcGIS) or web applications.
\\
\sphinxhline
\sphinxAtStartPar
\sphinxstylestrong{WFS}
&
\sphinxAtStartPar
OGC Web Feature Service
&
\sphinxAtStartPar
Returns raw vector feature data (geometries + attributes) for download or
analysis. Supports filtering by attribute or bounding box.
\\
\sphinxhline
\sphinxAtStartPar
\sphinxstylestrong{WCS}
&
\sphinxAtStartPar
OGC Web Coverage Service
&
\sphinxAtStartPar
Returns raw raster data (grids) for analysis. Available for raster datasets.
\\
\sphinxhline
\sphinxAtStartPar
\sphinxstylestrong{WMTS}
&
\sphinxAtStartPar
OGC Web Map Tile Service
&
\sphinxAtStartPar
Returns pre\sphinxhyphen{}rendered map tiles for fast display in web maps.
\\
\sphinxhline
\sphinxAtStartPar
\sphinxstylestrong{CSW}
&
\sphinxAtStartPar
OGC Catalogue Service
&
\sphinxAtStartPar
Exposes KMaP’s metadata catalogue for search and harvest by external portals.
\\
\sphinxbottomrule
\end{tabular}
\sphinxtableafterendhook\par
\sphinxattableend\end{savenotes}
\subsubsection*{How to get service URLs}

\sphinxAtStartPar
On any Dataset detail page, click the \sphinxstylestrong{OGC Services} tab (or look for the
\sphinxstyleemphasis{Download / OGC} section). The WMS, WFS, and WCS endpoint URLs are listed there
and can be copied directly into QGIS or any other OGC\sphinxhyphen{}compatible client.

\begin{figure}[htbp]
\centering
\capstart

\noindent\sphinxincludegraphics{{ogc_services}.png}
\caption{The ogc service links in metadata section}\label{\detokenize{en/01_welcome:id1}}\end{figure}

\begin{sphinxadmonition}{tip}{Tip:}
\sphinxAtStartPar
In QGIS, go to \sphinxstylestrong{Layer \(\rightarrow\) Add Layer \(\rightarrow\) Add WMS/WMTS Layer}, paste the WMS URL,
and browse the list of available layers to add KMaP data directly to your project
without downloading any files.
Consider also the \sphinxhref{https://geonode.org/QGISGeoNodePlugin/}{QGIS Geonode Plugin}
to connect directly with kmap resources.
\end{sphinxadmonition}


\subsection{Themes Available on KMaP}
\label{\detokenize{en/01_welcome:themes-available-on-kmap}}
\sphinxAtStartPar
KMaP organises its datasets into the the following categories according with UNEP\sphinxhyphen{}MAP themes:
\begin{itemize}
\item {} 
\sphinxAtStartPar
\sphinxincludegraphics[height=3.5em]{{01pollution}.svg} \sphinxstylestrong{Pollution}

\item {} 
\sphinxAtStartPar
\sphinxincludegraphics[height=3.5em]{{02biodiversity}.svg} \sphinxstylestrong{Marine Biodiversity}

\item {} 
\sphinxAtStartPar
\sphinxincludegraphics[height=3.5em]{{03sustainability}.svg} \sphinxstylestrong{Sustainability and Blue Economy}

\item {} 
\sphinxAtStartPar
\sphinxincludegraphics[height=3.5em]{{04climate_change}.svg} \sphinxstylestrong{Climate Change}

\item {} 
\sphinxAtStartPar
\sphinxincludegraphics[height=3.5em]{{05MSP}.svg} \sphinxstylestrong{Marine Spatial Planning}

\item {} 
\sphinxAtStartPar
\sphinxincludegraphics[height=3.5em]{{06fishery}.svg} \sphinxstylestrong{Fishery and Aquaculture}

\item {} 
\sphinxAtStartPar
\sphinxincludegraphics[height=3.5em]{{07governance}.svg} \sphinxstylestrong{Governance}

\end{itemize}

\sphinxAtStartPar
Resources related to the \sphinxhref{https://pelagos-sanctuary.org/scientific-studies/}{Scientific Studies financied by Pelagos Agreement}
are categorized under dtheir respective themes, mainly Marine Biodiversity  and Pollution.

\sphinxstepscope


\section{How to Find and Visualise Your Data}
\label{\detokenize{en/02_find_and_visualize:how-to-find-and-visualise-your-data}}\label{\detokenize{en/02_find_and_visualize:find-en}}\label{\detokenize{en/02_find_and_visualize::doc}}
\begin{sphinxcontents}
\sphinxstylecontentstitle{On this page}
\begin{itemize}
\item {} 
\sphinxAtStartPar
\phantomsection\label{\detokenize{en/02_find_and_visualize:id1}}{\hyperref[\detokenize{en/02_find_and_visualize:accessing-the-catalogue}]{\sphinxcrossref{Accessing the Catalogue}}}

\item {} 
\sphinxAtStartPar
\phantomsection\label{\detokenize{en/02_find_and_visualize:id2}}{\hyperref[\detokenize{en/02_find_and_visualize:filtering-and-searching-the-catalogue}]{\sphinxcrossref{Filtering and Searching the Catalogue}}}
\begin{itemize}
\item {} 
\sphinxAtStartPar
\phantomsection\label{\detokenize{en/02_find_and_visualize:id3}}{\hyperref[\detokenize{en/02_find_and_visualize:using-the-search-bar}]{\sphinxcrossref{Using the Search Bar}}}

\item {} 
\sphinxAtStartPar
\phantomsection\label{\detokenize{en/02_find_and_visualize:id4}}{\hyperref[\detokenize{en/02_find_and_visualize:using-filters}]{\sphinxcrossref{Using Filters}}}

\end{itemize}

\item {} 
\sphinxAtStartPar
\phantomsection\label{\detokenize{en/02_find_and_visualize:id5}}{\hyperref[\detokenize{en/02_find_and_visualize:step-by-step-example-finding-a-pelagos-dataset}]{\sphinxcrossref{Step\sphinxhyphen{}by\sphinxhyphen{}Step Example: Finding a Pelagos Dataset}}}
\begin{itemize}
\item {} 
\sphinxAtStartPar
\phantomsection\label{\detokenize{en/02_find_and_visualize:id6}}{\hyperref[\detokenize{en/02_find_and_visualize:step-1-open-the-catalogue}]{\sphinxcrossref{Step 1 \textendash{} Open the catalogue}}}

\item {} 
\sphinxAtStartPar
\phantomsection\label{\detokenize{en/02_find_and_visualize:id7}}{\hyperref[\detokenize{en/02_find_and_visualize:step-2-search-by-keyword}]{\sphinxcrossref{Step 2 \textendash{} Search by keyword}}}

\item {} 
\sphinxAtStartPar
\phantomsection\label{\detokenize{en/02_find_and_visualize:id8}}{\hyperref[\detokenize{en/02_find_and_visualize:step-3-filter-to-datasets-only}]{\sphinxcrossref{Step 3 \textendash{} Filter to Datasets only}}}

\item {} 
\sphinxAtStartPar
\phantomsection\label{\detokenize{en/02_find_and_visualize:id9}}{\hyperref[\detokenize{en/02_find_and_visualize:step-4-select-a-dataset}]{\sphinxcrossref{Step 4 \textendash{} Select a Dataset}}}

\end{itemize}

\item {} 
\sphinxAtStartPar
\phantomsection\label{\detokenize{en/02_find_and_visualize:id10}}{\hyperref[\detokenize{en/02_find_and_visualize:the-dataset-detail-page}]{\sphinxcrossref{The Dataset Detail Page}}}
\begin{itemize}
\item {} 
\sphinxAtStartPar
\phantomsection\label{\detokenize{en/02_find_and_visualize:id11}}{\hyperref[\detokenize{en/02_find_and_visualize:thumbnail-and-actions}]{\sphinxcrossref{Thumbnail and Actions}}}

\item {} 
\sphinxAtStartPar
\phantomsection\label{\detokenize{en/02_find_and_visualize:id12}}{\hyperref[\detokenize{en/02_find_and_visualize:interactive-preview}]{\sphinxcrossref{Interactive Preview}}}

\item {} 
\sphinxAtStartPar
\phantomsection\label{\detokenize{en/02_find_and_visualize:id13}}{\hyperref[\detokenize{en/02_find_and_visualize:metadata-panel}]{\sphinxcrossref{Metadata Panel}}}

\item {} 
\sphinxAtStartPar
\phantomsection\label{\detokenize{en/02_find_and_visualize:id14}}{\hyperref[\detokenize{en/02_find_and_visualize:full-metadata-view}]{\sphinxcrossref{Full Metadata View}}}

\item {} 
\sphinxAtStartPar
\phantomsection\label{\detokenize{en/02_find_and_visualize:id15}}{\hyperref[\detokenize{en/02_find_and_visualize:downloading-a-dataset}]{\sphinxcrossref{Downloading a Dataset}}}

\end{itemize}

\end{itemize}
\end{sphinxcontents}

\sphinxAtStartPar
This section explains how to discover resources in KMaP, apply filters to narrow your
search, and explore the detail page and preview of a Dataset.


\subsection{Accessing the Catalogue}
\label{\detokenize{en/02_find_and_visualize:accessing-the-catalogue}}
\sphinxAtStartPar
The KMaP catalogue lists all public resources and all resources shared with your group.
There are two main ways to reach it:
\begin{enumerate}
\sphinxsetlistlabels{\arabic}{enumi}{enumii}{}{.}%
\item {} 
\sphinxAtStartPar
\sphinxstylestrong{From the navigation bar} \textendash{} click \sphinxstylestrong{Maps} (for spatial resources) or \sphinxstylestrong{Library}
(for documents and publications) in the top menu.

\item {} 
\sphinxAtStartPar
\sphinxstylestrong{Directly via URL} \textendash{} navigate to \sphinxurl{https://kmap.info-rac.org/resources/}

\end{enumerate}

\sphinxAtStartPar
The catalogue page displays resources as cards, each showing a thumbnail, title, resource
type, and summary metadata.

\begin{sphinxadmonition}{tip}{Tip:}
\sphinxAtStartPar
Make sure you are \sphinxstylestrong{logged in} to see resources that are shared exclusively with the
Pelagos Agreement group. Resources restricted to your group are not visible to
unauthenticated users.
\end{sphinxadmonition}


\subsection{Filtering and Searching the Catalogue}
\label{\detokenize{en/02_find_and_visualize:filtering-and-searching-the-catalogue}}

\subsubsection{Using the Search Bar}
\label{\detokenize{en/02_find_and_visualize:using-the-search-bar}}
\sphinxAtStartPar
The search bar at the top of the catalogue page accepts free\sphinxhyphen{}text queries. It searches
across titles, abstracts, and keyword fields.

\sphinxAtStartPar
To find Pelagos\sphinxhyphen{}related content:
\begin{enumerate}
\sphinxsetlistlabels{\arabic}{enumi}{enumii}{}{.}%
\item {} 
\sphinxAtStartPar
Type \sphinxcode{\sphinxupquote{pelagos sanctuary}} in the search bar and press \sphinxstylestrong{Enter}.

\item {} 
\sphinxAtStartPar
Alternatively, search for \sphinxcode{\sphinxupquote{pelagos agreement}} to retrieve resources tagged
with that keyword.

\end{enumerate}

\begin{sphinxadmonition}{note}{Note:}
\sphinxAtStartPar
Both \sphinxcode{\sphinxupquote{pelagos sanctuary}} and \sphinxcode{\sphinxupquote{pelagos agreement}} are controlled keywords used by
the Pelagos Agreement group when publishing resources on KMaP. Using these exact terms
will return the most relevant results.
\end{sphinxadmonition}


\subsubsection{Using Filters}
\label{\detokenize{en/02_find_and_visualize:using-filters}}
\sphinxAtStartPar
On the left side of the catalogue page you will find a filter panel. Filters can be
combined freely and are applied immediately.


\begin{savenotes}\sphinxattablestart
\sphinxthistablewithglobalstyle
\centering
\begin{tabular}[t]{\X{25}{100}\X{75}{100}}
\sphinxtoprule
\sphinxstyletheadfamily 
\sphinxAtStartPar
Filter
&\sphinxstyletheadfamily 
\sphinxAtStartPar
How to use it
\\
\sphinxmidrule
\sphinxtableatstartofbodyhook
\sphinxAtStartPar
\sphinxstylestrong{Resource type}
&
\sphinxAtStartPar
Select \sphinxstyleemphasis{Dataset}, \sphinxstyleemphasis{Document}, \sphinxstyleemphasis{Map}, \sphinxstyleemphasis{GeoStory}, or \sphinxstyleemphasis{Dashboard} to restrict
results to one type.
\\
\sphinxhline
\sphinxAtStartPar
\sphinxstylestrong{Keywords}
&
\sphinxAtStartPar
Type or select a keyword (e.g. \sphinxcode{\sphinxupquote{pelagos sanctuary}}) to filter by thematic tag.
\\
\sphinxhline
\sphinxAtStartPar
\sphinxstylestrong{Category / Theme}
&
\sphinxAtStartPar
Select a thematic category such as \sphinxstyleemphasis{Marine Biodiversity} or \sphinxstyleemphasis{Governance}.
\\
\sphinxhline
\sphinxAtStartPar
\sphinxstylestrong{Date}
&
\sphinxAtStartPar
Filter by publication or modification date range.
\\
\sphinxhline
\sphinxAtStartPar
\sphinxstylestrong{Owner}
&
\sphinxAtStartPar
Filter by the user or organisation that published the resource.
\\
\sphinxhline
\sphinxAtStartPar
\sphinxstylestrong{Extent}
&
\sphinxAtStartPar
Draw a bounding box on the mini\sphinxhyphen{}map to filter by geographic area.
\\
\sphinxbottomrule
\end{tabular}
\sphinxtableafterendhook\par
\sphinxattableend\end{savenotes}

\begin{sphinxadmonition}{tip}{Tip:}
\sphinxAtStartPar
To find all Pelagos\sphinxhyphen{}related datasets quickly:
\begin{enumerate}
\sphinxsetlistlabels{\arabic}{enumi}{enumii}{}{.}%
\item {} 
\sphinxAtStartPar
Set \sphinxstylestrong{Resource type} \(\rightarrow\) \sphinxstyleemphasis{Dataset}

\item {} 
\sphinxAtStartPar
In the \sphinxstylestrong{Keywords} filter, type \sphinxcode{\sphinxupquote{pelagos sanctuary}}

\item {} 
\sphinxAtStartPar
The catalogue will update to show only datasets tagged with that keyword.

\end{enumerate}
\end{sphinxadmonition}


\subsection{Step\sphinxhyphen{}by\sphinxhyphen{}Step Example: Finding a Pelagos Dataset}
\label{\detokenize{en/02_find_and_visualize:step-by-step-example-finding-a-pelagos-dataset}}
\sphinxAtStartPar
The following example walks through finding, previewing, and reading the metadata of a
Dataset relevant to the Pelagos Agreement.


\subsubsection{Step 1 \textendash{} Open the catalogue}
\label{\detokenize{en/02_find_and_visualize:step-1-open-the-catalogue}}
\sphinxAtStartPar
Go to \sphinxurl{https://kmap.info-rac.org} and click \sphinxstylestrong{Maps} in the top navigation bar.
This opens the Data Hub catalogue, filtered to spatial resource types.


\subsubsection{Step 2 \textendash{} Search by keyword}
\label{\detokenize{en/02_find_and_visualize:step-2-search-by-keyword}}
\sphinxAtStartPar
In the search bar, type:

\begin{sphinxVerbatim}[commandchars=\\\{\}]
\PYG{n}{pelagos} \PYG{n}{sanctuary}
\end{sphinxVerbatim}

\sphinxAtStartPar
Press \sphinxstylestrong{Enter}. The results panel updates to show resources matching that keyword.


\subsubsection{Step 3 \textendash{} Filter to Datasets only}
\label{\detokenize{en/02_find_and_visualize:step-3-filter-to-datasets-only}}
\sphinxAtStartPar
In the left filter panel, under \sphinxstylestrong{Resource type}, tick \sphinxstylestrong{Dataset}.
The list now shows only datasets tagged with \sphinxstyleemphasis{pelagos sanctuary}.


\subsubsection{Step 4 \textendash{} Select a Dataset}
\label{\detokenize{en/02_find_and_visualize:step-4-select-a-dataset}}
\sphinxAtStartPar
Click on a dataset card to open its \sphinxstylestrong{detail page}.


\subsection{The Dataset Detail Page}
\label{\detokenize{en/02_find_and_visualize:the-dataset-detail-page}}
\sphinxAtStartPar
The detail page is the central information page for any resource. It is divided into
several sections.


\subsubsection{Thumbnail and Actions}
\label{\detokenize{en/02_find_and_visualize:thumbnail-and-actions}}
\sphinxAtStartPar
At the top of the page you will see:
\begin{itemize}
\item {} 
\sphinxAtStartPar
A \sphinxstylestrong{map thumbnail} or preview image of the dataset.

\item {} 
\sphinxAtStartPar
Action buttons:
\begin{itemize}
\item {} 
\sphinxAtStartPar
\sphinxstylestrong{View map} \textendash{} opens the dataset in the interactive map viewer.

\item {} 
\sphinxAtStartPar
\sphinxstylestrong{Download} \textendash{} downloads the dataset in your chosen format.

\item {} 
\sphinxAtStartPar
\sphinxstylestrong{Edit} (if you have permission) \textendash{} opens the metadata editor.

\item {} 
\sphinxAtStartPar
\sphinxstylestrong{Share} \textendash{} copies the permalink to the resource.

\end{itemize}

\end{itemize}


\subsubsection{Interactive Preview}
\label{\detokenize{en/02_find_and_visualize:interactive-preview}}
\sphinxAtStartPar
Below the thumbnail, the \sphinxstylestrong{map preview panel} displays the dataset rendered on a
basemap. You can:
\begin{itemize}
\item {} 
\sphinxAtStartPar
Pan and zoom using mouse or touch gestures.

\item {} 
\sphinxAtStartPar
Click on a feature to open a \sphinxstylestrong{pop\sphinxhyphen{}up} showing its attribute values.

\item {} 
\sphinxAtStartPar
Use the layer panel (☰ icon) to toggle visibility or inspect the legend.

\end{itemize}

\begin{sphinxadmonition}{note}{Note:}
\sphinxAtStartPar
The preview is a live, interactive map — not a static image. You can explore the
data spatially before deciding whether to download or use it.
\end{sphinxadmonition}


\subsubsection{Metadata Panel}
\label{\detokenize{en/02_find_and_visualize:metadata-panel}}
\sphinxAtStartPar
Below the preview, the \sphinxstylestrong{metadata panel} contains detailed information about the
dataset. The key fields are:


\begin{savenotes}\sphinxattablestart
\sphinxthistablewithglobalstyle
\centering
\begin{tabular}[t]{\X{30}{100}\X{70}{100}}
\sphinxtoprule
\sphinxstyletheadfamily 
\sphinxAtStartPar
Field
&\sphinxstyletheadfamily 
\sphinxAtStartPar
Description
\\
\sphinxmidrule
\sphinxtableatstartofbodyhook
\sphinxAtStartPar
\sphinxstylestrong{Title}
&
\sphinxAtStartPar
The full name of the dataset.
\\
\sphinxhline
\sphinxAtStartPar
\sphinxstylestrong{Abstract}
&
\sphinxAtStartPar
A plain\sphinxhyphen{}language description of what the dataset contains and its purpose.
\\
\sphinxhline
\sphinxAtStartPar
\sphinxstylestrong{Keywords}
&
\sphinxAtStartPar
Thematic and geographic tags (look for \sphinxcode{\sphinxupquote{pelagos sanctuary}} or
\sphinxcode{\sphinxupquote{pelagos agreement}} here).
\\
\sphinxhline
\sphinxAtStartPar
\sphinxstylestrong{Category}
&
\sphinxAtStartPar
The INSPIRE or thematic category (e.g. \sphinxstyleemphasis{Marine Biodiversity}).
\\
\sphinxhline
\sphinxAtStartPar
\sphinxstylestrong{Date}
&
\sphinxAtStartPar
Publication date and last modification date.
\\
\sphinxhline
\sphinxAtStartPar
\sphinxstylestrong{Spatial extent}
&
\sphinxAtStartPar
Bounding box or geographic region covered by the data.
\\
\sphinxhline
\sphinxAtStartPar
\sphinxstylestrong{Coordinate Reference System (CRS)}
&
\sphinxAtStartPar
The projection used (e.g. WGS 84 / EPSG:4326).
\\
\sphinxhline
\sphinxAtStartPar
\sphinxstylestrong{Owner}
&
\sphinxAtStartPar
The user or organisation responsible for the dataset.
\\
\sphinxhline
\sphinxAtStartPar
\sphinxstylestrong{Licence}
&
\sphinxAtStartPar
The usage licence (e.g. Creative Commons, ODbL, …).
\\
\sphinxhline
\sphinxAtStartPar
\sphinxstylestrong{Point of contact}
&
\sphinxAtStartPar
Who to contact for questions about the data.
\\
\sphinxhline
\sphinxAtStartPar
\sphinxstylestrong{OGC Services}
&
\sphinxAtStartPar
WMS, WFS, and WCS endpoint URLs for direct integration in GIS software.
\\
\sphinxhline
\sphinxAtStartPar
\sphinxstylestrong{Related resources}
&
\sphinxAtStartPar
Links to Maps, Documents, or other Datasets connected to this resource.
\\
\sphinxbottomrule
\end{tabular}
\sphinxtableafterendhook\par
\sphinxattableend\end{savenotes}


\subsubsection{Full Metadata View}
\label{\detokenize{en/02_find_and_visualize:full-metadata-view}}
\sphinxAtStartPar
To see the complete metadata record, click the \sphinxstylestrong{Metadata} tab (or the \sphinxstyleemphasis{View full
metadata} link near the bottom of the detail page). This view displays all metadata
fields in their full form, including ISO 19115 / INSPIRE\sphinxhyphen{}compliant fields.


\subsubsection{Downloading a Dataset}
\label{\detokenize{en/02_find_and_visualize:downloading-a-dataset}}
\sphinxAtStartPar
Click the \sphinxstylestrong{Download} button on the detail page. A format selector appears. Common
options include:
\begin{itemize}
\item {} 
\sphinxAtStartPar
\sphinxstylestrong{ESRI Shapefile} \textendash{} for use in QGIS, ArcGIS, and most desktop GIS.

\item {} 
\sphinxAtStartPar
\sphinxstylestrong{GeoJSON} \textendash{} for web mapping and lightweight GIS use.

\item {} 
\sphinxAtStartPar
\sphinxstylestrong{KML/KMZ} \textendash{} for Google Earth.

\item {} 
\sphinxAtStartPar
\sphinxstylestrong{GeoPackage} \textendash{} a single\sphinxhyphen{}file format supported by QGIS.

\item {} 
\sphinxAtStartPar
\sphinxstylestrong{GeoTIFF} \textendash{} for raster datasets.

\item {} 
\sphinxAtStartPar
\sphinxstylestrong{CSV} \textendash{} attribute table only (no geometry).

\end{itemize}

\sphinxAtStartPar
Select your preferred format and click \sphinxstylestrong{Download} to save the file.

\sphinxstepscope


\section{Managing Metadata}
\label{\detokenize{en/03_managing_metadata:managing-metadata}}\label{\detokenize{en/03_managing_metadata:metadata-en}}\label{\detokenize{en/03_managing_metadata::doc}}
\begin{sphinxcontents}
\sphinxstylecontentstitle{On this page}
\begin{itemize}
\item {} 
\sphinxAtStartPar
\phantomsection\label{\detokenize{en/03_managing_metadata:id1}}{\hyperref[\detokenize{en/03_managing_metadata:who-can-edit-metadata}]{\sphinxcrossref{Who Can Edit Metadata?}}}

\item {} 
\sphinxAtStartPar
\phantomsection\label{\detokenize{en/03_managing_metadata:id2}}{\hyperref[\detokenize{en/03_managing_metadata:opening-the-metadata-editor}]{\sphinxcrossref{Opening the Metadata Editor}}}

\item {} 
\sphinxAtStartPar
\phantomsection\label{\detokenize{en/03_managing_metadata:id3}}{\hyperref[\detokenize{en/03_managing_metadata:overview-of-the-metadata-editor}]{\sphinxcrossref{Overview of the Metadata Editor}}}

\item {} 
\sphinxAtStartPar
\phantomsection\label{\detokenize{en/03_managing_metadata:id4}}{\hyperref[\detokenize{en/03_managing_metadata:editing-key-metadata-fields}]{\sphinxcrossref{Editing Key Metadata Fields}}}
\begin{itemize}
\item {} 
\sphinxAtStartPar
\phantomsection\label{\detokenize{en/03_managing_metadata:id5}}{\hyperref[\detokenize{en/03_managing_metadata:title}]{\sphinxcrossref{Title}}}

\item {} 
\sphinxAtStartPar
\phantomsection\label{\detokenize{en/03_managing_metadata:id6}}{\hyperref[\detokenize{en/03_managing_metadata:abstract}]{\sphinxcrossref{Abstract}}}

\item {} 
\sphinxAtStartPar
\phantomsection\label{\detokenize{en/03_managing_metadata:id7}}{\hyperref[\detokenize{en/03_managing_metadata:keywords}]{\sphinxcrossref{Keywords}}}

\item {} 
\sphinxAtStartPar
\phantomsection\label{\detokenize{en/03_managing_metadata:id8}}{\hyperref[\detokenize{en/03_managing_metadata:category}]{\sphinxcrossref{Category}}}

\item {} 
\sphinxAtStartPar
\phantomsection\label{\detokenize{en/03_managing_metadata:id9}}{\hyperref[\detokenize{en/03_managing_metadata:spatial-extent}]{\sphinxcrossref{Spatial Extent}}}

\item {} 
\sphinxAtStartPar
\phantomsection\label{\detokenize{en/03_managing_metadata:id10}}{\hyperref[\detokenize{en/03_managing_metadata:date-fields}]{\sphinxcrossref{Date Fields}}}

\item {} 
\sphinxAtStartPar
\phantomsection\label{\detokenize{en/03_managing_metadata:id11}}{\hyperref[\detokenize{en/03_managing_metadata:licence}]{\sphinxcrossref{Licence}}}

\item {} 
\sphinxAtStartPar
\phantomsection\label{\detokenize{en/03_managing_metadata:id12}}{\hyperref[\detokenize{en/03_managing_metadata:point-of-contact}]{\sphinxcrossref{Point of Contact}}}

\item {} 
\sphinxAtStartPar
\phantomsection\label{\detokenize{en/03_managing_metadata:id13}}{\hyperref[\detokenize{en/03_managing_metadata:thumbnail}]{\sphinxcrossref{Thumbnail}}}

\end{itemize}

\item {} 
\sphinxAtStartPar
\phantomsection\label{\detokenize{en/03_managing_metadata:id14}}{\hyperref[\detokenize{en/03_managing_metadata:saving-your-changes}]{\sphinxcrossref{Saving Your Changes}}}

\item {} 
\sphinxAtStartPar
\phantomsection\label{\detokenize{en/03_managing_metadata:id15}}{\hyperref[\detokenize{en/03_managing_metadata:metadata-completeness-checklist}]{\sphinxcrossref{Metadata Completeness Checklist}}}

\end{itemize}
\end{sphinxcontents}

\sphinxAtStartPar
Good metadata is what makes a resource discoverable and trustworthy. This section
explains how to view and edit the metadata of a Dataset or Document in KMaP.

\begin{sphinxadmonition}{note}{Note:}
\sphinxAtStartPar
You can only edit the metadata of resources that you \sphinxstylestrong{own} or for which you have
been granted \sphinxstylestrong{edit permissions}. Members of the Pelagos Agreement group may have
edit rights on resources shared within the group — check with your group administrator
if you are unsure.
\end{sphinxadmonition}


\subsection{Who Can Edit Metadata?}
\label{\detokenize{en/03_managing_metadata:who-can-edit-metadata}}
\sphinxAtStartPar
KMaP has a permission model with several levels:


\begin{savenotes}\sphinxattablestart
\sphinxthistablewithglobalstyle
\centering
\begin{tabular}[t]{\X{25}{100}\X{75}{100}}
\sphinxtoprule
\sphinxstyletheadfamily 
\sphinxAtStartPar
Role
&\sphinxstyletheadfamily 
\sphinxAtStartPar
Metadata editing rights
\\
\sphinxmidrule
\sphinxtableatstartofbodyhook
\sphinxAtStartPar
\sphinxstylestrong{Resource owner}
&
\sphinxAtStartPar
Full edit rights on all metadata fields.
\\
\sphinxhline
\sphinxAtStartPar
\sphinxstylestrong{Group manager}
&
\sphinxAtStartPar
Can edit resources owned by any member of their group.
\\
\sphinxhline
\sphinxAtStartPar
\sphinxstylestrong{Group member (editor)}
&
\sphinxAtStartPar
Can edit resources explicitly shared with the group with \sphinxstyleemphasis{edit} permission.
\\
\sphinxhline
\sphinxAtStartPar
\sphinxstylestrong{Viewer / public}
&
\sphinxAtStartPar
Read\sphinxhyphen{}only access. Cannot edit metadata.
\\
\sphinxbottomrule
\end{tabular}
\sphinxtableafterendhook\par
\sphinxattableend\end{savenotes}

\sphinxAtStartPar
If you see an \sphinxstylestrong{Edit} button on the resource detail page, you have the right to edit
that resource’s metadata.


\subsection{Opening the Metadata Editor}
\label{\detokenize{en/03_managing_metadata:opening-the-metadata-editor}}
\sphinxAtStartPar
There are two ways to open the metadata editor:

\sphinxAtStartPar
\sphinxstylestrong{From the detail page:}
\begin{enumerate}
\sphinxsetlistlabels{\arabic}{enumi}{enumii}{}{.}%
\item {} 
\sphinxAtStartPar
Navigate to the resource detail page (see {\hyperref[\detokenize{en/02_find_and_visualize:find-en}]{\sphinxcrossref{\DUrole{std}{\DUrole{std-ref}{How to Find and Visualise Your Data}}}}} for how to find a resource).

\item {} 
\sphinxAtStartPar
Click the \sphinxstylestrong{Edit} button (pencil icon) in the top\sphinxhyphen{}right action bar.

\item {} 
\sphinxAtStartPar
A dropdown menu appears. Select \sphinxstylestrong{Edit metadata}.

\end{enumerate}

\sphinxAtStartPar
\sphinxstylestrong{From the catalogue:}
\begin{enumerate}
\sphinxsetlistlabels{\arabic}{enumi}{enumii}{}{.}%
\item {} 
\sphinxAtStartPar
In the catalogue, hover over a resource card.

\item {} 
\sphinxAtStartPar
Click the \sphinxstylestrong{⋮} (more options) menu that appears.

\item {} 
\sphinxAtStartPar
Select \sphinxstylestrong{Edit metadata}.

\end{enumerate}

\sphinxAtStartPar
The metadata editor opens as a multi\sphinxhyphen{}tab form.


\subsection{Overview of the Metadata Editor}
\label{\detokenize{en/03_managing_metadata:overview-of-the-metadata-editor}}
\sphinxAtStartPar
The metadata editor is organised into tabs. The most important tabs are:


\begin{savenotes}\sphinxattablestart
\sphinxthistablewithglobalstyle
\centering
\begin{tabular}[t]{\X{25}{100}\X{75}{100}}
\sphinxtoprule
\sphinxstyletheadfamily 
\sphinxAtStartPar
Tab
&\sphinxstyletheadfamily 
\sphinxAtStartPar
Content
\\
\sphinxmidrule
\sphinxtableatstartofbodyhook
\sphinxAtStartPar
\sphinxstylestrong{Basic info}
&
\sphinxAtStartPar
Title, abstract, language, category, keywords, point of contact.
\\
\sphinxhline
\sphinxAtStartPar
\sphinxstylestrong{Location and licences}
&
\sphinxAtStartPar
Spatial extent, CRS, date fields, licence, attribution.
\\
\sphinxhline
\sphinxAtStartPar
\sphinxstylestrong{Optional metadata}
&
\sphinxAtStartPar
Additional ISO 19115 fields (edition, purpose, supplemental information, …).
\\
\sphinxhline
\sphinxAtStartPar
\sphinxstylestrong{Settings}
&
\sphinxAtStartPar
Permissions, featured status, advertised status.
\\
\sphinxbottomrule
\end{tabular}
\sphinxtableafterendhook\par
\sphinxattableend\end{savenotes}


\subsection{Editing Key Metadata Fields}
\label{\detokenize{en/03_managing_metadata:editing-key-metadata-fields}}

\subsubsection{Title}
\label{\detokenize{en/03_managing_metadata:title}}
\sphinxAtStartPar
The \sphinxstylestrong{title} is the primary way users find your resource. It should be:
\begin{itemize}
\item {} 
\sphinxAtStartPar
Descriptive and specific (avoid generic titles like “Data” or “Map layer”).

\item {} 
\sphinxAtStartPar
Consistent with the naming conventions used by the Pelagos Agreement group.

\end{itemize}

\sphinxAtStartPar
To edit: click on the \sphinxstylestrong{Title} field in the \sphinxstyleemphasis{Basic info} tab and type your changes.


\subsubsection{Abstract}
\label{\detokenize{en/03_managing_metadata:abstract}}
\sphinxAtStartPar
The \sphinxstylestrong{abstract} is a plain\sphinxhyphen{}language description of the resource. A good abstract answers:
\begin{itemize}
\item {} 
\sphinxAtStartPar
What does this resource contain?

\item {} 
\sphinxAtStartPar
What geographic area does it cover?

\item {} 
\sphinxAtStartPar
What time period does it represent?

\item {} 
\sphinxAtStartPar
What is it used for?

\item {} 
\sphinxAtStartPar
What are its known limitations?

\end{itemize}

\sphinxAtStartPar
The abstract field supports basic formatting. Aim for at least two or three sentences.


\subsubsection{Keywords}
\label{\detokenize{en/03_managing_metadata:keywords}}
\sphinxAtStartPar
Keywords are the primary mechanism for filtering and discovering resources in KMaP.

\begin{sphinxadmonition}{important}{Important:}
\sphinxAtStartPar
All resources related to the Pelagos Agreement \sphinxstylestrong{must include} at least one of the
following keywords:
\begin{itemize}
\item {} 
\sphinxAtStartPar
\sphinxcode{\sphinxupquote{pelagos sanctuary}}

\item {} 
\sphinxAtStartPar
\sphinxcode{\sphinxupquote{pelagos agreement}}

\end{itemize}

\sphinxAtStartPar
Without these keywords, your resource will not appear when other group members
filter the catalogue by these terms.
\end{sphinxadmonition}

\sphinxAtStartPar
To add a keyword:
\begin{enumerate}
\sphinxsetlistlabels{\arabic}{enumi}{enumii}{}{.}%
\item {} 
\sphinxAtStartPar
In the \sphinxstylestrong{Keywords} field, start typing the keyword.

\item {} 
\sphinxAtStartPar
If the keyword already exists in the system, it will appear in a dropdown — select it.

\item {} 
\sphinxAtStartPar
If it is new, type the full keyword and press \sphinxstylestrong{Enter} or click \sphinxstylestrong{Add}.

\item {} 
\sphinxAtStartPar
To remove a keyword, click the \sphinxstylestrong{×} next to it.

\end{enumerate}


\subsubsection{Category}
\label{\detokenize{en/03_managing_metadata:category}}
\sphinxAtStartPar
The \sphinxstylestrong{category} assigns the resource to a thematic classification.
Select the most appropriate category from the dropdown:
\begin{itemize}
\item {} 
\sphinxAtStartPar
Fishery and Aquaculture

\item {} 
\sphinxAtStartPar
Marine Biodiversity

\item {} 
\sphinxAtStartPar
Pollution

\item {} 
\sphinxAtStartPar
Climate Change

\item {} 
\sphinxAtStartPar
Marine Spatial Planning

\item {} 
\sphinxAtStartPar
Sustainability and Blue Economy

\item {} 
\sphinxAtStartPar
Governance

\end{itemize}

\sphinxAtStartPar
A resource can belong to only one primary category. Use keywords for additional
thematic tagging.


\subsubsection{Spatial Extent}
\label{\detokenize{en/03_managing_metadata:spatial-extent}}
\sphinxAtStartPar
The \sphinxstylestrong{spatial extent} defines the geographic bounding box of the resource. For datasets,
this is usually set automatically when the data is uploaded. You can refine it manually:
\begin{enumerate}
\sphinxsetlistlabels{\arabic}{enumi}{enumii}{}{.}%
\item {} 
\sphinxAtStartPar
Go to the \sphinxstylestrong{Location and licences} tab.

\item {} 
\sphinxAtStartPar
In the \sphinxstyleemphasis{Spatial extent} section, use the map widget to draw or adjust the bounding box,
or enter coordinates (min/max longitude and latitude) manually.

\end{enumerate}

\begin{sphinxadmonition}{tip}{Tip:}
\sphinxAtStartPar
For Pelagos\sphinxhyphen{}related resources, the spatial extent should cover at least the Pelagos
Sanctuary area (approximately between France, Italy, and Monaco in the Ligurian Sea and
adjacent western Mediterranean).
\end{sphinxadmonition}


\subsubsection{Date Fields}
\label{\detokenize{en/03_managing_metadata:date-fields}}
\sphinxAtStartPar
Three date fields are available:


\begin{savenotes}\sphinxattablestart
\sphinxthistablewithglobalstyle
\centering
\begin{tabular}[t]{\X{30}{100}\X{70}{100}}
\sphinxtoprule
\sphinxstyletheadfamily 
\sphinxAtStartPar
Field
&\sphinxstyletheadfamily 
\sphinxAtStartPar
Description
\\
\sphinxmidrule
\sphinxtableatstartofbodyhook
\sphinxAtStartPar
\sphinxstylestrong{Date} (creation date)
&
\sphinxAtStartPar
When the resource was originally created or the data was collected.
\\
\sphinxhline
\sphinxAtStartPar
\sphinxstylestrong{Publication date}
&
\sphinxAtStartPar
When the resource was published on KMaP.
\\
\sphinxhline
\sphinxAtStartPar
\sphinxstylestrong{Revision date}
&
\sphinxAtStartPar
When the resource was last updated.
\\
\sphinxbottomrule
\end{tabular}
\sphinxtableafterendhook\par
\sphinxattableend\end{savenotes}

\sphinxAtStartPar
Keep these dates accurate — they are used to filter and sort resources in the catalogue.


\subsubsection{Licence}
\label{\detokenize{en/03_managing_metadata:licence}}
\sphinxAtStartPar
The \sphinxstylestrong{licence} defines how others may use the resource.
\begin{enumerate}
\sphinxsetlistlabels{\arabic}{enumi}{enumii}{}{.}%
\item {} 
\sphinxAtStartPar
Go to the \sphinxstylestrong{Location and licences} tab.

\item {} 
\sphinxAtStartPar
Select a licence from the dropdown. Common options include:
\begin{itemize}
\item {} 
\sphinxAtStartPar
\sphinxstyleemphasis{Creative Commons Attribution (CC BY)} \textendash{} free reuse with attribution.

\item {} 
\sphinxAtStartPar
\sphinxstyleemphasis{Creative Commons Attribution ShareAlike (CC BY\sphinxhyphen{}SA)} \textendash{} reuse with attribution and
same licence.

\item {} 
\sphinxAtStartPar
\sphinxstyleemphasis{Open Database Licence (ODbL)} \textendash{} for datasets.

\item {} 
\sphinxAtStartPar
\sphinxstyleemphasis{No restrictions} \textendash{} public domain.

\end{itemize}

\end{enumerate}

\begin{sphinxadmonition}{note}{Note:}
\sphinxAtStartPar
If you are unsure which licence to apply, consult your group administrator. Applying
the wrong licence may restrict the reuse of your data inappropriately.
\end{sphinxadmonition}


\subsubsection{Point of Contact}
\label{\detokenize{en/03_managing_metadata:point-of-contact}}
\sphinxAtStartPar
The \sphinxstylestrong{point of contact} identifies who to contact for questions about the resource.
It can be a person or an organisation.
\begin{enumerate}
\sphinxsetlistlabels{\arabic}{enumi}{enumii}{}{.}%
\item {} 
\sphinxAtStartPar
In the \sphinxstyleemphasis{Basic info} tab, find the \sphinxstylestrong{Point of contact} field.

\item {} 
\sphinxAtStartPar
Start typing a name to search registered KMaP users, or enter an organisation name.

\end{enumerate}


\subsubsection{Thumbnail}
\label{\detokenize{en/03_managing_metadata:thumbnail}}
\sphinxAtStartPar
The \sphinxstylestrong{thumbnail} is the preview image displayed in catalogue cards. It is generated
automatically for Datasets and Maps, but can be replaced manually:
\begin{enumerate}
\sphinxsetlistlabels{\arabic}{enumi}{enumii}{}{.}%
\item {} 
\sphinxAtStartPar
Click \sphinxstylestrong{Edit} on the detail page.

\item {} 
\sphinxAtStartPar
Select \sphinxstylestrong{Edit thumbnail}.

\item {} 
\sphinxAtStartPar
Upload a new image (recommended: 600 × 400 px, JPEG or PNG).

\end{enumerate}


\subsection{Saving Your Changes}
\label{\detokenize{en/03_managing_metadata:saving-your-changes}}
\sphinxAtStartPar
When you have finished editing, scroll to the bottom of the editor and click
\sphinxstylestrong{Save} (or \sphinxstylestrong{Update} — the label may vary). A confirmation message will appear.

\begin{sphinxadmonition}{warning}{Warning:}
\sphinxAtStartPar
Metadata changes are applied immediately and are visible to all users with access to
the resource. There is no draft or staging mode.
\end{sphinxadmonition}

\sphinxAtStartPar
After saving, you are returned to the resource detail page where you can verify that
all fields have been updated correctly.


\subsection{Metadata Completeness Checklist}
\label{\detokenize{en/03_managing_metadata:metadata-completeness-checklist}}
\sphinxAtStartPar
Use the following checklist before publishing or sharing a resource:


\begin{savenotes}\sphinxattablestart
\sphinxthistablewithglobalstyle
\centering
\begin{tabular}[t]{\X{10}{100}\X{90}{100}}
\sphinxtoprule
\sphinxstyletheadfamily 
\sphinxAtStartPar
\(\checkmark\)
&\sphinxstyletheadfamily 
\sphinxAtStartPar
Item
\\
\sphinxmidrule
\sphinxtableatstartofbodyhook
\sphinxAtStartPar
☐
&
\sphinxAtStartPar
Title is descriptive and specific.
\\
\sphinxhline
\sphinxAtStartPar
☐
&
\sphinxAtStartPar
Abstract explains what, where, when, and why.
\\
\sphinxhline
\sphinxAtStartPar
☐
&
\sphinxAtStartPar
At least one Pelagos keyword is present (\sphinxcode{\sphinxupquote{pelagos sanctuary}} or
\sphinxcode{\sphinxupquote{pelagos agreement}}).
\\
\sphinxhline
\sphinxAtStartPar
☐
&
\sphinxAtStartPar
Category is set.
\\
\sphinxhline
\sphinxAtStartPar
☐
&
\sphinxAtStartPar
Spatial extent is defined and covers the correct area.
\\
\sphinxhline
\sphinxAtStartPar
☐
&
\sphinxAtStartPar
Date fields are filled in.
\\
\sphinxhline
\sphinxAtStartPar
☐
&
\sphinxAtStartPar
Licence is set.
\\
\sphinxhline
\sphinxAtStartPar
☐
&
\sphinxAtStartPar
Point of contact is identified.
\\
\sphinxhline
\sphinxAtStartPar
☐
&
\sphinxAtStartPar
Thumbnail is representative.
\\
\sphinxbottomrule
\end{tabular}
\sphinxtableafterendhook\par
\sphinxattableend\end{savenotes}

\sphinxstepscope


\section{Exploring Maps}
\label{\detokenize{en/04_exploring_maps:exploring-maps}}\label{\detokenize{en/04_exploring_maps:maps-en}}\label{\detokenize{en/04_exploring_maps::doc}}
\begin{sphinxcontents}
\sphinxstylecontentstitle{On this page}
\begin{itemize}
\item {} 
\sphinxAtStartPar
\phantomsection\label{\detokenize{en/04_exploring_maps:id1}}{\hyperref[\detokenize{en/04_exploring_maps:finding-a-map}]{\sphinxcrossref{Finding a Map}}}

\item {} 
\sphinxAtStartPar
\phantomsection\label{\detokenize{en/04_exploring_maps:id2}}{\hyperref[\detokenize{en/04_exploring_maps:the-map-viewer-interface}]{\sphinxcrossref{The Map Viewer Interface}}}

\item {} 
\sphinxAtStartPar
\phantomsection\label{\detokenize{en/04_exploring_maps:id3}}{\hyperref[\detokenize{en/04_exploring_maps:navigating-the-map}]{\sphinxcrossref{Navigating the Map}}}
\begin{itemize}
\item {} 
\sphinxAtStartPar
\phantomsection\label{\detokenize{en/04_exploring_maps:id4}}{\hyperref[\detokenize{en/04_exploring_maps:basic-navigation}]{\sphinxcrossref{Basic Navigation}}}

\item {} 
\sphinxAtStartPar
\phantomsection\label{\detokenize{en/04_exploring_maps:id5}}{\hyperref[\detokenize{en/04_exploring_maps:going-to-a-specific-location}]{\sphinxcrossref{Going to a Specific Location}}}

\end{itemize}

\item {} 
\sphinxAtStartPar
\phantomsection\label{\detokenize{en/04_exploring_maps:id6}}{\hyperref[\detokenize{en/04_exploring_maps:working-with-layers}]{\sphinxcrossref{Working with Layers}}}
\begin{itemize}
\item {} 
\sphinxAtStartPar
\phantomsection\label{\detokenize{en/04_exploring_maps:id7}}{\hyperref[\detokenize{en/04_exploring_maps:opening-the-layer-panel}]{\sphinxcrossref{Opening the Layer Panel}}}

\item {} 
\sphinxAtStartPar
\phantomsection\label{\detokenize{en/04_exploring_maps:id8}}{\hyperref[\detokenize{en/04_exploring_maps:toggling-layer-visibility}]{\sphinxcrossref{Toggling Layer Visibility}}}

\item {} 
\sphinxAtStartPar
\phantomsection\label{\detokenize{en/04_exploring_maps:id9}}{\hyperref[\detokenize{en/04_exploring_maps:changing-layer-order}]{\sphinxcrossref{Changing Layer Order}}}

\item {} 
\sphinxAtStartPar
\phantomsection\label{\detokenize{en/04_exploring_maps:id10}}{\hyperref[\detokenize{en/04_exploring_maps:adjusting-layer-opacity}]{\sphinxcrossref{Adjusting Layer Opacity}}}

\item {} 
\sphinxAtStartPar
\phantomsection\label{\detokenize{en/04_exploring_maps:id11}}{\hyperref[\detokenize{en/04_exploring_maps:viewing-the-legend}]{\sphinxcrossref{Viewing the Legend}}}

\end{itemize}

\item {} 
\sphinxAtStartPar
\phantomsection\label{\detokenize{en/04_exploring_maps:id12}}{\hyperref[\detokenize{en/04_exploring_maps:querying-features}]{\sphinxcrossref{Querying Features}}}
\begin{itemize}
\item {} 
\sphinxAtStartPar
\phantomsection\label{\detokenize{en/04_exploring_maps:id13}}{\hyperref[\detokenize{en/04_exploring_maps:clicking-on-a-feature}]{\sphinxcrossref{Clicking on a Feature}}}

\item {} 
\sphinxAtStartPar
\phantomsection\label{\detokenize{en/04_exploring_maps:id14}}{\hyperref[\detokenize{en/04_exploring_maps:identify-tool}]{\sphinxcrossref{Identify Tool}}}

\item {} 
\sphinxAtStartPar
\phantomsection\label{\detokenize{en/04_exploring_maps:id15}}{\hyperref[\detokenize{en/04_exploring_maps:selecting-and-filtering-features}]{\sphinxcrossref{Selecting and Filtering Features}}}

\end{itemize}

\item {} 
\sphinxAtStartPar
\phantomsection\label{\detokenize{en/04_exploring_maps:id16}}{\hyperref[\detokenize{en/04_exploring_maps:measurement-tools}]{\sphinxcrossref{Measurement Tools}}}
\begin{itemize}
\item {} 
\sphinxAtStartPar
\phantomsection\label{\detokenize{en/04_exploring_maps:id17}}{\hyperref[\detokenize{en/04_exploring_maps:distance-measurement}]{\sphinxcrossref{Distance Measurement}}}

\item {} 
\sphinxAtStartPar
\phantomsection\label{\detokenize{en/04_exploring_maps:id18}}{\hyperref[\detokenize{en/04_exploring_maps:area-measurement}]{\sphinxcrossref{Area Measurement}}}

\end{itemize}

\item {} 
\sphinxAtStartPar
\phantomsection\label{\detokenize{en/04_exploring_maps:id19}}{\hyperref[\detokenize{en/04_exploring_maps:printing-and-exporting-the-map-view}]{\sphinxcrossref{Printing and Exporting the Map View}}}
\begin{itemize}
\item {} 
\sphinxAtStartPar
\phantomsection\label{\detokenize{en/04_exploring_maps:id20}}{\hyperref[\detokenize{en/04_exploring_maps:print}]{\sphinxcrossref{Print}}}

\item {} 
\sphinxAtStartPar
\phantomsection\label{\detokenize{en/04_exploring_maps:id21}}{\hyperref[\detokenize{en/04_exploring_maps:share-permalink}]{\sphinxcrossref{Share / Permalink}}}

\item {} 
\sphinxAtStartPar
\phantomsection\label{\detokenize{en/04_exploring_maps:id22}}{\hyperref[\detokenize{en/04_exploring_maps:embed}]{\sphinxcrossref{Embed}}}

\end{itemize}

\item {} 
\sphinxAtStartPar
\phantomsection\label{\detokenize{en/04_exploring_maps:id23}}{\hyperref[\detokenize{en/04_exploring_maps:saving-changes-to-a-map}]{\sphinxcrossref{Saving Changes to a Map}}}

\item {} 
\sphinxAtStartPar
\phantomsection\label{\detokenize{en/04_exploring_maps:id24}}{\hyperref[\detokenize{en/04_exploring_maps:creating-a-new-map}]{\sphinxcrossref{Creating a New Map}}}

\end{itemize}
\end{sphinxcontents}

\sphinxAtStartPar
A \sphinxstylestrong{Map} in KMaP is an interactive, multi\sphinxhyphen{}layer web map built from one or more Datasets.
This section explains how to find, open, and navigate Maps, how to work with layers, and
how to extract information from the data displayed.


\subsection{Finding a Map}
\label{\detokenize{en/04_exploring_maps:finding-a-map}}
\sphinxAtStartPar
Maps are listed in the catalogue alongside all other resource types. To filter the
catalogue to show only Maps:
\begin{enumerate}
\sphinxsetlistlabels{\arabic}{enumi}{enumii}{}{.}%
\item {} 
\sphinxAtStartPar
Open the catalogue: click \sphinxstylestrong{Maps} in the top navigation bar.

\item {} 
\sphinxAtStartPar
In the filter panel on the left, under \sphinxstylestrong{Resource type}, select \sphinxstylestrong{Map}.

\item {} 
\sphinxAtStartPar
To find Pelagos\sphinxhyphen{}related maps, also apply the keyword filter:
in the \sphinxstylestrong{Keywords} field, type \sphinxcode{\sphinxupquote{pelagos sanctuary}} or \sphinxcode{\sphinxupquote{pelagos agreement}}.

\end{enumerate}

\sphinxAtStartPar
Alternatively, search for a map by title using the search bar.

\sphinxAtStartPar
Click on a Map card to open its \sphinxstylestrong{detail page}. From there, click \sphinxstylestrong{View map} to open
it in the interactive map viewer.


\subsection{The Map Viewer Interface}
\label{\detokenize{en/04_exploring_maps:the-map-viewer-interface}}
\sphinxAtStartPar
The map viewer (powered by \sphinxstylestrong{MapStore}) opens in a full\sphinxhyphen{}screen layout.
Its main interface elements are:


\begin{savenotes}\sphinxattablestart
\sphinxthistablewithglobalstyle
\centering
\begin{tabular}[t]{\X{30}{100}\X{70}{100}}
\sphinxtoprule
\sphinxstyletheadfamily 
\sphinxAtStartPar
Element
&\sphinxstyletheadfamily 
\sphinxAtStartPar
Description
\\
\sphinxmidrule
\sphinxtableatstartofbodyhook
\sphinxAtStartPar
\sphinxstylestrong{Main map canvas}
&
\sphinxAtStartPar
The central interactive map area.
\\
\sphinxhline
\sphinxAtStartPar
\sphinxstylestrong{Layer panel} (☰)
&
\sphinxAtStartPar
Toggles layer visibility, reorders layers, accesses layer settings.
\\
\sphinxhline
\sphinxAtStartPar
\sphinxstylestrong{Basemap switcher}
&
\sphinxAtStartPar
Changes the background map (OpenStreetMap, satellite imagery, …).
\\
\sphinxhline
\sphinxAtStartPar
\sphinxstylestrong{Zoom controls} (+ / −)
&
\sphinxAtStartPar
Zoom in and out. Also supports mouse wheel and pinch gesture.
\\
\sphinxhline
\sphinxAtStartPar
\sphinxstylestrong{Toolbar}
&
\sphinxAtStartPar
Row of tool icons at the top or side of the map canvas.
\\
\sphinxhline
\sphinxAtStartPar
\sphinxstylestrong{Scale bar}
&
\sphinxAtStartPar
Displays the current map scale.
\\
\sphinxhline
\sphinxAtStartPar
\sphinxstylestrong{Coordinates panel}
&
\sphinxAtStartPar
Shows the geographic coordinates of the mouse cursor position.
\\
\sphinxbottomrule
\end{tabular}
\sphinxtableafterendhook\par
\sphinxattableend\end{savenotes}


\subsection{Navigating the Map}
\label{\detokenize{en/04_exploring_maps:navigating-the-map}}

\subsubsection{Basic Navigation}
\label{\detokenize{en/04_exploring_maps:basic-navigation}}

\begin{savenotes}\sphinxattablestart
\sphinxthistablewithglobalstyle
\centering
\begin{tabular}[t]{\X{35}{100}\X{65}{100}}
\sphinxtoprule
\sphinxstyletheadfamily 
\sphinxAtStartPar
Action
&\sphinxstyletheadfamily 
\sphinxAtStartPar
How to perform it
\\
\sphinxmidrule
\sphinxtableatstartofbodyhook
\sphinxAtStartPar
\sphinxstylestrong{Pan (move the map)}
&
\sphinxAtStartPar
Click and drag on the map canvas.
\\
\sphinxhline
\sphinxAtStartPar
\sphinxstylestrong{Zoom in}
&
\sphinxAtStartPar
Scroll the mouse wheel forward, click \sphinxstylestrong{+}, or double\sphinxhyphen{}click on the map.
\\
\sphinxhline
\sphinxAtStartPar
\sphinxstylestrong{Zoom out}
&
\sphinxAtStartPar
Scroll the mouse wheel backward or click \sphinxstylestrong{−}.
\\
\sphinxhline
\sphinxAtStartPar
\sphinxstylestrong{Zoom to full extent}
&
\sphinxAtStartPar
Click the \sphinxstyleemphasis{Home} icon (🏠) in the toolbar to return to the default view.
\\
\sphinxhline
\sphinxAtStartPar
\sphinxstylestrong{Zoom to a layer}
&
\sphinxAtStartPar
In the Layer panel, right\sphinxhyphen{}click a layer and select \sphinxstyleemphasis{Zoom to layer extent}.
\\
\sphinxbottomrule
\end{tabular}
\sphinxtableafterendhook\par
\sphinxattableend\end{savenotes}


\subsubsection{Going to a Specific Location}
\label{\detokenize{en/04_exploring_maps:going-to-a-specific-location}}
\sphinxAtStartPar
Use the \sphinxstylestrong{Search / Geocoder} tool in the toolbar (🔍 icon) to search for a place name.
Type the name and press Enter — the map will zoom to that location.


\subsection{Working with Layers}
\label{\detokenize{en/04_exploring_maps:working-with-layers}}

\subsubsection{Opening the Layer Panel}
\label{\detokenize{en/04_exploring_maps:opening-the-layer-panel}}
\sphinxAtStartPar
Click the \sphinxstylestrong{☰} (hamburger) icon to open the layer panel. It lists all layers included
in the map, grouped into:
\begin{itemize}
\item {} 
\sphinxAtStartPar
\sphinxstylestrong{Overlays} \textendash{} the thematic data layers (e.g. Pelagos Sanctuary boundary, cetacean
sighting points, shipping lanes).

\item {} 
\sphinxAtStartPar
\sphinxstylestrong{Background layers} \textendash{} the basemap options (only one can be active at a time).

\end{itemize}


\subsubsection{Toggling Layer Visibility}
\label{\detokenize{en/04_exploring_maps:toggling-layer-visibility}}
\sphinxAtStartPar
Each layer in the panel has an \sphinxstylestrong{eye icon} (👁). Click it to show or hide that layer.
Hidden layers remain in the map composition but are not displayed on the canvas.


\subsubsection{Changing Layer Order}
\label{\detokenize{en/04_exploring_maps:changing-layer-order}}
\sphinxAtStartPar
Layers are drawn in the order they appear in the panel (top layers are drawn on top).
To reorder:
\begin{enumerate}
\sphinxsetlistlabels{\arabic}{enumi}{enumii}{}{.}%
\item {} 
\sphinxAtStartPar
Click and hold a layer name in the panel.

\item {} 
\sphinxAtStartPar
Drag it up or down to the desired position.

\item {} 
\sphinxAtStartPar
Release to drop.

\end{enumerate}


\subsubsection{Adjusting Layer Opacity}
\label{\detokenize{en/04_exploring_maps:adjusting-layer-opacity}}\begin{enumerate}
\sphinxsetlistlabels{\arabic}{enumi}{enumii}{}{.}%
\item {} 
\sphinxAtStartPar
In the Layer panel, click the \sphinxstylestrong{⋮} (options) icon next to a layer.

\item {} 
\sphinxAtStartPar
Select \sphinxstylestrong{Settings} (or the slider icon).

\item {} 
\sphinxAtStartPar
Use the \sphinxstylestrong{Opacity} slider to set transparency between 0 \% (invisible) and 100 \% (solid).

\end{enumerate}


\subsubsection{Viewing the Legend}
\label{\detokenize{en/04_exploring_maps:viewing-the-legend}}\begin{enumerate}
\sphinxsetlistlabels{\arabic}{enumi}{enumii}{}{.}%
\item {} 
\sphinxAtStartPar
Click the \sphinxstylestrong{⋮} (options) icon next to a layer.

\item {} 
\sphinxAtStartPar
Select \sphinxstylestrong{Legend} to expand the symbol key for that layer.

\end{enumerate}

\sphinxAtStartPar
Alternatively, some map configurations display a persistent legend panel on the side
of the canvas.


\subsection{Querying Features}
\label{\detokenize{en/04_exploring_maps:querying-features}}

\subsubsection{Clicking on a Feature}
\label{\detokenize{en/04_exploring_maps:clicking-on-a-feature}}
\sphinxAtStartPar
Click on any visible feature (point, line, or polygon) on the map canvas to open an
\sphinxstylestrong{attribute pop\sphinxhyphen{}up}. The pop\sphinxhyphen{}up displays the attribute values associated with that
feature — for example, the name of a marine protected area, a species name, or a
monitoring value.

\sphinxAtStartPar
If multiple layers overlap at the clicked location, the pop\sphinxhyphen{}up may show results from
all overlapping layers. Use the arrow buttons in the pop\sphinxhyphen{}up to navigate between them.


\subsubsection{Identify Tool}
\label{\detokenize{en/04_exploring_maps:identify-tool}}
\sphinxAtStartPar
The \sphinxstylestrong{Identify} tool (ℹ️ icon in the toolbar) works similarly to clicking directly
on the map. It queries all visible layers at the clicked point and returns their
attributes in a panel.


\subsubsection{Selecting and Filtering Features}
\label{\detokenize{en/04_exploring_maps:selecting-and-filtering-features}}
\sphinxAtStartPar
For vector layers, MapStore offers a \sphinxstylestrong{Filter} tool:
\begin{enumerate}
\sphinxsetlistlabels{\arabic}{enumi}{enumii}{}{.}%
\item {} 
\sphinxAtStartPar
In the Layer panel, click \sphinxstylestrong{⋮} next to the desired layer.

\item {} 
\sphinxAtStartPar
Select \sphinxstylestrong{Filter}.

\item {} 
\sphinxAtStartPar
Define a filter condition using the attribute fields
(e.g. show only records where \sphinxstyleemphasis{country = “France”}).

\item {} 
\sphinxAtStartPar
Click \sphinxstylestrong{Apply}. Only features matching the filter are displayed.

\end{enumerate}

\begin{sphinxadmonition}{tip}{Tip:}
\sphinxAtStartPar
Use the filter tool to focus on a specific subset of Pelagos data — for example,
showing only cetacean sightings from a specific year, or only routes that pass
through the sanctuary.
\end{sphinxadmonition}


\subsection{Measurement Tools}
\label{\detokenize{en/04_exploring_maps:measurement-tools}}

\subsubsection{Distance Measurement}
\label{\detokenize{en/04_exploring_maps:distance-measurement}}\begin{enumerate}
\sphinxsetlistlabels{\arabic}{enumi}{enumii}{}{.}%
\item {} 
\sphinxAtStartPar
Click the \sphinxstylestrong{Measure} tool (📐 icon) in the toolbar.

\item {} 
\sphinxAtStartPar
Select \sphinxstylestrong{Distance}.

\item {} 
\sphinxAtStartPar
Click points on the map to define a line. Double\sphinxhyphen{}click to end.

\item {} 
\sphinxAtStartPar
The total distance is displayed in the unit of your choice (km, miles, nm).

\end{enumerate}


\subsubsection{Area Measurement}
\label{\detokenize{en/04_exploring_maps:area-measurement}}\begin{enumerate}
\sphinxsetlistlabels{\arabic}{enumi}{enumii}{}{.}%
\item {} 
\sphinxAtStartPar
Click the \sphinxstylestrong{Measure} tool.

\item {} 
\sphinxAtStartPar
Select \sphinxstylestrong{Area}.

\item {} 
\sphinxAtStartPar
Click to define the vertices of a polygon. Double\sphinxhyphen{}click to close.

\item {} 
\sphinxAtStartPar
The calculated area is displayed.

\end{enumerate}

\begin{sphinxadmonition}{note}{Note:}
\sphinxAtStartPar
Measurements are performed on the map projection and may have small geometric
inaccuracies for large areas. For precise measurements, use a desktop GIS application.
\end{sphinxadmonition}


\subsection{Printing and Exporting the Map View}
\label{\detokenize{en/04_exploring_maps:printing-and-exporting-the-map-view}}

\subsubsection{Print}
\label{\detokenize{en/04_exploring_maps:print}}\begin{enumerate}
\sphinxsetlistlabels{\arabic}{enumi}{enumii}{}{.}%
\item {} 
\sphinxAtStartPar
Click the \sphinxstylestrong{Print} tool (🖨 icon) in the toolbar.

\item {} 
\sphinxAtStartPar
A print dialog appears. Configure:
\begin{itemize}
\item {} 
\sphinxAtStartPar
\sphinxstylestrong{Title} \textendash{} a title to display on the printed map.

\item {} 
\sphinxAtStartPar
\sphinxstylestrong{Paper size} \textendash{} A4, A3, Letter, etc.

\item {} 
\sphinxAtStartPar
\sphinxstylestrong{Resolution} \textendash{} 72, 150, or 300 dpi.

\item {} 
\sphinxAtStartPar
\sphinxstylestrong{Layout} \textendash{} portrait or landscape.

\end{itemize}

\item {} 
\sphinxAtStartPar
Click \sphinxstylestrong{Print} to generate a PDF or image.

\end{enumerate}


\subsubsection{Share / Permalink}
\label{\detokenize{en/04_exploring_maps:share-permalink}}
\sphinxAtStartPar
To share the current map view (including zoom level, active layers, and centre point):
\begin{enumerate}
\sphinxsetlistlabels{\arabic}{enumi}{enumii}{}{.}%
\item {} 
\sphinxAtStartPar
Click the \sphinxstylestrong{Share} button (link icon) in the toolbar.

\item {} 
\sphinxAtStartPar
A permalink URL is generated. Copy and share this URL — recipients will open the map
in the same state.

\end{enumerate}


\subsubsection{Embed}
\label{\detokenize{en/04_exploring_maps:embed}}\begin{enumerate}
\sphinxsetlistlabels{\arabic}{enumi}{enumii}{}{.}%
\item {} 
\sphinxAtStartPar
Click the \sphinxstylestrong{Share} button.

\item {} 
\sphinxAtStartPar
Switch to the \sphinxstylestrong{Embed} tab.

\item {} 
\sphinxAtStartPar
Copy the HTML \sphinxcode{\sphinxupquote{\textless{}iframe\textgreater{}}} code and paste it into your website or report.

\end{enumerate}


\subsection{Saving Changes to a Map}
\label{\detokenize{en/04_exploring_maps:saving-changes-to-a-map}}
\sphinxAtStartPar
If you have \sphinxstylestrong{edit permission} on a Map, you can save changes to layers, styles, and
the initial view:
\begin{enumerate}
\sphinxsetlistlabels{\arabic}{enumi}{enumii}{}{.}%
\item {} 
\sphinxAtStartPar
Make your changes in the map viewer (reorder layers, change styles, set a new
initial extent, etc.).

\item {} 
\sphinxAtStartPar
Click the \sphinxstylestrong{Save} button (💾 icon) in the toolbar.

\item {} 
\sphinxAtStartPar
Choose \sphinxstylestrong{Save} (overwrite the existing map) or \sphinxstylestrong{Save as} (create a new map).

\end{enumerate}

\begin{sphinxadmonition}{warning}{Warning:}
\sphinxAtStartPar
Saving overwrites the map for all users who have access to it.
If you want to experiment without affecting others, use \sphinxstylestrong{Save as} to create a
personal copy first.
\end{sphinxadmonition}


\subsection{Creating a New Map}
\label{\detokenize{en/04_exploring_maps:creating-a-new-map}}
\sphinxAtStartPar
Members of the Pelagos Agreement group with appropriate permissions can create new Maps:
\begin{enumerate}
\sphinxsetlistlabels{\arabic}{enumi}{enumii}{}{.}%
\item {} 
\sphinxAtStartPar
Go to \sphinxstylestrong{Maps} in the navigation bar.

\item {} 
\sphinxAtStartPar
Click \sphinxstylestrong{+ Create map} (or the equivalent button at the top of the catalogue).

\item {} 
\sphinxAtStartPar
The map viewer opens with an empty canvas.

\item {} 
\sphinxAtStartPar
Use the \sphinxstylestrong{Add layer} button (+ icon in the Layer panel) to search for and add
Datasets from the KMaP catalogue.

\item {} 
\sphinxAtStartPar
Style each layer as needed.

\item {} 
\sphinxAtStartPar
Click \sphinxstylestrong{Save} and provide a title, abstract, and keywords
(remember to include \sphinxcode{\sphinxupquote{pelagos sanctuary}} or \sphinxcode{\sphinxupquote{pelagos agreement}}).

\end{enumerate}

\sphinxstepscope


\chapter{KMaP \textendash{} Knowledge Management Platform}
\label{\detokenize{it/index:kmap-knowledge-management-platform}}\label{\detokenize{it/index::doc}}
\sphinxAtStartPar
\sphinxstylestrong{Documentazione utente per il gruppo Accordo Pelagos}

\begin{sphinxadmonition}{note}{Note:}
\sphinxAtStartPar
Questa documentazione è rivolta principalmente agli utenti di Kmap che sono membri del gruppo
\sphinxstylestrong{Accordo Pelagos} sulla piattaforma KMaP. Alcuni elementi descritti qui
potrebbero essere mostrati in maniera differente agli utenti che non appartengono a questo gruppo.
\end{sphinxadmonition}

\sphinxAtStartPar
La piattaforma KMaP è accessibile all’indirizzo: \sphinxurl{https://kmap.info-rac.org}

\sphinxstepscope


\section{Benvenuto su KMaP}
\label{\detokenize{it/01_benvenuto:benvenuto-su-kmap}}\label{\detokenize{it/01_benvenuto:welcome-it}}\label{\detokenize{it/01_benvenuto::doc}}

\subsection{Presentazione di KMaP}
\label{\detokenize{it/01_benvenuto:presentazione-di-kmap}}
\sphinxAtStartPar
\sphinxstylestrong{KMaP} (Knowledge Management Platform) è l’Infrastruttura di Dati Spaziali sviluppata
da \sphinxstylestrong{INFO/RAC}, il Centro di Informazione e Comunicazione INFO/RAC del
Piano d’Azione per il Mediterraneo \sphinxhref{https://www.unep.org/unepmap/}{(Mediteranean Action Plan \sphinxhyphen{} MAP)}
del Programma Ambiente delle Nazioni Unite (United Nation Environment Program \sphinxhyphen{} UNEP)..

\sphinxAtStartPar
KMaP è basata su \sphinxhref{https://geonode.org/}{GeoNode}, una piattaforma web open\sphinxhyphen{}source per
la gestione e la condivisione di dati geografici.
Costituisce il punto di accesso principale al patrimonio di dati ambientali e spaziali del Mediterraneo.

\sphinxAtStartPar
La piattaforma è suddivisa su tre aree:


\begin{savenotes}\sphinxattablestart
\sphinxthistablewithglobalstyle
\centering
\begin{tabular}[t]{\X{20}{100}\X{15}{100}\X{65}{100}}
\sphinxtoprule
\sphinxstyletheadfamily 
\sphinxAtStartPar
Hub
&\sphinxstyletheadfamily 
\sphinxAtStartPar
Accesso
&\sphinxstyletheadfamily 
\sphinxAtStartPar
Descrizione
\\
\sphinxmidrule
\sphinxtableatstartofbodyhook
\sphinxAtStartPar
\sphinxstylestrong{Data Hub}
&
\sphinxAtStartPar
{\color{red}\bfseries{}|mapsbut|} \sphinxstyleemphasis{Maps}
&
\sphinxAtStartPar
Dati geografici: dataset, mappe interattive, dashboard e geostorie.
\\
\sphinxhline
\sphinxAtStartPar
\sphinxstylestrong{Knowledge Hub}
&
\sphinxAtStartPar
{\color{red}\bfseries{}|librarybut|} \sphinxstyleemphasis{Library}
&
\sphinxAtStartPar
Patrimonio documentale: rapporti, pubblicazioni e documenti di riferimento.
\\
\sphinxhline
\sphinxAtStartPar
\sphinxstylestrong{Exchange Hub}
&
\sphinxAtStartPar
{\color{red}\bfseries{}|networkbut|} \sphinxstyleemphasis{Network}
&
\sphinxAtStartPar
Strumenti di collaborazione e scambio di conoscenze.
\\
\sphinxbottomrule
\end{tabular}
\sphinxtableafterendhook\par
\sphinxattableend\end{savenotes}

\begin{sphinxadmonition}{note}{Note:}
\sphinxAtStartPar
In quanto membri del gruppo \sphinxstylestrong{Accordo Pelagos}, avrete accesso a risorse specifiche,
diritti di modifica dei metadati e autorizzazioni a livello di gruppo non disponibili
a tutti gli utenti.
\end{sphinxadmonition}


\subsection{Tipi di Risorse}
\label{\detokenize{it/01_benvenuto:tipi-di-risorse}}
\sphinxAtStartPar
KMaP organizza tutti i contenuti in cinque tipi di risorse, ciascuno con uno scopo distinto.


\subsubsection{Dataset}
\label{\detokenize{it/01_benvenuto:dataset}}
\sphinxAtStartPar
Un \sphinxstylestrong{Dataset} è uno strato di dati geografici che rappresenta elementi o misurazioni
del mondo reale. È il mattone fondamentale della piattaforma — la materia prima con cui
si costruiscono mappe, dashboard e analisi.

\sphinxAtStartPar
I dataset possono essere:
\begin{itemize}
\item {} 
\sphinxAtStartPar
\sphinxstylestrong{Vettoriali} \textendash{} elementi geografici discreti rappresentati come punti, linee o poligoni
(es. confini di aree protette, reti fluviali, localizzazione di stazioni di monitoraggio).

\item {} 
\sphinxAtStartPar
\sphinxstylestrong{Raster} \textendash{} superfici continue rappresentate come griglie di celle
(es. temperatura superficiale del mare, immagini satellitari, modelli batimetrici).

\end{itemize}

\sphinxAtStartPar
Con un Dataset puoi:
\begin{itemize}
\item {} 
\sphinxAtStartPar
Visualizzarlo interattivamente su una mappa web

\item {} 
\sphinxAtStartPar
Consultare la tabella degli attributi

\item {} 
\sphinxAtStartPar
Scaricarlo nei formati GIS standard (Shapefile, GeoTIFF, GeoJSON, KML, CSV, …)

\item {} 
\sphinxAtStartPar
Leggere e modificare i suoi metadati

\item {} 
\sphinxAtStartPar
Aggiungerlo come strato a una Mappa

\end{itemize}

\begin{sphinxadmonition}{tip}{Tip:}
\sphinxAtStartPar
I dataset pubblicati per l’Accordo Pelagos sono etichettati con le parole chiave
\sphinxcode{\sphinxupquote{pelagos sanctuary}} o \sphinxcode{\sphinxupquote{pelagos agreement}}. Usa queste parole chiave per filtrare
rapidamente il catalogo (vedi {\hyperref[\detokenize{it/02_trovare_visualizzare:find-it}]{\sphinxcrossref{\DUrole{std}{\DUrole{std-ref}{Come Trovare e Visualizzare i Tuoi Dati}}}}}).
\end{sphinxadmonition}


\subsubsection{Document}
\label{\detokenize{it/01_benvenuto:document}}
\sphinxAtStartPar
Un \sphinxstylestrong{Document} è qualsiasi file caricato che integra i dati spaziali: un rapporto PDF,
una fotografia, un foglio di calcolo, una presentazione o qualsiasi altro materiale
di riferimento.

\sphinxAtStartPar
I documenti sono risorse a pieno titolo in KMaP — hanno la propria pagina di metadati,
possono essere cercati e filtrati, e possono essere collegati a Dataset o Mappe.

\sphinxAtStartPar
I formati supportati includono: \sphinxcode{\sphinxupquote{.pdf}}, \sphinxcode{\sphinxupquote{.docx}}, \sphinxcode{\sphinxupquote{.xlsx}}, \sphinxcode{\sphinxupquote{.pptx}}, \sphinxcode{\sphinxupquote{.jpg}},
\sphinxcode{\sphinxupquote{.png}}, \sphinxcode{\sphinxupquote{.tif}}, \sphinxcode{\sphinxupquote{.mp4}}, \sphinxcode{\sphinxupquote{.zip}} e molti altri.


\subsubsection{Mappa}
\label{\detokenize{it/01_benvenuto:mappa}}
\sphinxAtStartPar
Una \sphinxstylestrong{Mappa} è una mappa web interattiva a più strati, creata componendo uno o più
Dataset. Le mappe sono create e visualizzate tramite il visualizzatore integrato \sphinxstylestrong{MapStore}.

\sphinxAtStartPar
Ogni Mappa:
\begin{itemize}
\item {} 
\sphinxAtStartPar
Combina più strati di Dataset con stili e impostazioni di visibilità individuali

\item {} 
\sphinxAtStartPar
Include una mappa di sfondo (es. OpenStreetMap o immagini satellitari)

\item {} 
\sphinxAtStartPar
Supporta strumenti di navigazione, interrogazione, misurazione e stampa

\item {} 
\sphinxAtStartPar
Può essere condivisa tramite un link permanente o incorporata in siti web esterni

\end{itemize}


\subsubsection{GeoStory}
\label{\detokenize{it/01_benvenuto:geostory}}
\sphinxAtStartPar
Una \sphinxstylestrong{GeoStory} è una narrazione multimediale scorrevole che intreccia testo, mappe
interattive, immagini e video — simile a una story map.

\sphinxAtStartPar
Le GeoStory sono composte da sezioni:


\begin{savenotes}\sphinxattablestart
\sphinxthistablewithglobalstyle
\centering
\begin{tabular}[t]{\X{25}{100}\X{75}{100}}
\sphinxtoprule
\sphinxstyletheadfamily 
\sphinxAtStartPar
Tipo di sezione
&\sphinxstyletheadfamily 
\sphinxAtStartPar
Descrizione
\\
\sphinxmidrule
\sphinxtableatstartofbodyhook
\sphinxAtStartPar
\sphinxstylestrong{Titolo}
&
\sphinxAtStartPar
Pannello iniziale a schermo intero con immagine o video di sfondo.
\\
\sphinxhline
\sphinxAtStartPar
\sphinxstylestrong{Banner}
&
\sphinxAtStartPar
Divisore orizzontale o intestazione di capitolo con testo e/o immagine.
\\
\sphinxhline
\sphinxAtStartPar
\sphinxstylestrong{Paragrafo}
&
\sphinxAtStartPar
Testo libero e contenuto multimediale (immagine, video, mappa).
\\
\sphinxhline
\sphinxAtStartPar
\sphinxstylestrong{Immersivo}
&
\sphinxAtStartPar
Mappa interattiva a schermo intero su cui scorre il testo narrativo.
\\
\sphinxhline
\sphinxAtStartPar
\sphinxstylestrong{GeoCarousel}
&
\sphinxAtStartPar
Presentazione a diapositive collegata a posizioni geografiche.
\\
\sphinxhline
\sphinxAtStartPar
\sphinxstylestrong{Pagina web}
&
\sphinxAtStartPar
Pagina web esterna incorporata direttamente nella storia.
\\
\sphinxbottomrule
\end{tabular}
\sphinxtableafterendhook\par
\sphinxattableend\end{savenotes}

\sphinxAtStartPar
Le GeoStory sono ideali per comunicare tematiche ambientali a un pubblico non specialista.


\subsubsection{Dashboard}
\label{\detokenize{it/01_benvenuto:dashboard}}
\sphinxAtStartPar
Una \sphinxstylestrong{Dashboard} è un pannello configurabile di widget interattivi — grafici, tabelle,
mappe, contatori e blocchi di testo — collegati a uno o più Dataset.

\sphinxAtStartPar
Tipi di widget disponibili:


\begin{savenotes}\sphinxattablestart
\sphinxthistablewithglobalstyle
\centering
\begin{tabular}[t]{\X{25}{100}\X{75}{100}}
\sphinxtoprule
\sphinxstyletheadfamily 
\sphinxAtStartPar
Widget
&\sphinxstyletheadfamily 
\sphinxAtStartPar
Descrizione
\\
\sphinxmidrule
\sphinxtableatstartofbodyhook
\sphinxAtStartPar
\sphinxstylestrong{Mappa}
&
\sphinxAtStartPar
Mini\sphinxhyphen{}mappa interattiva che visualizza strati di dataset.
\\
\sphinxhline
\sphinxAtStartPar
\sphinxstylestrong{Grafico}
&
\sphinxAtStartPar
Grafici a barre, linee, torta o dispersione derivati dagli attributi del dataset.
\\
\sphinxhline
\sphinxAtStartPar
\sphinxstylestrong{Tabella}
&
\sphinxAtStartPar
Vista tabulare ordinabile e filtrabile degli attributi del dataset.
\\
\sphinxhline
\sphinxAtStartPar
\sphinxstylestrong{Contatore}
&
\sphinxAtStartPar
Statistica riassuntiva in grande formato (conteggio, somma, media, …).
\\
\sphinxhline
\sphinxAtStartPar
\sphinxstylestrong{Testo}
&
\sphinxAtStartPar
Blocco di testo libero per titoli o commenti.
\\
\sphinxbottomrule
\end{tabular}
\sphinxtableafterendhook\par
\sphinxattableend\end{savenotes}

\sphinxAtStartPar
I widget in una Dashboard sono interconnessi: selezionando una regione sul widget mappa,
tutti gli altri widget si aggiornano automaticamente per mostrare solo i dati di
quell’area.


\subsection{Servizi Interoperabili}
\label{\detokenize{it/01_benvenuto:servizi-interoperabili}}
\sphinxAtStartPar
KMaP espone i propri dati attraverso \sphinxstylestrong{servizi web OGC} standard, permettendo a qualsiasi
software GIS compatibile o applicazione web di accedere direttamente ai dati senza
scaricare file.


\begin{savenotes}\sphinxattablestart
\sphinxthistablewithglobalstyle
\centering
\begin{tabular}[t]{\X{15}{100}\X{15}{100}\X{70}{100}}
\sphinxtoprule
\sphinxstyletheadfamily 
\sphinxAtStartPar
Servizio
&\sphinxstyletheadfamily 
\sphinxAtStartPar
Standard
&\sphinxstyletheadfamily 
\sphinxAtStartPar
Descrizione
\\
\sphinxmidrule
\sphinxtableatstartofbodyhook
\sphinxAtStartPar
\sphinxstylestrong{WMS}
&
\sphinxAtStartPar
OGC Web Map Service
&
\sphinxAtStartPar
Restituisce strati di dataset come immagini di mappa renderizzate. Ideale per la
visualizzazione in GIS desktop (QGIS, ArcGIS) o applicazioni web.
\\
\sphinxhline
\sphinxAtStartPar
\sphinxstylestrong{WFS}
&
\sphinxAtStartPar
OGC Web Feature Service
&
\sphinxAtStartPar
Restituisce dati vettoriali grezzi (geometrie + attributi) per il download o
l’analisi. Supporta il filtraggio per attributo o bounding box.
\\
\sphinxhline
\sphinxAtStartPar
\sphinxstylestrong{WCS}
&
\sphinxAtStartPar
OGC Web Coverage Service
&
\sphinxAtStartPar
Restituisce dati raster grezzi (griglie) per l’analisi. Disponibile per dataset raster.
\\
\sphinxhline
\sphinxAtStartPar
\sphinxstylestrong{WMTS}
&
\sphinxAtStartPar
OGC Web Map Tile Service
&
\sphinxAtStartPar
Restituisce tile di mappa pre\sphinxhyphen{}renderizzati per la visualizzazione rapida in mappe web.
\\
\sphinxhline
\sphinxAtStartPar
\sphinxstylestrong{CSW}
&
\sphinxAtStartPar
OGC Catalogue Service
&
\sphinxAtStartPar
Espone il catalogo dei metadati di KMaP per la ricerca e la raccolta da parte di
portali esterni.
\\
\sphinxbottomrule
\end{tabular}
\sphinxtableafterendhook\par
\sphinxattableend\end{savenotes}
\subsubsection*{Come ottenere gli URL dei servizi}

\sphinxAtStartPar
Nella pagina di dettaglio di qualsiasi Dataset, fai clic sulla scheda \sphinxstylestrong{Servizi OGC}
(o cerca la sezione \sphinxstyleemphasis{Download / OGC}). Gli URL degli endpoint WMS, WFS e WCS sono
elencati lì e possono essere copiati direttamente in QGIS o in qualsiasi altro client
compatibile con OGC.

\begin{sphinxadmonition}{tip}{Tip:}
\sphinxAtStartPar
In QGIS, vai su \sphinxstylestrong{Layer \(\rightarrow\) Aggiungi layer \(\rightarrow\) Aggiungi layer WMS/WMTS}, incolla l’URL
WMS e sfoglia l’elenco dei layer disponibili per aggiungere direttamente i dati KMaP
al tuo progetto senza scaricare alcun file.
\end{sphinxadmonition}


\subsection{Tematiche Disponibili su KMaP}
\label{\detokenize{it/01_benvenuto:tematiche-disponibili-su-kmap}}
\sphinxAtStartPar
KMaP organizza i suoi dataset nelle seguenti categorie tematiche:
\begin{itemize}
\item {} 
\sphinxAtStartPar
🐟  \sphinxstylestrong{Pesca e Acquacoltura}

\item {} 
\sphinxAtStartPar
🌊  \sphinxstylestrong{Biodiversità Marina}

\item {} 
\sphinxAtStartPar
🏭  \sphinxstylestrong{Inquinamento}

\item {} 
\sphinxAtStartPar
🌡  \sphinxstylestrong{Cambiamenti Climatici}

\item {} 
\sphinxAtStartPar
🗺  \sphinxstylestrong{Pianificazione Spaziale Marina}

\item {} 
\sphinxAtStartPar
💙  \sphinxstylestrong{Sostenibilità ed Economia Blu}

\item {} 
\sphinxAtStartPar
🏛  \sphinxstylestrong{Governance}

\end{itemize}

\sphinxAtStartPar
Le risorse relative al Santuario Pelagos e all’Accordo Pelagos sono distribuite tra
più tematiche, in particolare Biodiversità Marina, Governance e Pianificazione Spaziale
Marina.

\sphinxstepscope


\section{Come Trovare e Visualizzare i Tuoi Dati}
\label{\detokenize{it/02_trovare_visualizzare:come-trovare-e-visualizzare-i-tuoi-dati}}\label{\detokenize{it/02_trovare_visualizzare:find-it}}\label{\detokenize{it/02_trovare_visualizzare::doc}}
\begin{sphinxcontents}
\sphinxstylecontentstitle{In questa pagina}
\begin{itemize}
\item {} 
\sphinxAtStartPar
\phantomsection\label{\detokenize{it/02_trovare_visualizzare:id1}}{\hyperref[\detokenize{it/02_trovare_visualizzare:accedere-al-catalogo}]{\sphinxcrossref{Accedere al Catalogo}}}

\item {} 
\sphinxAtStartPar
\phantomsection\label{\detokenize{it/02_trovare_visualizzare:id2}}{\hyperref[\detokenize{it/02_trovare_visualizzare:filtrare-e-cercare-nel-catalogo}]{\sphinxcrossref{Filtrare e Cercare nel Catalogo}}}
\begin{itemize}
\item {} 
\sphinxAtStartPar
\phantomsection\label{\detokenize{it/02_trovare_visualizzare:id3}}{\hyperref[\detokenize{it/02_trovare_visualizzare:usare-la-barra-di-ricerca}]{\sphinxcrossref{Usare la Barra di Ricerca}}}

\item {} 
\sphinxAtStartPar
\phantomsection\label{\detokenize{it/02_trovare_visualizzare:id4}}{\hyperref[\detokenize{it/02_trovare_visualizzare:usare-i-filtri}]{\sphinxcrossref{Usare i Filtri}}}

\end{itemize}

\item {} 
\sphinxAtStartPar
\phantomsection\label{\detokenize{it/02_trovare_visualizzare:id5}}{\hyperref[\detokenize{it/02_trovare_visualizzare:esempio-guidato-trovare-un-dataset-pelagos}]{\sphinxcrossref{Esempio Guidato: Trovare un Dataset Pelagos}}}
\begin{itemize}
\item {} 
\sphinxAtStartPar
\phantomsection\label{\detokenize{it/02_trovare_visualizzare:id6}}{\hyperref[\detokenize{it/02_trovare_visualizzare:passo-1-aprire-il-catalogo}]{\sphinxcrossref{Passo 1 \textendash{} Aprire il catalogo}}}

\item {} 
\sphinxAtStartPar
\phantomsection\label{\detokenize{it/02_trovare_visualizzare:id7}}{\hyperref[\detokenize{it/02_trovare_visualizzare:passo-2-cercare-per-parola-chiave}]{\sphinxcrossref{Passo 2 \textendash{} Cercare per parola chiave}}}

\item {} 
\sphinxAtStartPar
\phantomsection\label{\detokenize{it/02_trovare_visualizzare:id8}}{\hyperref[\detokenize{it/02_trovare_visualizzare:passo-3-filtrare-solo-i-dataset}]{\sphinxcrossref{Passo 3 \textendash{} Filtrare solo i Dataset}}}

\item {} 
\sphinxAtStartPar
\phantomsection\label{\detokenize{it/02_trovare_visualizzare:id9}}{\hyperref[\detokenize{it/02_trovare_visualizzare:passo-4-selezionare-un-dataset}]{\sphinxcrossref{Passo 4 \textendash{} Selezionare un Dataset}}}

\end{itemize}

\item {} 
\sphinxAtStartPar
\phantomsection\label{\detokenize{it/02_trovare_visualizzare:id10}}{\hyperref[\detokenize{it/02_trovare_visualizzare:la-pagina-di-dettaglio-del-dataset}]{\sphinxcrossref{La Pagina di Dettaglio del Dataset}}}
\begin{itemize}
\item {} 
\sphinxAtStartPar
\phantomsection\label{\detokenize{it/02_trovare_visualizzare:id11}}{\hyperref[\detokenize{it/02_trovare_visualizzare:miniatura-e-azioni}]{\sphinxcrossref{Miniatura e Azioni}}}

\item {} 
\sphinxAtStartPar
\phantomsection\label{\detokenize{it/02_trovare_visualizzare:id12}}{\hyperref[\detokenize{it/02_trovare_visualizzare:anteprima-interattiva}]{\sphinxcrossref{Anteprima Interattiva}}}

\item {} 
\sphinxAtStartPar
\phantomsection\label{\detokenize{it/02_trovare_visualizzare:id13}}{\hyperref[\detokenize{it/02_trovare_visualizzare:pannello-dei-metadati}]{\sphinxcrossref{Pannello dei Metadati}}}

\item {} 
\sphinxAtStartPar
\phantomsection\label{\detokenize{it/02_trovare_visualizzare:id14}}{\hyperref[\detokenize{it/02_trovare_visualizzare:visualizzazione-completa-dei-metadati}]{\sphinxcrossref{Visualizzazione Completa dei Metadati}}}

\item {} 
\sphinxAtStartPar
\phantomsection\label{\detokenize{it/02_trovare_visualizzare:id15}}{\hyperref[\detokenize{it/02_trovare_visualizzare:scaricare-un-dataset}]{\sphinxcrossref{Scaricare un Dataset}}}

\end{itemize}

\end{itemize}
\end{sphinxcontents}

\sphinxAtStartPar
Questa sezione spiega come scoprire le risorse in KMaP, applicare filtri per restringere
la ricerca ed esplorare la pagina di dettaglio e l’anteprima di un Dataset.


\subsection{Accedere al Catalogo}
\label{\detokenize{it/02_trovare_visualizzare:accedere-al-catalogo}}
\sphinxAtStartPar
Il catalogo di KMaP elenca tutte le risorse pubbliche e tutte le risorse condivise con
il tuo gruppo. Esistono due modi principali per raggiungerlo:
\begin{enumerate}
\sphinxsetlistlabels{\arabic}{enumi}{enumii}{}{.}%
\item {} 
\sphinxAtStartPar
\sphinxstylestrong{Dalla barra di navigazione} \textendash{} clicca su \sphinxstylestrong{Maps} (per le risorse spaziali) o
\sphinxstylestrong{Library} (per documenti e pubblicazioni) nel menu in alto.

\item {} 
\sphinxAtStartPar
\sphinxstylestrong{Direttamente tramite URL} \textendash{} vai su \sphinxurl{https://kmap.info-rac.org/resources/}

\end{enumerate}

\sphinxAtStartPar
La pagina del catalogo mostra le risorse come schede, ciascuna con miniatura, titolo,
tipo di risorsa e metadati riassuntivi.

\begin{sphinxadmonition}{tip}{Tip:}
\sphinxAtStartPar
Assicurati di essere \sphinxstylestrong{connesso} per vedere le risorse condivise esclusivamente con
il gruppo Accordo Pelagos. Le risorse riservate al gruppo non sono visibili agli
utenti non autenticati.
\end{sphinxadmonition}


\subsection{Filtrare e Cercare nel Catalogo}
\label{\detokenize{it/02_trovare_visualizzare:filtrare-e-cercare-nel-catalogo}}

\subsubsection{Usare la Barra di Ricerca}
\label{\detokenize{it/02_trovare_visualizzare:usare-la-barra-di-ricerca}}
\sphinxAtStartPar
La barra di ricerca in cima alla pagina del catalogo accetta query in testo libero.
Effettua la ricerca su titoli, abstract e parole chiave.

\sphinxAtStartPar
Per trovare contenuti relativi a Pelagos:
\begin{enumerate}
\sphinxsetlistlabels{\arabic}{enumi}{enumii}{}{.}%
\item {} 
\sphinxAtStartPar
Digita \sphinxcode{\sphinxupquote{pelagos sanctuary}} nella barra di ricerca e premi \sphinxstylestrong{Invio}.

\item {} 
\sphinxAtStartPar
In alternativa, cerca \sphinxcode{\sphinxupquote{pelagos agreement}} per recuperare le risorse etichettate
con quella parola chiave.

\end{enumerate}

\begin{sphinxadmonition}{note}{Note:}
\sphinxAtStartPar
Sia \sphinxcode{\sphinxupquote{pelagos sanctuary}} che \sphinxcode{\sphinxupquote{pelagos agreement}} sono parole chiave controllate
utilizzate dal gruppo Accordo Pelagos per pubblicare risorse su KMaP. Usando questi
termini esatti otterrai i risultati più pertinenti.
\end{sphinxadmonition}


\subsubsection{Usare i Filtri}
\label{\detokenize{it/02_trovare_visualizzare:usare-i-filtri}}
\sphinxAtStartPar
Sul lato sinistro della pagina del catalogo si trova un pannello di filtri. I filtri
possono essere combinati liberamente e vengono applicati immediatamente.


\begin{savenotes}\sphinxattablestart
\sphinxthistablewithglobalstyle
\centering
\begin{tabular}[t]{\X{25}{100}\X{75}{100}}
\sphinxtoprule
\sphinxstyletheadfamily 
\sphinxAtStartPar
Filtro
&\sphinxstyletheadfamily 
\sphinxAtStartPar
Come utilizzarlo
\\
\sphinxmidrule
\sphinxtableatstartofbodyhook
\sphinxAtStartPar
\sphinxstylestrong{Tipo di risorsa}
&
\sphinxAtStartPar
Seleziona \sphinxstyleemphasis{Dataset}, \sphinxstyleemphasis{Document}, \sphinxstyleemphasis{Map}, \sphinxstyleemphasis{GeoStory} o \sphinxstyleemphasis{Dashboard} per limitare
i risultati a un solo tipo.
\\
\sphinxhline
\sphinxAtStartPar
\sphinxstylestrong{Parole chiave}
&
\sphinxAtStartPar
Digita o seleziona una parola chiave (es. \sphinxcode{\sphinxupquote{pelagos sanctuary}}) per filtrare
per tag tematico.
\\
\sphinxhline
\sphinxAtStartPar
\sphinxstylestrong{Categoria / Tema}
&
\sphinxAtStartPar
Seleziona una categoria tematica come \sphinxstyleemphasis{Biodiversità Marina} o \sphinxstyleemphasis{Governance}.
\\
\sphinxhline
\sphinxAtStartPar
\sphinxstylestrong{Data}
&
\sphinxAtStartPar
Filtra per intervallo di date di pubblicazione o modifica.
\\
\sphinxhline
\sphinxAtStartPar
\sphinxstylestrong{Proprietario}
&
\sphinxAtStartPar
Filtra per l’utente o l’organizzazione che ha pubblicato la risorsa.
\\
\sphinxhline
\sphinxAtStartPar
\sphinxstylestrong{Estensione}
&
\sphinxAtStartPar
Disegna un riquadro sulla mini\sphinxhyphen{}mappa per filtrare per area geografica.
\\
\sphinxbottomrule
\end{tabular}
\sphinxtableafterendhook\par
\sphinxattableend\end{savenotes}

\begin{sphinxadmonition}{tip}{Tip:}
\sphinxAtStartPar
Per trovare rapidamente tutti i dataset relativi a Pelagos:
\begin{enumerate}
\sphinxsetlistlabels{\arabic}{enumi}{enumii}{}{.}%
\item {} 
\sphinxAtStartPar
Imposta \sphinxstylestrong{Tipo di risorsa} \(\rightarrow\) \sphinxstyleemphasis{Dataset}

\item {} 
\sphinxAtStartPar
Nel filtro \sphinxstylestrong{Parole chiave}, digita \sphinxcode{\sphinxupquote{pelagos sanctuary}}

\item {} 
\sphinxAtStartPar
Il catalogo si aggiornerà mostrando solo i dataset etichettati con quella parola chiave.

\end{enumerate}
\end{sphinxadmonition}


\subsection{Esempio Guidato: Trovare un Dataset Pelagos}
\label{\detokenize{it/02_trovare_visualizzare:esempio-guidato-trovare-un-dataset-pelagos}}
\sphinxAtStartPar
Il seguente esempio mostra come trovare, visualizzare in anteprima e leggere i metadati
di un Dataset rilevante per l’Accordo Pelagos.


\subsubsection{Passo 1 \textendash{} Aprire il catalogo}
\label{\detokenize{it/02_trovare_visualizzare:passo-1-aprire-il-catalogo}}
\sphinxAtStartPar
Vai su \sphinxurl{https://kmap.info-rac.org} e clicca \sphinxstylestrong{Maps} nella barra di navigazione in alto.
Si apre il catalogo del Data Hub, filtrato per tipi di risorse spaziali.


\subsubsection{Passo 2 \textendash{} Cercare per parola chiave}
\label{\detokenize{it/02_trovare_visualizzare:passo-2-cercare-per-parola-chiave}}
\sphinxAtStartPar
Nella barra di ricerca, digita:

\begin{sphinxVerbatim}[commandchars=\\\{\}]
\PYG{n}{pelagos} \PYG{n}{sanctuary}
\end{sphinxVerbatim}

\sphinxAtStartPar
Premi \sphinxstylestrong{Invio}. Il pannello dei risultati si aggiorna mostrando le risorse corrispondenti
a quella parola chiave.


\subsubsection{Passo 3 \textendash{} Filtrare solo i Dataset}
\label{\detokenize{it/02_trovare_visualizzare:passo-3-filtrare-solo-i-dataset}}
\sphinxAtStartPar
Nel pannello di filtri a sinistra, sotto \sphinxstylestrong{Tipo di risorsa}, spunta \sphinxstylestrong{Dataset}.
L’elenco ora mostra solo i dataset etichettati con \sphinxstyleemphasis{pelagos sanctuary}.


\subsubsection{Passo 4 \textendash{} Selezionare un Dataset}
\label{\detokenize{it/02_trovare_visualizzare:passo-4-selezionare-un-dataset}}
\sphinxAtStartPar
Clicca su una scheda dataset per aprire la relativa \sphinxstylestrong{pagina di dettaglio}.


\subsection{La Pagina di Dettaglio del Dataset}
\label{\detokenize{it/02_trovare_visualizzare:la-pagina-di-dettaglio-del-dataset}}
\sphinxAtStartPar
La pagina di dettaglio è la pagina informativa centrale di qualsiasi risorsa. È divisa
in diverse sezioni.


\subsubsection{Miniatura e Azioni}
\label{\detokenize{it/02_trovare_visualizzare:miniatura-e-azioni}}
\sphinxAtStartPar
In cima alla pagina troverai:
\begin{itemize}
\item {} 
\sphinxAtStartPar
Una \sphinxstylestrong{miniatura mappa} o immagine di anteprima del dataset.

\item {} 
\sphinxAtStartPar
Pulsanti di azione:
\begin{itemize}
\item {} 
\sphinxAtStartPar
\sphinxstylestrong{View map} \textendash{} apre il dataset nel visualizzatore mappa interattivo.

\item {} 
\sphinxAtStartPar
\sphinxstylestrong{Download} \textendash{} scarica il dataset nel formato scelto.

\item {} 
\sphinxAtStartPar
\sphinxstylestrong{Edit} (se hai l’autorizzazione) \textendash{} apre l’editor dei metadati.

\item {} 
\sphinxAtStartPar
\sphinxstylestrong{Share} \textendash{} copia il link permanente alla risorsa.

\end{itemize}

\end{itemize}


\subsubsection{Anteprima Interattiva}
\label{\detokenize{it/02_trovare_visualizzare:anteprima-interattiva}}
\sphinxAtStartPar
Sotto la miniatura, il \sphinxstylestrong{pannello di anteprima mappa} mostra il dataset visualizzato
su una mappa di sfondo. Puoi:
\begin{itemize}
\item {} 
\sphinxAtStartPar
Spostarti e zoomare con il mouse o il touch.

\item {} 
\sphinxAtStartPar
Cliccare su un elemento per aprire un \sphinxstylestrong{pop\sphinxhyphen{}up} con i suoi valori di attributo.

\item {} 
\sphinxAtStartPar
Usare il pannello dei layer (icona ☰) per attivare/disattivare la visibilità o
consultare la legenda.

\end{itemize}

\begin{sphinxadmonition}{note}{Note:}
\sphinxAtStartPar
L’anteprima è una mappa interattiva — non un’immagine statica. Puoi esplorare i
dati nello spazio prima di decidere se scaricarli o utilizzarli.
\end{sphinxadmonition}


\subsubsection{Pannello dei Metadati}
\label{\detokenize{it/02_trovare_visualizzare:pannello-dei-metadati}}
\sphinxAtStartPar
Sotto l’anteprima, il \sphinxstylestrong{pannello dei metadati} contiene informazioni dettagliate sul
dataset. I campi principali sono:


\begin{savenotes}\sphinxattablestart
\sphinxthistablewithglobalstyle
\centering
\begin{tabular}[t]{\X{30}{100}\X{70}{100}}
\sphinxtoprule
\sphinxstyletheadfamily 
\sphinxAtStartPar
Campo
&\sphinxstyletheadfamily 
\sphinxAtStartPar
Descrizione
\\
\sphinxmidrule
\sphinxtableatstartofbodyhook
\sphinxAtStartPar
\sphinxstylestrong{Titolo}
&
\sphinxAtStartPar
Il nome completo del dataset.
\\
\sphinxhline
\sphinxAtStartPar
\sphinxstylestrong{Abstract}
&
\sphinxAtStartPar
Una descrizione in linguaggio semplice del contenuto e dello scopo del dataset.
\\
\sphinxhline
\sphinxAtStartPar
\sphinxstylestrong{Parole chiave}
&
\sphinxAtStartPar
Tag tematici e geografici (cerca \sphinxcode{\sphinxupquote{pelagos sanctuary}} o \sphinxcode{\sphinxupquote{pelagos agreement}}).
\\
\sphinxhline
\sphinxAtStartPar
\sphinxstylestrong{Categoria}
&
\sphinxAtStartPar
La categoria INSPIRE o tematica (es. \sphinxstyleemphasis{Biodiversità Marina}).
\\
\sphinxhline
\sphinxAtStartPar
\sphinxstylestrong{Data}
&
\sphinxAtStartPar
Data di pubblicazione e ultima data di modifica.
\\
\sphinxhline
\sphinxAtStartPar
\sphinxstylestrong{Estensione spaziale}
&
\sphinxAtStartPar
Riquadro di delimitazione o regione geografica coperta dai dati.
\\
\sphinxhline
\sphinxAtStartPar
\sphinxstylestrong{Sistema di Riferimento delle Coordinate (CRS)}
&
\sphinxAtStartPar
La proiezione utilizzata (es. WGS 84 / EPSG:4326).
\\
\sphinxhline
\sphinxAtStartPar
\sphinxstylestrong{Proprietario}
&
\sphinxAtStartPar
L’utente o l’organizzazione responsabile del dataset.
\\
\sphinxhline
\sphinxAtStartPar
\sphinxstylestrong{Licenza}
&
\sphinxAtStartPar
La licenza d’uso (es. Creative Commons, ODbL, …).
\\
\sphinxhline
\sphinxAtStartPar
\sphinxstylestrong{Punto di contatto}
&
\sphinxAtStartPar
Chi contattare per domande sui dati.
\\
\sphinxhline
\sphinxAtStartPar
\sphinxstylestrong{Servizi OGC}
&
\sphinxAtStartPar
URL degli endpoint WMS, WFS e WCS per l’integrazione in software GIS.
\\
\sphinxhline
\sphinxAtStartPar
\sphinxstylestrong{Risorse correlate}
&
\sphinxAtStartPar
Link a Mappe, Documenti o altri Dataset collegati a questa risorsa.
\\
\sphinxbottomrule
\end{tabular}
\sphinxtableafterendhook\par
\sphinxattableend\end{savenotes}


\subsubsection{Visualizzazione Completa dei Metadati}
\label{\detokenize{it/02_trovare_visualizzare:visualizzazione-completa-dei-metadati}}
\sphinxAtStartPar
Per vedere il record completo dei metadati, clicca sulla scheda \sphinxstylestrong{Metadati} (o il
link \sphinxstyleemphasis{Visualizza metadati completi} vicino alla parte inferiore della pagina di
dettaglio). Questa vista mostra tutti i campi dei metadati nella loro forma completa,
compresi i campi conformi agli standard ISO 19115 / INSPIRE.


\subsubsection{Scaricare un Dataset}
\label{\detokenize{it/02_trovare_visualizzare:scaricare-un-dataset}}
\sphinxAtStartPar
Clicca il pulsante \sphinxstylestrong{Download} nella pagina di dettaglio. Appare un selettore di
formato. Le opzioni comuni includono:
\begin{itemize}
\item {} 
\sphinxAtStartPar
\sphinxstylestrong{ESRI Shapefile} \textendash{} per l’uso in QGIS, ArcGIS e la maggior parte dei GIS desktop.

\item {} 
\sphinxAtStartPar
\sphinxstylestrong{GeoJSON} \textendash{} per la mappatura web e l’uso GIS leggero.

\item {} 
\sphinxAtStartPar
\sphinxstylestrong{KML/KMZ} \textendash{} per Google Earth.

\item {} 
\sphinxAtStartPar
\sphinxstylestrong{GeoPackage} \textendash{} formato a file singolo supportato da QGIS.

\item {} 
\sphinxAtStartPar
\sphinxstylestrong{GeoTIFF} \textendash{} per dataset raster.

\item {} 
\sphinxAtStartPar
\sphinxstylestrong{CSV} \textendash{} solo tabella degli attributi (senza geometria).

\end{itemize}

\sphinxAtStartPar
Seleziona il formato preferito e clicca \sphinxstylestrong{Download} per salvare il file.

\sphinxstepscope


\section{Gestire i Metadati}
\label{\detokenize{it/03_gestire_metadati:gestire-i-metadati}}\label{\detokenize{it/03_gestire_metadati:metadata-it}}\label{\detokenize{it/03_gestire_metadati::doc}}
\begin{sphinxcontents}
\sphinxstylecontentstitle{In questa pagina}
\begin{itemize}
\item {} 
\sphinxAtStartPar
\phantomsection\label{\detokenize{it/03_gestire_metadati:id1}}{\hyperref[\detokenize{it/03_gestire_metadati:chi-puo-modificare-i-metadati}]{\sphinxcrossref{Chi Può Modificare i Metadati?}}}

\item {} 
\sphinxAtStartPar
\phantomsection\label{\detokenize{it/03_gestire_metadati:id2}}{\hyperref[\detokenize{it/03_gestire_metadati:aprire-l-editor-dei-metadati}]{\sphinxcrossref{Aprire l’Editor dei Metadati}}}

\item {} 
\sphinxAtStartPar
\phantomsection\label{\detokenize{it/03_gestire_metadati:id3}}{\hyperref[\detokenize{it/03_gestire_metadati:panoramica-dell-editor-dei-metadati}]{\sphinxcrossref{Panoramica dell’Editor dei Metadati}}}

\item {} 
\sphinxAtStartPar
\phantomsection\label{\detokenize{it/03_gestire_metadati:id4}}{\hyperref[\detokenize{it/03_gestire_metadati:modifica-dei-campi-principali-dei-metadati}]{\sphinxcrossref{Modifica dei Campi Principali dei Metadati}}}
\begin{itemize}
\item {} 
\sphinxAtStartPar
\phantomsection\label{\detokenize{it/03_gestire_metadati:id5}}{\hyperref[\detokenize{it/03_gestire_metadati:titolo}]{\sphinxcrossref{Titolo}}}

\item {} 
\sphinxAtStartPar
\phantomsection\label{\detokenize{it/03_gestire_metadati:id6}}{\hyperref[\detokenize{it/03_gestire_metadati:abstract}]{\sphinxcrossref{Abstract}}}

\item {} 
\sphinxAtStartPar
\phantomsection\label{\detokenize{it/03_gestire_metadati:id7}}{\hyperref[\detokenize{it/03_gestire_metadati:parole-chiave}]{\sphinxcrossref{Parole Chiave}}}

\item {} 
\sphinxAtStartPar
\phantomsection\label{\detokenize{it/03_gestire_metadati:id8}}{\hyperref[\detokenize{it/03_gestire_metadati:categoria}]{\sphinxcrossref{Categoria}}}

\item {} 
\sphinxAtStartPar
\phantomsection\label{\detokenize{it/03_gestire_metadati:id9}}{\hyperref[\detokenize{it/03_gestire_metadati:estensione-spaziale}]{\sphinxcrossref{Estensione Spaziale}}}

\item {} 
\sphinxAtStartPar
\phantomsection\label{\detokenize{it/03_gestire_metadati:id10}}{\hyperref[\detokenize{it/03_gestire_metadati:campi-data}]{\sphinxcrossref{Campi Data}}}

\item {} 
\sphinxAtStartPar
\phantomsection\label{\detokenize{it/03_gestire_metadati:id11}}{\hyperref[\detokenize{it/03_gestire_metadati:licenza}]{\sphinxcrossref{Licenza}}}

\item {} 
\sphinxAtStartPar
\phantomsection\label{\detokenize{it/03_gestire_metadati:id12}}{\hyperref[\detokenize{it/03_gestire_metadati:punto-di-contatto}]{\sphinxcrossref{Punto di Contatto}}}

\item {} 
\sphinxAtStartPar
\phantomsection\label{\detokenize{it/03_gestire_metadati:id13}}{\hyperref[\detokenize{it/03_gestire_metadati:miniatura}]{\sphinxcrossref{Miniatura}}}

\end{itemize}

\item {} 
\sphinxAtStartPar
\phantomsection\label{\detokenize{it/03_gestire_metadati:id14}}{\hyperref[\detokenize{it/03_gestire_metadati:salvare-le-modifiche}]{\sphinxcrossref{Salvare le Modifiche}}}

\item {} 
\sphinxAtStartPar
\phantomsection\label{\detokenize{it/03_gestire_metadati:id15}}{\hyperref[\detokenize{it/03_gestire_metadati:lista-di-controllo-per-la-completezza-dei-metadati}]{\sphinxcrossref{Lista di Controllo per la Completezza dei Metadati}}}

\end{itemize}
\end{sphinxcontents}

\sphinxAtStartPar
I metadati di qualità rendono una risorsa trovabile e affidabile. Questa sezione spiega
come visualizzare e modificare i metadati di un Dataset o di un Document in KMaP.

\begin{sphinxadmonition}{note}{Note:}
\sphinxAtStartPar
Puoi modificare i metadati solo delle risorse di cui sei \sphinxstylestrong{proprietario} o per cui
ti sono stati concessi \sphinxstylestrong{permessi di modifica}. I membri del gruppo Accordo Pelagos
possono avere diritti di modifica sulle risorse condivise nel gruppo — contatta il
tuo amministratore di gruppo se hai dubbi.
\end{sphinxadmonition}


\subsection{Chi Può Modificare i Metadati?}
\label{\detokenize{it/03_gestire_metadati:chi-puo-modificare-i-metadati}}
\sphinxAtStartPar
KMaP dispone di un modello di permessi a più livelli:


\begin{savenotes}\sphinxattablestart
\sphinxthistablewithglobalstyle
\centering
\begin{tabular}[t]{\X{25}{100}\X{75}{100}}
\sphinxtoprule
\sphinxstyletheadfamily 
\sphinxAtStartPar
Ruolo
&\sphinxstyletheadfamily 
\sphinxAtStartPar
Diritti di modifica dei metadati
\\
\sphinxmidrule
\sphinxtableatstartofbodyhook
\sphinxAtStartPar
\sphinxstylestrong{Proprietario della risorsa}
&
\sphinxAtStartPar
Diritti di modifica completi su tutti i campi dei metadati.
\\
\sphinxhline
\sphinxAtStartPar
\sphinxstylestrong{Gestore del gruppo}
&
\sphinxAtStartPar
Può modificare le risorse di proprietà di qualsiasi membro del proprio gruppo.
\\
\sphinxhline
\sphinxAtStartPar
\sphinxstylestrong{Membro del gruppo (editor)}
&
\sphinxAtStartPar
Può modificare le risorse condivise esplicitamente con il gruppo con permesso di
\sphinxstyleemphasis{modifica}.
\\
\sphinxhline
\sphinxAtStartPar
\sphinxstylestrong{Visualizzatore / pubblico}
&
\sphinxAtStartPar
Accesso in sola lettura. Non può modificare i metadati.
\\
\sphinxbottomrule
\end{tabular}
\sphinxtableafterendhook\par
\sphinxattableend\end{savenotes}

\sphinxAtStartPar
Se vedi il pulsante \sphinxstylestrong{Edit} nella pagina di dettaglio di una risorsa, hai il diritto
di modificare i metadati di quella risorsa.


\subsection{Aprire l’Editor dei Metadati}
\label{\detokenize{it/03_gestire_metadati:aprire-l-editor-dei-metadati}}
\sphinxAtStartPar
Esistono due modi per aprire l’editor dei metadati:

\sphinxAtStartPar
\sphinxstylestrong{Dalla pagina di dettaglio:}
\begin{enumerate}
\sphinxsetlistlabels{\arabic}{enumi}{enumii}{}{.}%
\item {} 
\sphinxAtStartPar
Naviga alla pagina di dettaglio della risorsa (vedi {\hyperref[\detokenize{it/02_trovare_visualizzare:find-it}]{\sphinxcrossref{\DUrole{std}{\DUrole{std-ref}{Come Trovare e Visualizzare i Tuoi Dati}}}}} per trovare
una risorsa).

\item {} 
\sphinxAtStartPar
Clicca il pulsante \sphinxstylestrong{Edit} (icona matita) nella barra delle azioni in alto a destra.

\item {} 
\sphinxAtStartPar
Appare un menu a discesa. Seleziona \sphinxstylestrong{Edit metadata}.

\end{enumerate}

\sphinxAtStartPar
\sphinxstylestrong{Dal catalogo:}
\begin{enumerate}
\sphinxsetlistlabels{\arabic}{enumi}{enumii}{}{.}%
\item {} 
\sphinxAtStartPar
Nel catalogo, passa il cursore su una scheda risorsa.

\item {} 
\sphinxAtStartPar
Clicca il menu \sphinxstylestrong{⋮} (altre opzioni) che appare.

\item {} 
\sphinxAtStartPar
Seleziona \sphinxstylestrong{Edit metadata}.

\end{enumerate}

\sphinxAtStartPar
L’editor dei metadati si apre come un modulo a più schede.


\subsection{Panoramica dell’Editor dei Metadati}
\label{\detokenize{it/03_gestire_metadati:panoramica-dell-editor-dei-metadati}}
\sphinxAtStartPar
L’editor dei metadati è organizzato in schede. Le più importanti sono:


\begin{savenotes}\sphinxattablestart
\sphinxthistablewithglobalstyle
\centering
\begin{tabular}[t]{\X{25}{100}\X{75}{100}}
\sphinxtoprule
\sphinxstyletheadfamily 
\sphinxAtStartPar
Scheda
&\sphinxstyletheadfamily 
\sphinxAtStartPar
Contenuto
\\
\sphinxmidrule
\sphinxtableatstartofbodyhook
\sphinxAtStartPar
\sphinxstylestrong{Basic info}
&
\sphinxAtStartPar
Titolo, abstract, lingua, categoria, parole chiave, punto di contatto.
\\
\sphinxhline
\sphinxAtStartPar
\sphinxstylestrong{Location and licences}
&
\sphinxAtStartPar
Estensione spaziale, CRS, campi data, licenza, attribuzione.
\\
\sphinxhline
\sphinxAtStartPar
\sphinxstylestrong{Optional metadata}
&
\sphinxAtStartPar
Campi ISO 19115 aggiuntivi (edizione, scopo, informazioni supplementari, …).
\\
\sphinxhline
\sphinxAtStartPar
\sphinxstylestrong{Settings}
&
\sphinxAtStartPar
Permessi, stato in evidenza, stato pubblicato.
\\
\sphinxbottomrule
\end{tabular}
\sphinxtableafterendhook\par
\sphinxattableend\end{savenotes}


\subsection{Modifica dei Campi Principali dei Metadati}
\label{\detokenize{it/03_gestire_metadati:modifica-dei-campi-principali-dei-metadati}}

\subsubsection{Titolo}
\label{\detokenize{it/03_gestire_metadati:titolo}}
\sphinxAtStartPar
Il \sphinxstylestrong{titolo} è il modo principale con cui gli utenti trovano la tua risorsa. Deve essere:
\begin{itemize}
\item {} 
\sphinxAtStartPar
Descrittivo e specifico (evita titoli generici come “Dati” o “Layer cartografico”).

\item {} 
\sphinxAtStartPar
Coerente con le convenzioni di denominazione usate dal gruppo Accordo Pelagos.

\end{itemize}

\sphinxAtStartPar
Per modificarlo: clicca sul campo \sphinxstylestrong{Titolo} nella scheda \sphinxstyleemphasis{Basic info} e digita le
modifiche.


\subsubsection{Abstract}
\label{\detokenize{it/03_gestire_metadati:abstract}}
\sphinxAtStartPar
L’\sphinxstylestrong{abstract} è una descrizione in linguaggio semplice della risorsa. Un buon abstract
risponde a:
\begin{itemize}
\item {} 
\sphinxAtStartPar
Cosa contiene questa risorsa?

\item {} 
\sphinxAtStartPar
Quale area geografica copre?

\item {} 
\sphinxAtStartPar
Quale periodo temporale rappresenta?

\item {} 
\sphinxAtStartPar
A cosa serve?

\item {} 
\sphinxAtStartPar
Quali sono le sue limitazioni note?

\end{itemize}

\sphinxAtStartPar
Il campo abstract supporta la formattazione di base. Mira ad almeno due o tre frasi.


\subsubsection{Parole Chiave}
\label{\detokenize{it/03_gestire_metadati:parole-chiave}}
\sphinxAtStartPar
Le parole chiave sono il meccanismo principale per filtrare e scoprire risorse in KMaP.

\begin{sphinxadmonition}{important}{Important:}
\sphinxAtStartPar
Tutte le risorse relative all’Accordo Pelagos \sphinxstylestrong{devono includere} almeno una delle
seguenti parole chiave:
\begin{itemize}
\item {} 
\sphinxAtStartPar
\sphinxcode{\sphinxupquote{pelagos sanctuary}}

\item {} 
\sphinxAtStartPar
\sphinxcode{\sphinxupquote{pelagos agreement}}

\end{itemize}

\sphinxAtStartPar
Senza queste parole chiave, la risorsa non apparirà quando altri membri del gruppo
filtrano il catalogo con questi termini.
\end{sphinxadmonition}

\sphinxAtStartPar
Per aggiungere una parola chiave:
\begin{enumerate}
\sphinxsetlistlabels{\arabic}{enumi}{enumii}{}{.}%
\item {} 
\sphinxAtStartPar
Nel campo \sphinxstylestrong{Parole chiave}, inizia a digitare la parola chiave.

\item {} 
\sphinxAtStartPar
Se la parola chiave esiste già nel sistema, apparirà in un menu a discesa — selezionala.

\item {} 
\sphinxAtStartPar
Se è nuova, digita la parola chiave completa e premi \sphinxstylestrong{Invio} o clicca \sphinxstylestrong{Aggiungi}.

\item {} 
\sphinxAtStartPar
Per rimuovere una parola chiave, clicca la \sphinxstylestrong{×} accanto ad essa.

\end{enumerate}


\subsubsection{Categoria}
\label{\detokenize{it/03_gestire_metadati:categoria}}
\sphinxAtStartPar
La \sphinxstylestrong{categoria} assegna la risorsa a una classificazione tematica.
Seleziona la categoria più appropriata dal menu a discesa:
\begin{itemize}
\item {} 
\sphinxAtStartPar
Pesca e Acquacoltura

\item {} 
\sphinxAtStartPar
Biodiversità Marina

\item {} 
\sphinxAtStartPar
Inquinamento

\item {} 
\sphinxAtStartPar
Cambiamenti Climatici

\item {} 
\sphinxAtStartPar
Pianificazione Spaziale Marina

\item {} 
\sphinxAtStartPar
Sostenibilità ed Economia Blu

\item {} 
\sphinxAtStartPar
Governance

\end{itemize}

\sphinxAtStartPar
Una risorsa può appartenere a una sola categoria principale. Usa le parole chiave per
la categorizzazione tematica aggiuntiva.


\subsubsection{Estensione Spaziale}
\label{\detokenize{it/03_gestire_metadati:estensione-spaziale}}
\sphinxAtStartPar
L’\sphinxstylestrong{estensione spaziale} definisce il riquadro di delimitazione geografica della risorsa.
Per i dataset, viene solitamente impostata automaticamente al caricamento. Puoi
perfezionarla manualmente:
\begin{enumerate}
\sphinxsetlistlabels{\arabic}{enumi}{enumii}{}{.}%
\item {} 
\sphinxAtStartPar
Vai alla scheda \sphinxstylestrong{Location and licences}.

\item {} 
\sphinxAtStartPar
Nella sezione \sphinxstyleemphasis{Estensione spaziale}, usa il widget mappa per disegnare o regolare
il riquadro, oppure inserisci le coordinate (longitudine e latitudine min/max)
manualmente.

\end{enumerate}

\begin{sphinxadmonition}{tip}{Tip:}
\sphinxAtStartPar
Per le risorse relative a Pelagos, l’estensione spaziale dovrebbe coprire almeno
l’area del Santuario Pelagos (approssimativamente tra Francia, Italia e Monaco nel
Mar Ligure e nel Mediterraneo occidentale adiacente).
\end{sphinxadmonition}


\subsubsection{Campi Data}
\label{\detokenize{it/03_gestire_metadati:campi-data}}
\sphinxAtStartPar
Sono disponibili tre campi data:


\begin{savenotes}\sphinxattablestart
\sphinxthistablewithglobalstyle
\centering
\begin{tabular}[t]{\X{30}{100}\X{70}{100}}
\sphinxtoprule
\sphinxstyletheadfamily 
\sphinxAtStartPar
Campo
&\sphinxstyletheadfamily 
\sphinxAtStartPar
Descrizione
\\
\sphinxmidrule
\sphinxtableatstartofbodyhook
\sphinxAtStartPar
\sphinxstylestrong{Data} (data di creazione)
&
\sphinxAtStartPar
Quando la risorsa è stata originariamente creata o i dati sono stati raccolti.
\\
\sphinxhline
\sphinxAtStartPar
\sphinxstylestrong{Data di pubblicazione}
&
\sphinxAtStartPar
Quando la risorsa è stata pubblicata su KMaP.
\\
\sphinxhline
\sphinxAtStartPar
\sphinxstylestrong{Data di revisione}
&
\sphinxAtStartPar
Quando la risorsa è stata aggiornata l’ultima volta.
\\
\sphinxbottomrule
\end{tabular}
\sphinxtableafterendhook\par
\sphinxattableend\end{savenotes}

\sphinxAtStartPar
Mantieni queste date accurate — vengono usate per filtrare e ordinare le risorse
nel catalogo.


\subsubsection{Licenza}
\label{\detokenize{it/03_gestire_metadati:licenza}}
\sphinxAtStartPar
La \sphinxstylestrong{licenza} definisce come altri possono utilizzare la risorsa.
\begin{enumerate}
\sphinxsetlistlabels{\arabic}{enumi}{enumii}{}{.}%
\item {} 
\sphinxAtStartPar
Vai alla scheda \sphinxstylestrong{Location and licences}.

\item {} 
\sphinxAtStartPar
Seleziona una licenza dal menu a discesa. Le opzioni comuni includono:
\begin{itemize}
\item {} 
\sphinxAtStartPar
\sphinxstyleemphasis{Creative Commons Attribuzione (CC BY)} \textendash{} riuso libero con attribuzione.

\item {} 
\sphinxAtStartPar
\sphinxstyleemphasis{Creative Commons Attribuzione\sphinxhyphen{}CondividiAlloStessoModo (CC BY\sphinxhyphen{}SA)} \textendash{} riuso con
attribuzione e stessa licenza.

\item {} 
\sphinxAtStartPar
\sphinxstyleemphasis{Open Database Licence (ODbL)} \textendash{} per dataset.

\item {} 
\sphinxAtStartPar
\sphinxstyleemphasis{Nessuna restrizione} \textendash{} dominio pubblico.

\end{itemize}

\end{enumerate}

\begin{sphinxadmonition}{note}{Note:}
\sphinxAtStartPar
Se non sei sicuro di quale licenza applicare, consulta il tuo amministratore di gruppo.
Applicare la licenza errata può limitare il riuso dei dati in modo non appropriato.
\end{sphinxadmonition}


\subsubsection{Punto di Contatto}
\label{\detokenize{it/03_gestire_metadati:punto-di-contatto}}
\sphinxAtStartPar
Il \sphinxstylestrong{punto di contatto} identifica chi contattare per domande sulla risorsa.
Può essere una persona o un’organizzazione.
\begin{enumerate}
\sphinxsetlistlabels{\arabic}{enumi}{enumii}{}{.}%
\item {} 
\sphinxAtStartPar
Nella scheda \sphinxstyleemphasis{Basic info}, trova il campo \sphinxstylestrong{Punto di contatto}.

\item {} 
\sphinxAtStartPar
Inizia a digitare un nome per cercare tra gli utenti KMaP registrati, oppure inserisci
il nome di un’organizzazione.

\end{enumerate}


\subsubsection{Miniatura}
\label{\detokenize{it/03_gestire_metadati:miniatura}}
\sphinxAtStartPar
La \sphinxstylestrong{miniatura} è l’immagine di anteprima mostrata nelle schede del catalogo. Viene
generata automaticamente per Dataset e Mappe, ma può essere sostituita manualmente:
\begin{enumerate}
\sphinxsetlistlabels{\arabic}{enumi}{enumii}{}{.}%
\item {} 
\sphinxAtStartPar
Clicca \sphinxstylestrong{Edit} nella pagina di dettaglio.

\item {} 
\sphinxAtStartPar
Seleziona \sphinxstylestrong{Edit thumbnail}.

\item {} 
\sphinxAtStartPar
Carica una nuova immagine (consigliato: 600 × 400 px, JPEG o PNG).

\end{enumerate}


\subsection{Salvare le Modifiche}
\label{\detokenize{it/03_gestire_metadati:salvare-le-modifiche}}
\sphinxAtStartPar
Quando hai finito di modificare, scorri fino in fondo all’editor e clicca \sphinxstylestrong{Salva}
(o \sphinxstylestrong{Aggiorna} — l’etichetta può variare). Apparirà un messaggio di conferma.

\begin{sphinxadmonition}{warning}{Warning:}
\sphinxAtStartPar
Le modifiche ai metadati vengono applicate immediatamente e sono visibili a tutti
gli utenti con accesso alla risorsa. Non esiste una modalità bozza o staging.
\end{sphinxadmonition}

\sphinxAtStartPar
Dopo il salvataggio, vieni reindirizzato alla pagina di dettaglio della risorsa dove
puoi verificare che tutti i campi siano stati aggiornati correttamente.


\subsection{Lista di Controllo per la Completezza dei Metadati}
\label{\detokenize{it/03_gestire_metadati:lista-di-controllo-per-la-completezza-dei-metadati}}
\sphinxAtStartPar
Usa la seguente lista di controllo prima di pubblicare o condividere una risorsa:


\begin{savenotes}\sphinxattablestart
\sphinxthistablewithglobalstyle
\centering
\begin{tabular}[t]{\X{10}{100}\X{90}{100}}
\sphinxtoprule
\sphinxstyletheadfamily 
\sphinxAtStartPar
\(\checkmark\)
&\sphinxstyletheadfamily 
\sphinxAtStartPar
Elemento
\\
\sphinxmidrule
\sphinxtableatstartofbodyhook
\sphinxAtStartPar
☐
&
\sphinxAtStartPar
Il titolo è descrittivo e specifico.
\\
\sphinxhline
\sphinxAtStartPar
☐
&
\sphinxAtStartPar
L’abstract spiega cosa, dove, quando e perché.
\\
\sphinxhline
\sphinxAtStartPar
☐
&
\sphinxAtStartPar
È presente almeno una parola chiave Pelagos (\sphinxcode{\sphinxupquote{pelagos sanctuary}} o
\sphinxcode{\sphinxupquote{pelagos agreement}}).
\\
\sphinxhline
\sphinxAtStartPar
☐
&
\sphinxAtStartPar
La categoria è impostata.
\\
\sphinxhline
\sphinxAtStartPar
☐
&
\sphinxAtStartPar
L’estensione spaziale è definita e copre l’area corretta.
\\
\sphinxhline
\sphinxAtStartPar
☐
&
\sphinxAtStartPar
I campi data sono compilati.
\\
\sphinxhline
\sphinxAtStartPar
☐
&
\sphinxAtStartPar
La licenza è impostata.
\\
\sphinxhline
\sphinxAtStartPar
☐
&
\sphinxAtStartPar
Il punto di contatto è identificato.
\\
\sphinxhline
\sphinxAtStartPar
☐
&
\sphinxAtStartPar
La miniatura è rappresentativa.
\\
\sphinxbottomrule
\end{tabular}
\sphinxtableafterendhook\par
\sphinxattableend\end{savenotes}

\sphinxstepscope


\section{Esplorare le Mappe}
\label{\detokenize{it/04_esplorare_mappe:esplorare-le-mappe}}\label{\detokenize{it/04_esplorare_mappe:maps-it}}\label{\detokenize{it/04_esplorare_mappe::doc}}
\begin{sphinxcontents}
\sphinxstylecontentstitle{In questa pagina}
\begin{itemize}
\item {} 
\sphinxAtStartPar
\phantomsection\label{\detokenize{it/04_esplorare_mappe:id1}}{\hyperref[\detokenize{it/04_esplorare_mappe:trovare-una-mappa}]{\sphinxcrossref{Trovare una Mappa}}}

\item {} 
\sphinxAtStartPar
\phantomsection\label{\detokenize{it/04_esplorare_mappe:id2}}{\hyperref[\detokenize{it/04_esplorare_mappe:l-interfaccia-del-visualizzatore-mappa}]{\sphinxcrossref{L’Interfaccia del Visualizzatore Mappa}}}

\item {} 
\sphinxAtStartPar
\phantomsection\label{\detokenize{it/04_esplorare_mappe:id3}}{\hyperref[\detokenize{it/04_esplorare_mappe:navigare-la-mappa}]{\sphinxcrossref{Navigare la Mappa}}}
\begin{itemize}
\item {} 
\sphinxAtStartPar
\phantomsection\label{\detokenize{it/04_esplorare_mappe:id4}}{\hyperref[\detokenize{it/04_esplorare_mappe:navigazione-di-base}]{\sphinxcrossref{Navigazione di Base}}}

\item {} 
\sphinxAtStartPar
\phantomsection\label{\detokenize{it/04_esplorare_mappe:id5}}{\hyperref[\detokenize{it/04_esplorare_mappe:andare-a-una-posizione-specifica}]{\sphinxcrossref{Andare a una Posizione Specifica}}}

\end{itemize}

\item {} 
\sphinxAtStartPar
\phantomsection\label{\detokenize{it/04_esplorare_mappe:id6}}{\hyperref[\detokenize{it/04_esplorare_mappe:lavorare-con-i-layer}]{\sphinxcrossref{Lavorare con i Layer}}}
\begin{itemize}
\item {} 
\sphinxAtStartPar
\phantomsection\label{\detokenize{it/04_esplorare_mappe:id7}}{\hyperref[\detokenize{it/04_esplorare_mappe:aprire-il-pannello-layer}]{\sphinxcrossref{Aprire il Pannello Layer}}}

\item {} 
\sphinxAtStartPar
\phantomsection\label{\detokenize{it/04_esplorare_mappe:id8}}{\hyperref[\detokenize{it/04_esplorare_mappe:attivare-disattivare-la-visibilita-dei-layer}]{\sphinxcrossref{Attivare/Disattivare la Visibilità dei Layer}}}

\item {} 
\sphinxAtStartPar
\phantomsection\label{\detokenize{it/04_esplorare_mappe:id9}}{\hyperref[\detokenize{it/04_esplorare_mappe:cambiare-l-ordine-dei-layer}]{\sphinxcrossref{Cambiare l’Ordine dei Layer}}}

\item {} 
\sphinxAtStartPar
\phantomsection\label{\detokenize{it/04_esplorare_mappe:id10}}{\hyperref[\detokenize{it/04_esplorare_mappe:regolare-l-opacita-del-layer}]{\sphinxcrossref{Regolare l’Opacità del Layer}}}

\item {} 
\sphinxAtStartPar
\phantomsection\label{\detokenize{it/04_esplorare_mappe:id11}}{\hyperref[\detokenize{it/04_esplorare_mappe:visualizzare-la-legenda}]{\sphinxcrossref{Visualizzare la Legenda}}}

\end{itemize}

\item {} 
\sphinxAtStartPar
\phantomsection\label{\detokenize{it/04_esplorare_mappe:id12}}{\hyperref[\detokenize{it/04_esplorare_mappe:interrogare-gli-elementi}]{\sphinxcrossref{Interrogare gli Elementi}}}
\begin{itemize}
\item {} 
\sphinxAtStartPar
\phantomsection\label{\detokenize{it/04_esplorare_mappe:id13}}{\hyperref[\detokenize{it/04_esplorare_mappe:cliccare-su-un-elemento}]{\sphinxcrossref{Cliccare su un Elemento}}}

\item {} 
\sphinxAtStartPar
\phantomsection\label{\detokenize{it/04_esplorare_mappe:id14}}{\hyperref[\detokenize{it/04_esplorare_mappe:strumento-identifica}]{\sphinxcrossref{Strumento Identifica}}}

\item {} 
\sphinxAtStartPar
\phantomsection\label{\detokenize{it/04_esplorare_mappe:id15}}{\hyperref[\detokenize{it/04_esplorare_mappe:selezionare-e-filtrare-gli-elementi}]{\sphinxcrossref{Selezionare e Filtrare gli Elementi}}}

\end{itemize}

\item {} 
\sphinxAtStartPar
\phantomsection\label{\detokenize{it/04_esplorare_mappe:id16}}{\hyperref[\detokenize{it/04_esplorare_mappe:strumenti-di-misurazione}]{\sphinxcrossref{Strumenti di Misurazione}}}
\begin{itemize}
\item {} 
\sphinxAtStartPar
\phantomsection\label{\detokenize{it/04_esplorare_mappe:id17}}{\hyperref[\detokenize{it/04_esplorare_mappe:misurazione-della-distanza}]{\sphinxcrossref{Misurazione della Distanza}}}

\item {} 
\sphinxAtStartPar
\phantomsection\label{\detokenize{it/04_esplorare_mappe:id18}}{\hyperref[\detokenize{it/04_esplorare_mappe:misurazione-dell-area}]{\sphinxcrossref{Misurazione dell’Area}}}

\end{itemize}

\item {} 
\sphinxAtStartPar
\phantomsection\label{\detokenize{it/04_esplorare_mappe:id19}}{\hyperref[\detokenize{it/04_esplorare_mappe:stampare-ed-esportare-la-vista-mappa}]{\sphinxcrossref{Stampare ed Esportare la Vista Mappa}}}
\begin{itemize}
\item {} 
\sphinxAtStartPar
\phantomsection\label{\detokenize{it/04_esplorare_mappe:id20}}{\hyperref[\detokenize{it/04_esplorare_mappe:stampa}]{\sphinxcrossref{Stampa}}}

\item {} 
\sphinxAtStartPar
\phantomsection\label{\detokenize{it/04_esplorare_mappe:id21}}{\hyperref[\detokenize{it/04_esplorare_mappe:condividi-permalink}]{\sphinxcrossref{Condividi / Permalink}}}

\item {} 
\sphinxAtStartPar
\phantomsection\label{\detokenize{it/04_esplorare_mappe:id22}}{\hyperref[\detokenize{it/04_esplorare_mappe:incorpora}]{\sphinxcrossref{Incorpora}}}

\end{itemize}

\item {} 
\sphinxAtStartPar
\phantomsection\label{\detokenize{it/04_esplorare_mappe:id23}}{\hyperref[\detokenize{it/04_esplorare_mappe:salvare-le-modifiche-a-una-mappa}]{\sphinxcrossref{Salvare le Modifiche a una Mappa}}}

\item {} 
\sphinxAtStartPar
\phantomsection\label{\detokenize{it/04_esplorare_mappe:id24}}{\hyperref[\detokenize{it/04_esplorare_mappe:creare-una-nuova-mappa}]{\sphinxcrossref{Creare una Nuova Mappa}}}

\end{itemize}
\end{sphinxcontents}

\sphinxAtStartPar
Una \sphinxstylestrong{Mappa} in KMaP è una mappa web interattiva a più strati, costruita da uno o più
Dataset. Questa sezione spiega come trovare, aprire e navigare le Mappe, come lavorare
con i layer e come estrarre informazioni dai dati visualizzati.


\subsection{Trovare una Mappa}
\label{\detokenize{it/04_esplorare_mappe:trovare-una-mappa}}
\sphinxAtStartPar
Le Mappe sono elencate nel catalogo insieme a tutti gli altri tipi di risorse.
Per filtrare il catalogo e mostrare solo le Mappe:
\begin{enumerate}
\sphinxsetlistlabels{\arabic}{enumi}{enumii}{}{.}%
\item {} 
\sphinxAtStartPar
Apri il catalogo: clicca \sphinxstylestrong{Maps} nella barra di navigazione in alto.

\item {} 
\sphinxAtStartPar
Nel pannello dei filtri a sinistra, sotto \sphinxstylestrong{Tipo di risorsa}, seleziona \sphinxstylestrong{Map}.

\item {} 
\sphinxAtStartPar
Per trovare mappe relative a Pelagos, applica anche il filtro per parola chiave:
nel campo \sphinxstylestrong{Parole chiave}, digita \sphinxcode{\sphinxupquote{pelagos sanctuary}} o \sphinxcode{\sphinxupquote{pelagos agreement}}.

\end{enumerate}

\sphinxAtStartPar
In alternativa, cerca una mappa per titolo usando la barra di ricerca.

\sphinxAtStartPar
Clicca su una scheda Mappa per aprire la relativa \sphinxstylestrong{pagina di dettaglio}. Da lì,
clicca \sphinxstylestrong{View map} per aprirla nel visualizzatore mappa interattivo.


\subsection{L’Interfaccia del Visualizzatore Mappa}
\label{\detokenize{it/04_esplorare_mappe:l-interfaccia-del-visualizzatore-mappa}}
\sphinxAtStartPar
Il visualizzatore mappa (basato su \sphinxstylestrong{MapStore}) si apre in un layout a schermo intero.
I suoi elementi principali sono:


\begin{savenotes}\sphinxattablestart
\sphinxthistablewithglobalstyle
\centering
\begin{tabular}[t]{\X{30}{100}\X{70}{100}}
\sphinxtoprule
\sphinxstyletheadfamily 
\sphinxAtStartPar
Elemento
&\sphinxstyletheadfamily 
\sphinxAtStartPar
Descrizione
\\
\sphinxmidrule
\sphinxtableatstartofbodyhook
\sphinxAtStartPar
\sphinxstylestrong{Area mappa principale}
&
\sphinxAtStartPar
L’area mappa interattiva centrale.
\\
\sphinxhline
\sphinxAtStartPar
\sphinxstylestrong{Pannello layer} (☰)
&
\sphinxAtStartPar
Attiva/disattiva la visibilità dei layer, riordina i layer, accede alle impostazioni.
\\
\sphinxhline
\sphinxAtStartPar
\sphinxstylestrong{Selettore basemap}
&
\sphinxAtStartPar
Cambia la mappa di sfondo (OpenStreetMap, immagini satellitari, …).
\\
\sphinxhline
\sphinxAtStartPar
\sphinxstylestrong{Controlli zoom} (+ / −)
&
\sphinxAtStartPar
Zoom in e out. Supporta anche la rotella del mouse e il gesto di pizzico.
\\
\sphinxhline
\sphinxAtStartPar
\sphinxstylestrong{Barra degli strumenti}
&
\sphinxAtStartPar
Riga di icone degli strumenti in cima o sul lato del canvas.
\\
\sphinxhline
\sphinxAtStartPar
\sphinxstylestrong{Scala}
&
\sphinxAtStartPar
Mostra la scala corrente della mappa.
\\
\sphinxhline
\sphinxAtStartPar
\sphinxstylestrong{Pannello coordinate}
&
\sphinxAtStartPar
Mostra le coordinate geografiche della posizione del cursore.
\\
\sphinxbottomrule
\end{tabular}
\sphinxtableafterendhook\par
\sphinxattableend\end{savenotes}


\subsection{Navigare la Mappa}
\label{\detokenize{it/04_esplorare_mappe:navigare-la-mappa}}

\subsubsection{Navigazione di Base}
\label{\detokenize{it/04_esplorare_mappe:navigazione-di-base}}

\begin{savenotes}\sphinxattablestart
\sphinxthistablewithglobalstyle
\centering
\begin{tabular}[t]{\X{35}{100}\X{65}{100}}
\sphinxtoprule
\sphinxstyletheadfamily 
\sphinxAtStartPar
Azione
&\sphinxstyletheadfamily 
\sphinxAtStartPar
Come eseguirla
\\
\sphinxmidrule
\sphinxtableatstartofbodyhook
\sphinxAtStartPar
\sphinxstylestrong{Spostare la mappa (pan)}
&
\sphinxAtStartPar
Clicca e trascina sul canvas della mappa.
\\
\sphinxhline
\sphinxAtStartPar
\sphinxstylestrong{Zoom in}
&
\sphinxAtStartPar
Scorri la rotella del mouse in avanti, clicca \sphinxstylestrong{+} o fai doppio clic sulla mappa.
\\
\sphinxhline
\sphinxAtStartPar
\sphinxstylestrong{Zoom out}
&
\sphinxAtStartPar
Scorri la rotella del mouse all’indietro o clicca \sphinxstylestrong{−}.
\\
\sphinxhline
\sphinxAtStartPar
\sphinxstylestrong{Torna all’estensione completa}
&
\sphinxAtStartPar
Clicca l’icona \sphinxstyleemphasis{Home} (🏠) nella barra degli strumenti per tornare alla vista
predefinita.
\\
\sphinxhline
\sphinxAtStartPar
\sphinxstylestrong{Zoom su un layer}
&
\sphinxAtStartPar
Nel pannello layer, fai clic destro su un layer e seleziona \sphinxstyleemphasis{Zoom all’estensione
del layer}.
\\
\sphinxbottomrule
\end{tabular}
\sphinxtableafterendhook\par
\sphinxattableend\end{savenotes}


\subsubsection{Andare a una Posizione Specifica}
\label{\detokenize{it/04_esplorare_mappe:andare-a-una-posizione-specifica}}
\sphinxAtStartPar
Usa lo strumento \sphinxstylestrong{Ricerca / Geocodificatore} nella barra degli strumenti (icona 🔍)
per cercare un nome di luogo. Digita il nome e premi Invio — la mappa si sposterà
su quella posizione.


\subsection{Lavorare con i Layer}
\label{\detokenize{it/04_esplorare_mappe:lavorare-con-i-layer}}

\subsubsection{Aprire il Pannello Layer}
\label{\detokenize{it/04_esplorare_mappe:aprire-il-pannello-layer}}
\sphinxAtStartPar
Clicca l’icona \sphinxstylestrong{☰} (hamburger) per aprire il pannello dei layer. Elenca tutti i
layer inclusi nella mappa, raggruppati in:
\begin{itemize}
\item {} 
\sphinxAtStartPar
\sphinxstylestrong{Overlay} \textendash{} gli strati di dati tematici (es. confine del Santuario Pelagos, punti
di avvistamento di cetacei, rotte navali).

\item {} 
\sphinxAtStartPar
\sphinxstylestrong{Layer di sfondo} \textendash{} le opzioni basemap (solo una può essere attiva alla volta).

\end{itemize}


\subsubsection{Attivare/Disattivare la Visibilità dei Layer}
\label{\detokenize{it/04_esplorare_mappe:attivare-disattivare-la-visibilita-dei-layer}}
\sphinxAtStartPar
Ogni layer nel pannello ha un’\sphinxstylestrong{icona occhio} (👁). Cliccala per mostrare o nascondere
quel layer. I layer nascosti rimangono nella composizione della mappa ma non vengono
visualizzati sul canvas.


\subsubsection{Cambiare l’Ordine dei Layer}
\label{\detokenize{it/04_esplorare_mappe:cambiare-l-ordine-dei-layer}}
\sphinxAtStartPar
I layer sono disegnati nell’ordine in cui appaiono nel pannello (i layer in cima sono
disegnati sopra). Per riordinare:
\begin{enumerate}
\sphinxsetlistlabels{\arabic}{enumi}{enumii}{}{.}%
\item {} 
\sphinxAtStartPar
Clicca e tieni premuto un nome di layer nel pannello.

\item {} 
\sphinxAtStartPar
Trascinalo su o giù nella posizione desiderata.

\item {} 
\sphinxAtStartPar
Rilascia per posizionarlo.

\end{enumerate}


\subsubsection{Regolare l’Opacità del Layer}
\label{\detokenize{it/04_esplorare_mappe:regolare-l-opacita-del-layer}}\begin{enumerate}
\sphinxsetlistlabels{\arabic}{enumi}{enumii}{}{.}%
\item {} 
\sphinxAtStartPar
Nel pannello layer, clicca l’icona \sphinxstylestrong{⋮} (opzioni) accanto a un layer.

\item {} 
\sphinxAtStartPar
Seleziona \sphinxstylestrong{Impostazioni} (o l’icona del cursore).

\item {} 
\sphinxAtStartPar
Usa il cursore \sphinxstylestrong{Opacità} per impostare la trasparenza tra 0\% (invisibile) e 100\%
(solido).

\end{enumerate}


\subsubsection{Visualizzare la Legenda}
\label{\detokenize{it/04_esplorare_mappe:visualizzare-la-legenda}}\begin{enumerate}
\sphinxsetlistlabels{\arabic}{enumi}{enumii}{}{.}%
\item {} 
\sphinxAtStartPar
Clicca l’icona \sphinxstylestrong{⋮} (opzioni) accanto a un layer.

\item {} 
\sphinxAtStartPar
Seleziona \sphinxstylestrong{Legenda} per espandere la chiave dei simboli per quel layer.

\end{enumerate}

\sphinxAtStartPar
In alternativa, alcune configurazioni di mappa mostrano un pannello legenda persistente
sul lato del canvas.


\subsection{Interrogare gli Elementi}
\label{\detokenize{it/04_esplorare_mappe:interrogare-gli-elementi}}

\subsubsection{Cliccare su un Elemento}
\label{\detokenize{it/04_esplorare_mappe:cliccare-su-un-elemento}}
\sphinxAtStartPar
Clicca su qualsiasi elemento visibile (punto, linea o poligono) sul canvas della mappa
per aprire un \sphinxstylestrong{pop\sphinxhyphen{}up con gli attributi}. Il pop\sphinxhyphen{}up mostra i valori degli attributi
associati a quell’elemento — ad esempio, il nome di un’area marina protetta, il nome
di una specie o un valore di monitoraggio.

\sphinxAtStartPar
Se più layer si sovrappongono nel punto cliccato, il pop\sphinxhyphen{}up può mostrare risultati
da tutti i layer sovrapposti. Usa i pulsanti freccia nel pop\sphinxhyphen{}up per navigare tra di essi.


\subsubsection{Strumento Identifica}
\label{\detokenize{it/04_esplorare_mappe:strumento-identifica}}
\sphinxAtStartPar
Lo strumento \sphinxstylestrong{Identifica} (icona ℹ️ nella barra degli strumenti) funziona in modo
simile al clic diretto sulla mappa. Interroga tutti i layer visibili nel punto cliccato
e restituisce i loro attributi in un pannello.


\subsubsection{Selezionare e Filtrare gli Elementi}
\label{\detokenize{it/04_esplorare_mappe:selezionare-e-filtrare-gli-elementi}}
\sphinxAtStartPar
Per i layer vettoriali, MapStore offre uno strumento \sphinxstylestrong{Filtro}:
\begin{enumerate}
\sphinxsetlistlabels{\arabic}{enumi}{enumii}{}{.}%
\item {} 
\sphinxAtStartPar
Nel pannello layer, clicca \sphinxstylestrong{⋮} accanto al layer desiderato.

\item {} 
\sphinxAtStartPar
Seleziona \sphinxstylestrong{Filtro}.

\item {} 
\sphinxAtStartPar
Definisci una condizione di filtro usando i campi attributo
(es. mostra solo i record dove \sphinxstyleemphasis{paese = “Italia”}).

\item {} 
\sphinxAtStartPar
Clicca \sphinxstylestrong{Applica}. Vengono visualizzati solo gli elementi che corrispondono al filtro.

\end{enumerate}

\begin{sphinxadmonition}{tip}{Tip:}
\sphinxAtStartPar
Usa lo strumento filtro per concentrarti su un sottoinsieme specifico di dati Pelagos —
ad esempio, mostrando solo gli avvistamenti di cetacei di un anno specifico, o solo le
rotte che attraversano il santuario.
\end{sphinxadmonition}


\subsection{Strumenti di Misurazione}
\label{\detokenize{it/04_esplorare_mappe:strumenti-di-misurazione}}

\subsubsection{Misurazione della Distanza}
\label{\detokenize{it/04_esplorare_mappe:misurazione-della-distanza}}\begin{enumerate}
\sphinxsetlistlabels{\arabic}{enumi}{enumii}{}{.}%
\item {} 
\sphinxAtStartPar
Clicca lo strumento \sphinxstylestrong{Misura} (icona 📐) nella barra degli strumenti.

\item {} 
\sphinxAtStartPar
Seleziona \sphinxstylestrong{Distanza}.

\item {} 
\sphinxAtStartPar
Clicca punti sulla mappa per definire una linea. Doppio clic per terminare.

\item {} 
\sphinxAtStartPar
La distanza totale viene mostrata nell’unità scelta (km, miglia, mn).

\end{enumerate}


\subsubsection{Misurazione dell’Area}
\label{\detokenize{it/04_esplorare_mappe:misurazione-dell-area}}\begin{enumerate}
\sphinxsetlistlabels{\arabic}{enumi}{enumii}{}{.}%
\item {} 
\sphinxAtStartPar
Clicca lo strumento \sphinxstylestrong{Misura}.

\item {} 
\sphinxAtStartPar
Seleziona \sphinxstylestrong{Area}.

\item {} 
\sphinxAtStartPar
Clicca per definire i vertici di un poligono. Doppio clic per chiudere.

\item {} 
\sphinxAtStartPar
L’area calcolata viene visualizzata.

\end{enumerate}

\begin{sphinxadmonition}{note}{Note:}
\sphinxAtStartPar
Le misurazioni vengono eseguite sulla proiezione della mappa e possono presentare
piccole imprecisioni geometriche per aree grandi. Per misurazioni precise, usa
un’applicazione GIS desktop.
\end{sphinxadmonition}


\subsection{Stampare ed Esportare la Vista Mappa}
\label{\detokenize{it/04_esplorare_mappe:stampare-ed-esportare-la-vista-mappa}}

\subsubsection{Stampa}
\label{\detokenize{it/04_esplorare_mappe:stampa}}\begin{enumerate}
\sphinxsetlistlabels{\arabic}{enumi}{enumii}{}{.}%
\item {} 
\sphinxAtStartPar
Clicca lo strumento \sphinxstylestrong{Stampa} (icona 🖨) nella barra degli strumenti.

\item {} 
\sphinxAtStartPar
Appare una finestra di dialogo per la stampa. Configura:
\begin{itemize}
\item {} 
\sphinxAtStartPar
\sphinxstylestrong{Titolo} \textendash{} un titolo da visualizzare sulla mappa stampata.

\item {} 
\sphinxAtStartPar
\sphinxstylestrong{Formato carta} \textendash{} A4, A3, Letter, ecc.

\item {} 
\sphinxAtStartPar
\sphinxstylestrong{Risoluzione} \textendash{} 72, 150 o 300 dpi.

\item {} 
\sphinxAtStartPar
\sphinxstylestrong{Layout} \textendash{} verticale o orizzontale.

\end{itemize}

\item {} 
\sphinxAtStartPar
Clicca \sphinxstylestrong{Stampa} per generare un PDF o un’immagine.

\end{enumerate}


\subsubsection{Condividi / Permalink}
\label{\detokenize{it/04_esplorare_mappe:condividi-permalink}}
\sphinxAtStartPar
Per condividere la vista corrente della mappa (inclusi livello di zoom, layer attivi
e punto centrale):
\begin{enumerate}
\sphinxsetlistlabels{\arabic}{enumi}{enumii}{}{.}%
\item {} 
\sphinxAtStartPar
Clicca il pulsante \sphinxstylestrong{Condividi} (icona link) nella barra degli strumenti.

\item {} 
\sphinxAtStartPar
Viene generato un URL permalink. Copia e condividi questo URL — chi lo riceve aprirà
la mappa nello stesso stato.

\end{enumerate}


\subsubsection{Incorpora}
\label{\detokenize{it/04_esplorare_mappe:incorpora}}\begin{enumerate}
\sphinxsetlistlabels{\arabic}{enumi}{enumii}{}{.}%
\item {} 
\sphinxAtStartPar
Clicca il pulsante \sphinxstylestrong{Condividi}.

\item {} 
\sphinxAtStartPar
Passa alla scheda \sphinxstylestrong{Incorpora}.

\item {} 
\sphinxAtStartPar
Copia il codice HTML \sphinxcode{\sphinxupquote{\textless{}iframe\textgreater{}}} e incollalo nel tuo sito web o rapporto.

\end{enumerate}


\subsection{Salvare le Modifiche a una Mappa}
\label{\detokenize{it/04_esplorare_mappe:salvare-le-modifiche-a-una-mappa}}
\sphinxAtStartPar
Se hai il \sphinxstylestrong{permesso di modifica} su una Mappa, puoi salvare le modifiche a layer,
stili e vista iniziale:
\begin{enumerate}
\sphinxsetlistlabels{\arabic}{enumi}{enumii}{}{.}%
\item {} 
\sphinxAtStartPar
Effettua le modifiche nel visualizzatore mappa (riordina i layer, cambia gli stili,
imposta una nuova estensione iniziale, ecc.).

\item {} 
\sphinxAtStartPar
Clicca il pulsante \sphinxstylestrong{Salva} (icona 💾) nella barra degli strumenti.

\item {} 
\sphinxAtStartPar
Scegli \sphinxstylestrong{Salva} (sovrascrivi la mappa esistente) o \sphinxstylestrong{Salva come} (crea una nuova mappa).

\end{enumerate}

\begin{sphinxadmonition}{warning}{Warning:}
\sphinxAtStartPar
Il salvataggio sovrascrive la mappa per tutti gli utenti che vi hanno accesso.
Se vuoi sperimentare senza influenzare gli altri, usa prima \sphinxstylestrong{Salva come} per
creare una copia personale.
\end{sphinxadmonition}


\subsection{Creare una Nuova Mappa}
\label{\detokenize{it/04_esplorare_mappe:creare-una-nuova-mappa}}
\sphinxAtStartPar
I membri del gruppo Accordo Pelagos con le autorizzazioni appropriate possono creare
nuove Mappe:
\begin{enumerate}
\sphinxsetlistlabels{\arabic}{enumi}{enumii}{}{.}%
\item {} 
\sphinxAtStartPar
Vai su \sphinxstylestrong{Maps} nella barra di navigazione.

\item {} 
\sphinxAtStartPar
Clicca \sphinxstylestrong{+ Crea mappa} (o il pulsante equivalente in cima al catalogo).

\item {} 
\sphinxAtStartPar
Il visualizzatore mappa si apre con un canvas vuoto.

\item {} 
\sphinxAtStartPar
Usa il pulsante \sphinxstylestrong{Aggiungi layer} (icona + nel pannello layer) per cercare e aggiungere
Dataset dal catalogo KMaP.

\item {} 
\sphinxAtStartPar
Personalizza lo stile di ogni layer secondo necessità.

\item {} 
\sphinxAtStartPar
Clicca \sphinxstylestrong{Salva} e fornisci un titolo, un abstract e delle parole chiave
(ricorda di includere \sphinxcode{\sphinxupquote{pelagos sanctuary}} o \sphinxcode{\sphinxupquote{pelagos agreement}}).

\end{enumerate}

\sphinxstepscope


\chapter{KMaP \textendash{} Knowledge Management Platform}
\label{\detokenize{fr/index:kmap-knowledge-management-platform}}\label{\detokenize{fr/index::doc}}
\sphinxAtStartPar
\sphinxstylestrong{Documentation utilisateur pour le groupe Accord Pelagos}

\begin{sphinxadmonition}{note}{Note:}
\sphinxAtStartPar
Cette documentation est destinée exclusivement aux membres du groupe
\sphinxstylestrong{Accord Pelagos} sur la plateforme KMaP. Certaines fonctionnalités décrites ici
peuvent ne pas être visibles aux utilisateurs extérieurs à ce groupe.
\end{sphinxadmonition}

\sphinxAtStartPar
La plateforme KMaP est accessible à l’adresse : \sphinxurl{https://kmap.info-rac.org}

\sphinxstepscope


\section{Bienvenue sur KMaP}
\label{\detokenize{fr/01_bienvenue:bienvenue-sur-kmap}}\label{\detokenize{fr/01_bienvenue:welcome-fr}}\label{\detokenize{fr/01_bienvenue::doc}}
\begin{sphinxcontents}
\sphinxstylecontentstitle{Sur cette page}
\begin{itemize}
\item {} 
\sphinxAtStartPar
\phantomsection\label{\detokenize{fr/01_bienvenue:id1}}{\hyperref[\detokenize{fr/01_bienvenue:qu-est-ce-que-kmap}]{\sphinxcrossref{Qu’est\sphinxhyphen{}ce que KMaP ?}}}

\item {} 
\sphinxAtStartPar
\phantomsection\label{\detokenize{fr/01_bienvenue:id2}}{\hyperref[\detokenize{fr/01_bienvenue:types-de-ressources}]{\sphinxcrossref{Types de Ressources}}}
\begin{itemize}
\item {} 
\sphinxAtStartPar
\phantomsection\label{\detokenize{fr/01_bienvenue:id3}}{\hyperref[\detokenize{fr/01_bienvenue:jeu-de-donnees-dataset}]{\sphinxcrossref{Jeu de données (Dataset)}}}

\item {} 
\sphinxAtStartPar
\phantomsection\label{\detokenize{fr/01_bienvenue:id4}}{\hyperref[\detokenize{fr/01_bienvenue:document}]{\sphinxcrossref{Document}}}

\item {} 
\sphinxAtStartPar
\phantomsection\label{\detokenize{fr/01_bienvenue:id5}}{\hyperref[\detokenize{fr/01_bienvenue:carte-map}]{\sphinxcrossref{Carte (Map)}}}

\item {} 
\sphinxAtStartPar
\phantomsection\label{\detokenize{fr/01_bienvenue:id6}}{\hyperref[\detokenize{fr/01_bienvenue:geostory}]{\sphinxcrossref{GeoStory}}}

\item {} 
\sphinxAtStartPar
\phantomsection\label{\detokenize{fr/01_bienvenue:id7}}{\hyperref[\detokenize{fr/01_bienvenue:tableau-de-bord-dashboard}]{\sphinxcrossref{Tableau de bord (Dashboard)}}}

\end{itemize}

\item {} 
\sphinxAtStartPar
\phantomsection\label{\detokenize{fr/01_bienvenue:id8}}{\hyperref[\detokenize{fr/01_bienvenue:services-interoperables}]{\sphinxcrossref{Services Interopérables}}}

\item {} 
\sphinxAtStartPar
\phantomsection\label{\detokenize{fr/01_bienvenue:id9}}{\hyperref[\detokenize{fr/01_bienvenue:thematiques-disponibles-sur-kmap}]{\sphinxcrossref{Thématiques Disponibles sur KMaP}}}

\end{itemize}
\end{sphinxcontents}


\subsection{Qu’est\sphinxhyphen{}ce que KMaP ?}
\label{\detokenize{fr/01_bienvenue:qu-est-ce-que-kmap}}
\sphinxAtStartPar
\sphinxstylestrong{KMaP} (Knowledge Management Platform) est l’Infrastructure de Données Spatiales
développée par \sphinxstylestrong{INFO/RAC}, le Centre d’Activités Régionales pour l’Information et la
Communication du \sphinxhref{https://www.unep.org/unepmap/}{Plan d’Action pour la Méditerranée}
(UNEP/MAP).

\sphinxAtStartPar
KMaP est basée sur \sphinxhref{https://geonode.org/}{GeoNode}, une plateforme web open\sphinxhyphen{}source
pour la gestion et le partage d’informations géospatiales. Elle constitue un point
d’accès unifié au patrimoine de connaissances environnementales et spatiales de la
Méditerranée.

\sphinxAtStartPar
La plateforme est structurée autour de trois hubs :


\begin{savenotes}\sphinxattablestart
\sphinxthistablewithglobalstyle
\centering
\begin{tabular}[t]{\X{20}{100}\X{15}{100}\X{65}{100}}
\sphinxtoprule
\sphinxstyletheadfamily 
\sphinxAtStartPar
Hub
&\sphinxstyletheadfamily 
\sphinxAtStartPar
Accès
&\sphinxstyletheadfamily 
\sphinxAtStartPar
Description
\\
\sphinxmidrule
\sphinxtableatstartofbodyhook
\sphinxAtStartPar
\sphinxstylestrong{Data Hub}
&
\sphinxAtStartPar
Bouton \sphinxstyleemphasis{Maps}
&
\sphinxAtStartPar
Données géospatiales : jeux de données, cartes interactives, tableaux de bord et
géohistoires.
\\
\sphinxhline
\sphinxAtStartPar
\sphinxstylestrong{Knowledge Hub}
&
\sphinxAtStartPar
Bouton \sphinxstyleemphasis{Library}
&
\sphinxAtStartPar
Patrimoine documentaire : rapports, publications et documents de référence.
\\
\sphinxhline
\sphinxAtStartPar
\sphinxstylestrong{Exchange Hub}
&
\sphinxAtStartPar
Bouton \sphinxstyleemphasis{Network}
&
\sphinxAtStartPar
Collaboration et échange de connaissances (en cours de préparation).
\\
\sphinxbottomrule
\end{tabular}
\sphinxtableafterendhook\par
\sphinxattableend\end{savenotes}

\begin{sphinxadmonition}{note}{Note:}
\sphinxAtStartPar
En tant que membre du groupe \sphinxstylestrong{Accord Pelagos}, vous avez accès à des ressources
spécifiques, des droits de modification des métadonnées et des autorisations au
niveau du groupe qui ne sont pas disponibles pour les utilisateurs publics.
\end{sphinxadmonition}


\subsection{Types de Ressources}
\label{\detokenize{fr/01_bienvenue:types-de-ressources}}
\sphinxAtStartPar
KMaP organise tout le contenu en cinq types de ressources, chacun ayant un objectif
distinct.


\subsubsection{Jeu de données (Dataset)}
\label{\detokenize{fr/01_bienvenue:jeu-de-donnees-dataset}}
\sphinxAtStartPar
Un \sphinxstylestrong{Dataset} est une couche de données géographiques représentant des éléments ou
des mesures du monde réel. Il constitue la brique fondamentale de la plateforme —
la matière première à partir de laquelle sont construites les cartes, les tableaux de
bord et les analyses.

\sphinxAtStartPar
Les datasets peuvent être :
\begin{itemize}
\item {} 
\sphinxAtStartPar
\sphinxstylestrong{Vecteur} \textendash{} éléments géographiques discrets représentés sous forme de points,
lignes ou polygones (ex. limites de zones protégées, réseaux hydrographiques,
emplacements de stations de surveillance).

\item {} 
\sphinxAtStartPar
\sphinxstylestrong{Raster} \textendash{} surfaces continues représentées sous forme de grille de cellules
(ex. température de surface de la mer, imagerie satellitaire, modèles bathymétriques).

\end{itemize}

\sphinxAtStartPar
Ce que vous pouvez faire avec un Dataset :
\begin{itemize}
\item {} 
\sphinxAtStartPar
Le visualiser de manière interactive sur une carte web

\item {} 
\sphinxAtStartPar
Consulter sa table attributaire

\item {} 
\sphinxAtStartPar
Le télécharger dans des formats SIG standard (Shapefile, GeoTIFF, GeoJSON, KML, CSV, …)

\item {} 
\sphinxAtStartPar
Lire et modifier ses métadonnées

\item {} 
\sphinxAtStartPar
L’ajouter comme couche à une Carte

\end{itemize}

\begin{sphinxadmonition}{tip}{Tip:}
\sphinxAtStartPar
Les datasets publiés pour l’Accord Pelagos sont étiquetés avec les mots\sphinxhyphen{}clés
\sphinxcode{\sphinxupquote{pelagos sanctuary}} ou \sphinxcode{\sphinxupquote{pelagos agreement}}. Utilisez ces mots\sphinxhyphen{}clés pour filtrer
rapidement le catalogue (voir {\hyperref[\detokenize{fr/02_trouver_visualiser:find-fr}]{\sphinxcrossref{\DUrole{std}{\DUrole{std-ref}{Comment Trouver et Visualiser Vos Données}}}}}).
\end{sphinxadmonition}


\subsubsection{Document}
\label{\detokenize{fr/01_bienvenue:document}}
\sphinxAtStartPar
Un \sphinxstylestrong{Document} est tout fichier téléversé qui complète les données spatiales : un
rapport PDF, une photographie de terrain, un tableur, une présentation ou tout autre
matériel de référence.

\sphinxAtStartPar
Les documents sont des ressources à part entière dans KMaP — ils ont leur propre page
de métadonnées, peuvent être recherchés et filtrés, et peuvent être liés à des Datasets
ou des Cartes.

\sphinxAtStartPar
Les formats pris en charge comprennent : \sphinxcode{\sphinxupquote{.pdf}}, \sphinxcode{\sphinxupquote{.docx}}, \sphinxcode{\sphinxupquote{.xlsx}}, \sphinxcode{\sphinxupquote{.pptx}},
\sphinxcode{\sphinxupquote{.jpg}}, \sphinxcode{\sphinxupquote{.png}}, \sphinxcode{\sphinxupquote{.tif}}, \sphinxcode{\sphinxupquote{.mp4}}, \sphinxcode{\sphinxupquote{.zip}} et bien d’autres.


\subsubsection{Carte (Map)}
\label{\detokenize{fr/01_bienvenue:carte-map}}
\sphinxAtStartPar
Une \sphinxstylestrong{Carte} est une carte web interactive à plusieurs couches, créée en composant
un ou plusieurs Datasets. Les cartes sont créées et visualisées via le visualiseur de
carte intégré \sphinxstylestrong{MapStore}.

\sphinxAtStartPar
Chaque Carte :
\begin{itemize}
\item {} 
\sphinxAtStartPar
Combine plusieurs couches de Datasets avec des styles et des paramètres de visibilité
individuels

\item {} 
\sphinxAtStartPar
Inclut un fond de carte (ex. OpenStreetMap ou imagerie satellitaire)

\item {} 
\sphinxAtStartPar
Prend en charge les outils de navigation, d’interrogation, de mesure et d’impression

\item {} 
\sphinxAtStartPar
Peut être partagée via un lien permanent ou intégrée dans des sites web externes

\end{itemize}


\subsubsection{GeoStory}
\label{\detokenize{fr/01_bienvenue:geostory}}
\sphinxAtStartPar
Une \sphinxstylestrong{GeoStory} est un récit multimédia défilant qui entrelace texte, cartes
interactives, images et vidéos — similaire à une story map.

\sphinxAtStartPar
Les GeoStories sont composées de sections :


\begin{savenotes}\sphinxattablestart
\sphinxthistablewithglobalstyle
\centering
\begin{tabular}[t]{\X{25}{100}\X{75}{100}}
\sphinxtoprule
\sphinxstyletheadfamily 
\sphinxAtStartPar
Type de section
&\sphinxstyletheadfamily 
\sphinxAtStartPar
Description
\\
\sphinxmidrule
\sphinxtableatstartofbodyhook
\sphinxAtStartPar
\sphinxstylestrong{Titre}
&
\sphinxAtStartPar
Panneau d’ouverture plein écran avec image ou vidéo de fond.
\\
\sphinxhline
\sphinxAtStartPar
\sphinxstylestrong{Bannière}
&
\sphinxAtStartPar
Séparateur horizontal ou en\sphinxhyphen{}tête de chapitre avec texte et/ou image.
\\
\sphinxhline
\sphinxAtStartPar
\sphinxstylestrong{Paragraphe}
&
\sphinxAtStartPar
Texte libre et contenu multimédia (image, vidéo, carte).
\\
\sphinxhline
\sphinxAtStartPar
\sphinxstylestrong{Immersif}
&
\sphinxAtStartPar
Carte interactive plein écran sur laquelle le texte narratif défile.
\\
\sphinxhline
\sphinxAtStartPar
\sphinxstylestrong{GeoCarousel}
&
\sphinxAtStartPar
Diaporama lié à des emplacements géographiques.
\\
\sphinxhline
\sphinxAtStartPar
\sphinxstylestrong{Page web}
&
\sphinxAtStartPar
Page web externe intégrée directement dans le récit.
\\
\sphinxbottomrule
\end{tabular}
\sphinxtableafterendhook\par
\sphinxattableend\end{savenotes}

\sphinxAtStartPar
Les GeoStories sont idéales pour communiquer des thématiques environnementales à un
public non spécialisé.


\subsubsection{Tableau de bord (Dashboard)}
\label{\detokenize{fr/01_bienvenue:tableau-de-bord-dashboard}}
\sphinxAtStartPar
Un \sphinxstylestrong{Dashboard} est un panneau configurable de widgets interactifs — graphiques,
tableaux, cartes, compteurs et blocs de texte — liés à un ou plusieurs Datasets.

\sphinxAtStartPar
Types de widgets disponibles :


\begin{savenotes}\sphinxattablestart
\sphinxthistablewithglobalstyle
\centering
\begin{tabular}[t]{\X{25}{100}\X{75}{100}}
\sphinxtoprule
\sphinxstyletheadfamily 
\sphinxAtStartPar
Widget
&\sphinxstyletheadfamily 
\sphinxAtStartPar
Description
\\
\sphinxmidrule
\sphinxtableatstartofbodyhook
\sphinxAtStartPar
\sphinxstylestrong{Carte}
&
\sphinxAtStartPar
Mini\sphinxhyphen{}carte interactive affichant des couches de datasets.
\\
\sphinxhline
\sphinxAtStartPar
\sphinxstylestrong{Graphique}
&
\sphinxAtStartPar
Graphiques en barres, en courbes, circulaires ou en nuage de points dérivés
des attributs du dataset.
\\
\sphinxhline
\sphinxAtStartPar
\sphinxstylestrong{Tableau}
&
\sphinxAtStartPar
Vue tabulaire triable et filtrable des attributs du dataset.
\\
\sphinxhline
\sphinxAtStartPar
\sphinxstylestrong{Compteur}
&
\sphinxAtStartPar
Statistique de synthèse en grand format (comptage, somme, moyenne, …).
\\
\sphinxhline
\sphinxAtStartPar
\sphinxstylestrong{Texte}
&
\sphinxAtStartPar
Bloc de texte libre pour les titres ou les commentaires.
\\
\sphinxbottomrule
\end{tabular}
\sphinxtableafterendhook\par
\sphinxattableend\end{savenotes}

\sphinxAtStartPar
Les widgets d’un Dashboard sont interconnectés : sélectionner une région sur le widget
carte met automatiquement à jour tous les autres widgets pour n’afficher que les données
de cette zone.


\subsection{Services Interopérables}
\label{\detokenize{fr/01_bienvenue:services-interoperables}}
\sphinxAtStartPar
KMaP expose ses données via des \sphinxstylestrong{services web OGC} standard, permettant à tout
logiciel SIG compatible ou application web d’accéder directement aux données sans
télécharger de fichiers.


\begin{savenotes}\sphinxattablestart
\sphinxthistablewithglobalstyle
\centering
\begin{tabular}[t]{\X{15}{100}\X{15}{100}\X{70}{100}}
\sphinxtoprule
\sphinxstyletheadfamily 
\sphinxAtStartPar
Service
&\sphinxstyletheadfamily 
\sphinxAtStartPar
Standard
&\sphinxstyletheadfamily 
\sphinxAtStartPar
Description
\\
\sphinxmidrule
\sphinxtableatstartofbodyhook
\sphinxAtStartPar
\sphinxstylestrong{WMS}
&
\sphinxAtStartPar
OGC Web Map Service
&
\sphinxAtStartPar
Renvoie les couches de datasets sous forme d’images de carte rendues. Idéal pour
la visualisation dans les SIG de bureau (QGIS, ArcGIS) ou les applications web.
\\
\sphinxhline
\sphinxAtStartPar
\sphinxstylestrong{WFS}
&
\sphinxAtStartPar
OGC Web Feature Service
&
\sphinxAtStartPar
Renvoie les données vectorielles brutes (géométries + attributs) pour le
téléchargement ou l’analyse. Prend en charge le filtrage par attribut ou par
emprise.
\\
\sphinxhline
\sphinxAtStartPar
\sphinxstylestrong{WCS}
&
\sphinxAtStartPar
OGC Web Coverage Service
&
\sphinxAtStartPar
Renvoie les données raster brutes (grilles) pour l’analyse. Disponible pour les
datasets raster.
\\
\sphinxhline
\sphinxAtStartPar
\sphinxstylestrong{WMTS}
&
\sphinxAtStartPar
OGC Web Map Tile Service
&
\sphinxAtStartPar
Renvoie des tuiles de carte pré\sphinxhyphen{}rendues pour un affichage rapide dans les cartes web.
\\
\sphinxhline
\sphinxAtStartPar
\sphinxstylestrong{CSW}
&
\sphinxAtStartPar
OGC Catalogue Service
&
\sphinxAtStartPar
Expose le catalogue de métadonnées de KMaP pour la recherche et la collecte par
des portails externes.
\\
\sphinxbottomrule
\end{tabular}
\sphinxtableafterendhook\par
\sphinxattableend\end{savenotes}
\subsubsection*{Comment obtenir les URL des services}

\sphinxAtStartPar
Sur la page de détail de tout Dataset, cliquez sur l’onglet \sphinxstylestrong{Services OGC} (ou
cherchez la section \sphinxstyleemphasis{Télécharger / OGC}). Les URL des points de terminaison WMS, WFS et
WCS y sont listées et peuvent être copiées directement dans QGIS ou tout autre client
compatible OGC.

\begin{sphinxadmonition}{tip}{Tip:}
\sphinxAtStartPar
Dans QGIS, allez dans \sphinxstylestrong{Couche \(\rightarrow\) Ajouter une couche \(\rightarrow\) Ajouter une couche WMS/WMTS},
collez l’URL WMS et parcourez la liste des couches disponibles pour ajouter directement
les données KMaP à votre projet sans télécharger aucun fichier.
\end{sphinxadmonition}


\subsection{Thématiques Disponibles sur KMaP}
\label{\detokenize{fr/01_bienvenue:thematiques-disponibles-sur-kmap}}
\sphinxAtStartPar
KMaP organise ses datasets dans les catégories thématiques suivantes :
\begin{itemize}
\item {} 
\sphinxAtStartPar
🐟  \sphinxstylestrong{Pêche et Aquaculture}

\item {} 
\sphinxAtStartPar
🌊  \sphinxstylestrong{Biodiversité Marine}

\item {} 
\sphinxAtStartPar
🏭  \sphinxstylestrong{Pollution}

\item {} 
\sphinxAtStartPar
🌡  \sphinxstylestrong{Changements Climatiques}

\item {} 
\sphinxAtStartPar
🗺  \sphinxstylestrong{Planification Spatiale Marine}

\item {} 
\sphinxAtStartPar
💙  \sphinxstylestrong{Durabilité et Économie Bleue}

\item {} 
\sphinxAtStartPar
🏛  \sphinxstylestrong{Gouvernance}

\end{itemize}

\sphinxAtStartPar
Les ressources relatives au Sanctuaire Pelagos et à l’Accord Pelagos apparaissent dans
plusieurs thématiques, notamment Biodiversité Marine, Gouvernance et Planification
Spatiale Marine.

\sphinxstepscope


\section{Comment Trouver et Visualiser Vos Données}
\label{\detokenize{fr/02_trouver_visualiser:comment-trouver-et-visualiser-vos-donnees}}\label{\detokenize{fr/02_trouver_visualiser:find-fr}}\label{\detokenize{fr/02_trouver_visualiser::doc}}
\begin{sphinxcontents}
\sphinxstylecontentstitle{Sur cette page}
\begin{itemize}
\item {} 
\sphinxAtStartPar
\phantomsection\label{\detokenize{fr/02_trouver_visualiser:id1}}{\hyperref[\detokenize{fr/02_trouver_visualiser:acceder-au-catalogue}]{\sphinxcrossref{Accéder au Catalogue}}}

\item {} 
\sphinxAtStartPar
\phantomsection\label{\detokenize{fr/02_trouver_visualiser:id2}}{\hyperref[\detokenize{fr/02_trouver_visualiser:filtrer-et-rechercher-dans-le-catalogue}]{\sphinxcrossref{Filtrer et Rechercher dans le Catalogue}}}
\begin{itemize}
\item {} 
\sphinxAtStartPar
\phantomsection\label{\detokenize{fr/02_trouver_visualiser:id3}}{\hyperref[\detokenize{fr/02_trouver_visualiser:utiliser-la-barre-de-recherche}]{\sphinxcrossref{Utiliser la Barre de Recherche}}}

\item {} 
\sphinxAtStartPar
\phantomsection\label{\detokenize{fr/02_trouver_visualiser:id4}}{\hyperref[\detokenize{fr/02_trouver_visualiser:utiliser-les-filtres}]{\sphinxcrossref{Utiliser les Filtres}}}

\end{itemize}

\item {} 
\sphinxAtStartPar
\phantomsection\label{\detokenize{fr/02_trouver_visualiser:id5}}{\hyperref[\detokenize{fr/02_trouver_visualiser:exemple-pas-a-pas-trouver-un-dataset-pelagos}]{\sphinxcrossref{Exemple Pas à Pas : Trouver un Dataset Pelagos}}}
\begin{itemize}
\item {} 
\sphinxAtStartPar
\phantomsection\label{\detokenize{fr/02_trouver_visualiser:id6}}{\hyperref[\detokenize{fr/02_trouver_visualiser:etape-1-ouvrir-le-catalogue}]{\sphinxcrossref{Étape 1 \textendash{} Ouvrir le catalogue}}}

\item {} 
\sphinxAtStartPar
\phantomsection\label{\detokenize{fr/02_trouver_visualiser:id7}}{\hyperref[\detokenize{fr/02_trouver_visualiser:etape-2-rechercher-par-mot-cle}]{\sphinxcrossref{Étape 2 \textendash{} Rechercher par mot\sphinxhyphen{}clé}}}

\item {} 
\sphinxAtStartPar
\phantomsection\label{\detokenize{fr/02_trouver_visualiser:id8}}{\hyperref[\detokenize{fr/02_trouver_visualiser:etape-3-filtrer-uniquement-les-datasets}]{\sphinxcrossref{Étape 3 \textendash{} Filtrer uniquement les Datasets}}}

\item {} 
\sphinxAtStartPar
\phantomsection\label{\detokenize{fr/02_trouver_visualiser:id9}}{\hyperref[\detokenize{fr/02_trouver_visualiser:etape-4-selectionner-un-dataset}]{\sphinxcrossref{Étape 4 \textendash{} Sélectionner un Dataset}}}

\end{itemize}

\item {} 
\sphinxAtStartPar
\phantomsection\label{\detokenize{fr/02_trouver_visualiser:id10}}{\hyperref[\detokenize{fr/02_trouver_visualiser:la-page-de-detail-d-un-dataset}]{\sphinxcrossref{La Page de Détail d’un Dataset}}}
\begin{itemize}
\item {} 
\sphinxAtStartPar
\phantomsection\label{\detokenize{fr/02_trouver_visualiser:id11}}{\hyperref[\detokenize{fr/02_trouver_visualiser:miniature-et-actions}]{\sphinxcrossref{Miniature et Actions}}}

\item {} 
\sphinxAtStartPar
\phantomsection\label{\detokenize{fr/02_trouver_visualiser:id12}}{\hyperref[\detokenize{fr/02_trouver_visualiser:apercu-interactif}]{\sphinxcrossref{Aperçu Interactif}}}

\item {} 
\sphinxAtStartPar
\phantomsection\label{\detokenize{fr/02_trouver_visualiser:id13}}{\hyperref[\detokenize{fr/02_trouver_visualiser:panneau-des-metadonnees}]{\sphinxcrossref{Panneau des Métadonnées}}}

\item {} 
\sphinxAtStartPar
\phantomsection\label{\detokenize{fr/02_trouver_visualiser:id14}}{\hyperref[\detokenize{fr/02_trouver_visualiser:vue-complete-des-metadonnees}]{\sphinxcrossref{Vue Complète des Métadonnées}}}

\item {} 
\sphinxAtStartPar
\phantomsection\label{\detokenize{fr/02_trouver_visualiser:id15}}{\hyperref[\detokenize{fr/02_trouver_visualiser:telecharger-un-dataset}]{\sphinxcrossref{Télécharger un Dataset}}}

\end{itemize}

\end{itemize}
\end{sphinxcontents}

\sphinxAtStartPar
Cette section explique comment découvrir des ressources dans KMaP, appliquer des filtres
pour affiner votre recherche, et explorer la page de détail et l’aperçu d’un Dataset.


\subsection{Accéder au Catalogue}
\label{\detokenize{fr/02_trouver_visualiser:acceder-au-catalogue}}
\sphinxAtStartPar
Le catalogue KMaP liste toutes les ressources publiques et toutes les ressources
partagées avec votre groupe. Il existe deux façons principales d’y accéder :
\begin{enumerate}
\sphinxsetlistlabels{\arabic}{enumi}{enumii}{}{.}%
\item {} 
\sphinxAtStartPar
\sphinxstylestrong{Depuis la barre de navigation} \textendash{} cliquez sur \sphinxstylestrong{Maps} (pour les ressources
spatiales) ou \sphinxstylestrong{Library} (pour les documents et publications) dans le menu du haut.

\item {} 
\sphinxAtStartPar
\sphinxstylestrong{Directement via URL} \textendash{} rendez\sphinxhyphen{}vous sur \sphinxurl{https://kmap.info-rac.org/resources/}

\end{enumerate}

\sphinxAtStartPar
La page du catalogue affiche les ressources sous forme de cartes, chacune montrant une
miniature, un titre, le type de ressource et des métadonnées récapitulatives.

\begin{sphinxadmonition}{tip}{Tip:}
\sphinxAtStartPar
Assurez\sphinxhyphen{}vous d’être \sphinxstylestrong{connecté} pour voir les ressources partagées exclusivement
avec le groupe Accord Pelagos. Les ressources restreintes au groupe ne sont pas
visibles pour les utilisateurs non authentifiés.
\end{sphinxadmonition}


\subsection{Filtrer et Rechercher dans le Catalogue}
\label{\detokenize{fr/02_trouver_visualiser:filtrer-et-rechercher-dans-le-catalogue}}

\subsubsection{Utiliser la Barre de Recherche}
\label{\detokenize{fr/02_trouver_visualiser:utiliser-la-barre-de-recherche}}
\sphinxAtStartPar
La barre de recherche en haut de la page du catalogue accepte des requêtes en texte
libre. Elle effectue la recherche dans les titres, les résumés et les champs de
mots\sphinxhyphen{}clés.

\sphinxAtStartPar
Pour trouver des contenus liés à Pelagos :
\begin{enumerate}
\sphinxsetlistlabels{\arabic}{enumi}{enumii}{}{.}%
\item {} 
\sphinxAtStartPar
Tapez \sphinxcode{\sphinxupquote{pelagos sanctuary}} dans la barre de recherche et appuyez sur \sphinxstylestrong{Entrée}.

\item {} 
\sphinxAtStartPar
Alternativement, recherchez \sphinxcode{\sphinxupquote{pelagos agreement}} pour récupérer les ressources
étiquetées avec ce mot\sphinxhyphen{}clé.

\end{enumerate}

\begin{sphinxadmonition}{note}{Note:}
\sphinxAtStartPar
\sphinxcode{\sphinxupquote{pelagos sanctuary}} et \sphinxcode{\sphinxupquote{pelagos agreement}} sont des mots\sphinxhyphen{}clés contrôlés utilisés
par le groupe Accord Pelagos lors de la publication de ressources sur KMaP. Utiliser
ces termes exacts donnera les résultats les plus pertinents.
\end{sphinxadmonition}


\subsubsection{Utiliser les Filtres}
\label{\detokenize{fr/02_trouver_visualiser:utiliser-les-filtres}}
\sphinxAtStartPar
Sur le côté gauche de la page du catalogue se trouve un panneau de filtres. Les filtres
peuvent être combinés librement et sont appliqués immédiatement.


\begin{savenotes}\sphinxattablestart
\sphinxthistablewithglobalstyle
\centering
\begin{tabular}[t]{\X{25}{100}\X{75}{100}}
\sphinxtoprule
\sphinxstyletheadfamily 
\sphinxAtStartPar
Filtre
&\sphinxstyletheadfamily 
\sphinxAtStartPar
Comment l’utiliser
\\
\sphinxmidrule
\sphinxtableatstartofbodyhook
\sphinxAtStartPar
\sphinxstylestrong{Type de ressource}
&
\sphinxAtStartPar
Sélectionnez \sphinxstyleemphasis{Dataset}, \sphinxstyleemphasis{Document}, \sphinxstyleemphasis{Map}, \sphinxstyleemphasis{GeoStory} ou \sphinxstyleemphasis{Dashboard} pour
limiter les résultats à un seul type.
\\
\sphinxhline
\sphinxAtStartPar
\sphinxstylestrong{Mots\sphinxhyphen{}clés}
&
\sphinxAtStartPar
Tapez ou sélectionnez un mot\sphinxhyphen{}clé (ex. \sphinxcode{\sphinxupquote{pelagos sanctuary}}) pour filtrer par
étiquette thématique.
\\
\sphinxhline
\sphinxAtStartPar
\sphinxstylestrong{Catégorie / Thème}
&
\sphinxAtStartPar
Sélectionnez une catégorie thématique comme \sphinxstyleemphasis{Biodiversité Marine} ou
\sphinxstyleemphasis{Gouvernance}.
\\
\sphinxhline
\sphinxAtStartPar
\sphinxstylestrong{Date}
&
\sphinxAtStartPar
Filtrez par plage de dates de publication ou de modification.
\\
\sphinxhline
\sphinxAtStartPar
\sphinxstylestrong{Propriétaire}
&
\sphinxAtStartPar
Filtrez par l’utilisateur ou l’organisation qui a publié la ressource.
\\
\sphinxhline
\sphinxAtStartPar
\sphinxstylestrong{Emprise}
&
\sphinxAtStartPar
Dessinez un rectangle sur la mini\sphinxhyphen{}carte pour filtrer par zone géographique.
\\
\sphinxbottomrule
\end{tabular}
\sphinxtableafterendhook\par
\sphinxattableend\end{savenotes}

\begin{sphinxadmonition}{tip}{Tip:}
\sphinxAtStartPar
Pour trouver rapidement tous les datasets liés à Pelagos :
\begin{enumerate}
\sphinxsetlistlabels{\arabic}{enumi}{enumii}{}{.}%
\item {} 
\sphinxAtStartPar
Définissez \sphinxstylestrong{Type de ressource} \(\rightarrow\) \sphinxstyleemphasis{Dataset}

\item {} 
\sphinxAtStartPar
Dans le filtre \sphinxstylestrong{Mots\sphinxhyphen{}clés}, tapez \sphinxcode{\sphinxupquote{pelagos sanctuary}}

\item {} 
\sphinxAtStartPar
Le catalogue se mettra à jour pour n’afficher que les datasets étiquetés avec
ce mot\sphinxhyphen{}clé.

\end{enumerate}
\end{sphinxadmonition}


\subsection{Exemple Pas à Pas : Trouver un Dataset Pelagos}
\label{\detokenize{fr/02_trouver_visualiser:exemple-pas-a-pas-trouver-un-dataset-pelagos}}
\sphinxAtStartPar
L’exemple suivant montre comment trouver, prévisualiser et lire les métadonnées d’un
Dataset pertinent pour l’Accord Pelagos.


\subsubsection{Étape 1 \textendash{} Ouvrir le catalogue}
\label{\detokenize{fr/02_trouver_visualiser:etape-1-ouvrir-le-catalogue}}
\sphinxAtStartPar
Rendez\sphinxhyphen{}vous sur \sphinxurl{https://kmap.info-rac.org} et cliquez sur \sphinxstylestrong{Maps} dans la barre de
navigation en haut. Le catalogue du Data Hub s’ouvre, filtré sur les types de ressources
spatiales.


\subsubsection{Étape 2 \textendash{} Rechercher par mot\sphinxhyphen{}clé}
\label{\detokenize{fr/02_trouver_visualiser:etape-2-rechercher-par-mot-cle}}
\sphinxAtStartPar
Dans la barre de recherche, tapez

\begin{sphinxVerbatim}[commandchars=\\\{\}]
\PYG{n}{pelagos} \PYG{n}{sanctuary}
\end{sphinxVerbatim}

\sphinxAtStartPar
Appuyez sur \sphinxstylestrong{Entrée}. Le panneau des résultats se met à jour pour afficher les
ressources correspondant à ce mot\sphinxhyphen{}clé.


\subsubsection{Étape 3 \textendash{} Filtrer uniquement les Datasets}
\label{\detokenize{fr/02_trouver_visualiser:etape-3-filtrer-uniquement-les-datasets}}
\sphinxAtStartPar
Dans le panneau de filtres à gauche, sous \sphinxstylestrong{Type de ressource}, cochez \sphinxstylestrong{Dataset}.
La liste n’affiche maintenant que les datasets étiquetés avec \sphinxstyleemphasis{pelagos sanctuary}.


\subsubsection{Étape 4 \textendash{} Sélectionner un Dataset}
\label{\detokenize{fr/02_trouver_visualiser:etape-4-selectionner-un-dataset}}
\sphinxAtStartPar
Cliquez sur une carte de dataset pour ouvrir sa \sphinxstylestrong{page de détail}.


\subsection{La Page de Détail d’un Dataset}
\label{\detokenize{fr/02_trouver_visualiser:la-page-de-detail-d-un-dataset}}
\sphinxAtStartPar
La page de détail est la page d’information centrale de toute ressource. Elle est
divisée en plusieurs sections.


\subsubsection{Miniature et Actions}
\label{\detokenize{fr/02_trouver_visualiser:miniature-et-actions}}
\sphinxAtStartPar
En haut de la page, vous trouverez :
\begin{itemize}
\item {} 
\sphinxAtStartPar
Une \sphinxstylestrong{miniature de carte} ou une image d’aperçu du dataset.

\item {} 
\sphinxAtStartPar
Des boutons d’action :
\begin{itemize}
\item {} 
\sphinxAtStartPar
\sphinxstylestrong{View map} \textendash{} ouvre le dataset dans le visualiseur de carte interactif.

\item {} 
\sphinxAtStartPar
\sphinxstylestrong{Download} \textendash{} télécharge le dataset dans le format choisi.

\item {} 
\sphinxAtStartPar
\sphinxstylestrong{Edit} (si vous avez l’autorisation) \textendash{} ouvre l’éditeur de métadonnées.

\item {} 
\sphinxAtStartPar
\sphinxstylestrong{Share} \textendash{} copie le lien permanent vers la ressource.

\end{itemize}

\end{itemize}


\subsubsection{Aperçu Interactif}
\label{\detokenize{fr/02_trouver_visualiser:apercu-interactif}}
\sphinxAtStartPar
Sous la miniature, le \sphinxstylestrong{panneau d’aperçu de carte} affiche le dataset rendu sur un
fond de carte. Vous pouvez :
\begin{itemize}
\item {} 
\sphinxAtStartPar
Vous déplacer et zoomer avec la souris ou le toucher.

\item {} 
\sphinxAtStartPar
Cliquer sur un élément pour ouvrir une \sphinxstylestrong{fenêtre contextuelle} affichant ses valeurs
d’attributs.

\item {} 
\sphinxAtStartPar
Utiliser le panneau de couches (icône ☰) pour basculer la visibilité ou consulter
la légende.

\end{itemize}

\begin{sphinxadmonition}{note}{Note:}
\sphinxAtStartPar
L’aperçu est une carte interactive — pas une image statique. Vous pouvez explorer
les données spatialement avant de décider de les télécharger ou de les utiliser.
\end{sphinxadmonition}


\subsubsection{Panneau des Métadonnées}
\label{\detokenize{fr/02_trouver_visualiser:panneau-des-metadonnees}}
\sphinxAtStartPar
Sous l’aperçu, le \sphinxstylestrong{panneau des métadonnées} contient des informations détaillées sur
le dataset. Les champs principaux sont :


\begin{savenotes}\sphinxattablestart
\sphinxthistablewithglobalstyle
\centering
\begin{tabular}[t]{\X{30}{100}\X{70}{100}}
\sphinxtoprule
\sphinxstyletheadfamily 
\sphinxAtStartPar
Champ
&\sphinxstyletheadfamily 
\sphinxAtStartPar
Description
\\
\sphinxmidrule
\sphinxtableatstartofbodyhook
\sphinxAtStartPar
\sphinxstylestrong{Titre}
&
\sphinxAtStartPar
Le nom complet du dataset.
\\
\sphinxhline
\sphinxAtStartPar
\sphinxstylestrong{Résumé (Abstract)}
&
\sphinxAtStartPar
Une description en langage courant du contenu et de l’objet du dataset.
\\
\sphinxhline
\sphinxAtStartPar
\sphinxstylestrong{Mots\sphinxhyphen{}clés}
&
\sphinxAtStartPar
Étiquettes thématiques et géographiques (cherchez \sphinxcode{\sphinxupquote{pelagos sanctuary}} ou
\sphinxcode{\sphinxupquote{pelagos agreement}}).
\\
\sphinxhline
\sphinxAtStartPar
\sphinxstylestrong{Catégorie}
&
\sphinxAtStartPar
La catégorie INSPIRE ou thématique (ex. \sphinxstyleemphasis{Biodiversité Marine}).
\\
\sphinxhline
\sphinxAtStartPar
\sphinxstylestrong{Date}
&
\sphinxAtStartPar
Date de publication et date de dernière modification.
\\
\sphinxhline
\sphinxAtStartPar
\sphinxstylestrong{Emprise spatiale}
&
\sphinxAtStartPar
Rectangle d’emprise ou région géographique couverte par les données.
\\
\sphinxhline
\sphinxAtStartPar
\sphinxstylestrong{Système de Référence de Coordonnées (SRC)}
&
\sphinxAtStartPar
La projection utilisée (ex. WGS 84 / EPSG:4326).
\\
\sphinxhline
\sphinxAtStartPar
\sphinxstylestrong{Propriétaire}
&
\sphinxAtStartPar
L’utilisateur ou l’organisation responsable du dataset.
\\
\sphinxhline
\sphinxAtStartPar
\sphinxstylestrong{Licence}
&
\sphinxAtStartPar
La licence d’utilisation (ex. Creative Commons, ODbL, …).
\\
\sphinxhline
\sphinxAtStartPar
\sphinxstylestrong{Point de contact}
&
\sphinxAtStartPar
Qui contacter pour des questions sur les données.
\\
\sphinxhline
\sphinxAtStartPar
\sphinxstylestrong{Services OGC}
&
\sphinxAtStartPar
URL des points de terminaison WMS, WFS et WCS pour l’intégration dans un
logiciel SIG.
\\
\sphinxhline
\sphinxAtStartPar
\sphinxstylestrong{Ressources liées}
&
\sphinxAtStartPar
Liens vers des Cartes, Documents ou autres Datasets connectés à cette ressource.
\\
\sphinxbottomrule
\end{tabular}
\sphinxtableafterendhook\par
\sphinxattableend\end{savenotes}


\subsubsection{Vue Complète des Métadonnées}
\label{\detokenize{fr/02_trouver_visualiser:vue-complete-des-metadonnees}}
\sphinxAtStartPar
Pour voir l’enregistrement complet des métadonnées, cliquez sur l’onglet \sphinxstylestrong{Métadonnées}
(ou le lien \sphinxstyleemphasis{Voir les métadonnées complètes} près du bas de la page de détail). Cette
vue affiche tous les champs de métadonnées dans leur forme complète, y compris les
champs conformes aux normes ISO 19115 / INSPIRE.


\subsubsection{Télécharger un Dataset}
\label{\detokenize{fr/02_trouver_visualiser:telecharger-un-dataset}}
\sphinxAtStartPar
Cliquez sur le bouton \sphinxstylestrong{Download} de la page de détail. Un sélecteur de format
apparaît. Les options courantes comprennent :
\begin{itemize}
\item {} 
\sphinxAtStartPar
\sphinxstylestrong{ESRI Shapefile} \textendash{} pour une utilisation dans QGIS, ArcGIS et la plupart des SIG
de bureau.

\item {} 
\sphinxAtStartPar
\sphinxstylestrong{GeoJSON} \textendash{} pour la cartographie web et une utilisation SIG légère.

\item {} 
\sphinxAtStartPar
\sphinxstylestrong{KML/KMZ} \textendash{} pour Google Earth.

\item {} 
\sphinxAtStartPar
\sphinxstylestrong{GeoPackage} \textendash{} format mono\sphinxhyphen{}fichier pris en charge par QGIS.

\item {} 
\sphinxAtStartPar
\sphinxstylestrong{GeoTIFF} \textendash{} pour les datasets raster.

\item {} 
\sphinxAtStartPar
\sphinxstylestrong{CSV} \textendash{} table attributaire uniquement (sans géométrie).

\end{itemize}

\sphinxAtStartPar
Sélectionnez votre format préféré et cliquez sur \sphinxstylestrong{Download} pour enregistrer le fichier.

\sphinxstepscope


\section{Gérer les Métadonnées}
\label{\detokenize{fr/03_gerer_metadonnees:gerer-les-metadonnees}}\label{\detokenize{fr/03_gerer_metadonnees:metadata-fr}}\label{\detokenize{fr/03_gerer_metadonnees::doc}}
\begin{sphinxcontents}
\sphinxstylecontentstitle{Sur cette page}
\begin{itemize}
\item {} 
\sphinxAtStartPar
\phantomsection\label{\detokenize{fr/03_gerer_metadonnees:id1}}{\hyperref[\detokenize{fr/03_gerer_metadonnees:qui-peut-modifier-les-metadonnees}]{\sphinxcrossref{Qui Peut Modifier les Métadonnées ?}}}

\item {} 
\sphinxAtStartPar
\phantomsection\label{\detokenize{fr/03_gerer_metadonnees:id2}}{\hyperref[\detokenize{fr/03_gerer_metadonnees:ouvrir-l-editeur-de-metadonnees}]{\sphinxcrossref{Ouvrir l’Éditeur de Métadonnées}}}

\item {} 
\sphinxAtStartPar
\phantomsection\label{\detokenize{fr/03_gerer_metadonnees:id3}}{\hyperref[\detokenize{fr/03_gerer_metadonnees:apercu-de-l-editeur-de-metadonnees}]{\sphinxcrossref{Aperçu de l’Éditeur de Métadonnées}}}

\item {} 
\sphinxAtStartPar
\phantomsection\label{\detokenize{fr/03_gerer_metadonnees:id4}}{\hyperref[\detokenize{fr/03_gerer_metadonnees:modification-des-champs-de-metadonnees-principaux}]{\sphinxcrossref{Modification des Champs de Métadonnées Principaux}}}
\begin{itemize}
\item {} 
\sphinxAtStartPar
\phantomsection\label{\detokenize{fr/03_gerer_metadonnees:id5}}{\hyperref[\detokenize{fr/03_gerer_metadonnees:titre}]{\sphinxcrossref{Titre}}}

\item {} 
\sphinxAtStartPar
\phantomsection\label{\detokenize{fr/03_gerer_metadonnees:id6}}{\hyperref[\detokenize{fr/03_gerer_metadonnees:resume-abstract}]{\sphinxcrossref{Résumé (Abstract)}}}

\item {} 
\sphinxAtStartPar
\phantomsection\label{\detokenize{fr/03_gerer_metadonnees:id7}}{\hyperref[\detokenize{fr/03_gerer_metadonnees:mots-cles}]{\sphinxcrossref{Mots\sphinxhyphen{}clés}}}

\item {} 
\sphinxAtStartPar
\phantomsection\label{\detokenize{fr/03_gerer_metadonnees:id8}}{\hyperref[\detokenize{fr/03_gerer_metadonnees:categorie}]{\sphinxcrossref{Catégorie}}}

\item {} 
\sphinxAtStartPar
\phantomsection\label{\detokenize{fr/03_gerer_metadonnees:id9}}{\hyperref[\detokenize{fr/03_gerer_metadonnees:emprise-spatiale}]{\sphinxcrossref{Emprise Spatiale}}}

\item {} 
\sphinxAtStartPar
\phantomsection\label{\detokenize{fr/03_gerer_metadonnees:id10}}{\hyperref[\detokenize{fr/03_gerer_metadonnees:champs-de-date}]{\sphinxcrossref{Champs de Date}}}

\item {} 
\sphinxAtStartPar
\phantomsection\label{\detokenize{fr/03_gerer_metadonnees:id11}}{\hyperref[\detokenize{fr/03_gerer_metadonnees:licence}]{\sphinxcrossref{Licence}}}

\item {} 
\sphinxAtStartPar
\phantomsection\label{\detokenize{fr/03_gerer_metadonnees:id12}}{\hyperref[\detokenize{fr/03_gerer_metadonnees:point-de-contact}]{\sphinxcrossref{Point de Contact}}}

\item {} 
\sphinxAtStartPar
\phantomsection\label{\detokenize{fr/03_gerer_metadonnees:id13}}{\hyperref[\detokenize{fr/03_gerer_metadonnees:miniature}]{\sphinxcrossref{Miniature}}}

\end{itemize}

\item {} 
\sphinxAtStartPar
\phantomsection\label{\detokenize{fr/03_gerer_metadonnees:id14}}{\hyperref[\detokenize{fr/03_gerer_metadonnees:enregistrer-les-modifications}]{\sphinxcrossref{Enregistrer les Modifications}}}

\item {} 
\sphinxAtStartPar
\phantomsection\label{\detokenize{fr/03_gerer_metadonnees:id15}}{\hyperref[\detokenize{fr/03_gerer_metadonnees:liste-de-controle-de-la-completude-des-metadonnees}]{\sphinxcrossref{Liste de Contrôle de la Complétude des Métadonnées}}}

\end{itemize}
\end{sphinxcontents}

\sphinxAtStartPar
Des métadonnées de qualité rendent une ressource trouvable et fiable. Cette section
explique comment consulter et modifier les métadonnées d’un Dataset ou d’un Document
dans KMaP.

\begin{sphinxadmonition}{note}{Note:}
\sphinxAtStartPar
Vous pouvez uniquement modifier les métadonnées des ressources dont vous êtes le
\sphinxstylestrong{propriétaire} ou pour lesquelles des \sphinxstylestrong{droits de modification} vous ont été
accordés. Les membres du groupe Accord Pelagos peuvent avoir des droits de
modification sur les ressources partagées au sein du groupe — consultez votre
administrateur de groupe en cas de doute.
\end{sphinxadmonition}


\subsection{Qui Peut Modifier les Métadonnées ?}
\label{\detokenize{fr/03_gerer_metadonnees:qui-peut-modifier-les-metadonnees}}
\sphinxAtStartPar
KMaP dispose d’un modèle de permissions à plusieurs niveaux :


\begin{savenotes}\sphinxattablestart
\sphinxthistablewithglobalstyle
\centering
\begin{tabular}[t]{\X{25}{100}\X{75}{100}}
\sphinxtoprule
\sphinxstyletheadfamily 
\sphinxAtStartPar
Rôle
&\sphinxstyletheadfamily 
\sphinxAtStartPar
Droits de modification des métadonnées
\\
\sphinxmidrule
\sphinxtableatstartofbodyhook
\sphinxAtStartPar
\sphinxstylestrong{Propriétaire de la ressource}
&
\sphinxAtStartPar
Droits de modification complets sur tous les champs de métadonnées.
\\
\sphinxhline
\sphinxAtStartPar
\sphinxstylestrong{Gestionnaire de groupe}
&
\sphinxAtStartPar
Peut modifier les ressources appartenant à n’importe quel membre de son groupe.
\\
\sphinxhline
\sphinxAtStartPar
\sphinxstylestrong{Membre du groupe (éditeur)}
&
\sphinxAtStartPar
Peut modifier les ressources explicitement partagées avec le groupe avec la
permission de \sphinxstyleemphasis{modification}.
\\
\sphinxhline
\sphinxAtStartPar
\sphinxstylestrong{Visualiseur / public}
&
\sphinxAtStartPar
Accès en lecture seule. Ne peut pas modifier les métadonnées.
\\
\sphinxbottomrule
\end{tabular}
\sphinxtableafterendhook\par
\sphinxattableend\end{savenotes}

\sphinxAtStartPar
Si vous voyez un bouton \sphinxstylestrong{Edit} sur la page de détail d’une ressource, vous avez le
droit de modifier les métadonnées de cette ressource.


\subsection{Ouvrir l’Éditeur de Métadonnées}
\label{\detokenize{fr/03_gerer_metadonnees:ouvrir-l-editeur-de-metadonnees}}
\sphinxAtStartPar
Il existe deux façons d’ouvrir l’éditeur de métadonnées :

\sphinxAtStartPar
\sphinxstylestrong{Depuis la page de détail :}
\begin{enumerate}
\sphinxsetlistlabels{\arabic}{enumi}{enumii}{}{.}%
\item {} 
\sphinxAtStartPar
Naviguez vers la page de détail de la ressource (voir {\hyperref[\detokenize{fr/02_trouver_visualiser:find-fr}]{\sphinxcrossref{\DUrole{std}{\DUrole{std-ref}{Comment Trouver et Visualiser Vos Données}}}}} pour trouver
une ressource).

\item {} 
\sphinxAtStartPar
Cliquez sur le bouton \sphinxstylestrong{Edit} (icône crayon) dans la barre d’actions en haut à
droite.

\item {} 
\sphinxAtStartPar
Un menu déroulant apparaît. Sélectionnez \sphinxstylestrong{Edit metadata}.

\end{enumerate}

\sphinxAtStartPar
\sphinxstylestrong{Depuis le catalogue :}
\begin{enumerate}
\sphinxsetlistlabels{\arabic}{enumi}{enumii}{}{.}%
\item {} 
\sphinxAtStartPar
Dans le catalogue, survolez une carte de ressource.

\item {} 
\sphinxAtStartPar
Cliquez sur le menu \sphinxstylestrong{⋮} (autres options) qui apparaît.

\item {} 
\sphinxAtStartPar
Sélectionnez \sphinxstylestrong{Edit metadata}.

\end{enumerate}

\sphinxAtStartPar
L’éditeur de métadonnées s’ouvre sous forme de formulaire à plusieurs onglets.


\subsection{Aperçu de l’Éditeur de Métadonnées}
\label{\detokenize{fr/03_gerer_metadonnees:apercu-de-l-editeur-de-metadonnees}}
\sphinxAtStartPar
L’éditeur de métadonnées est organisé en onglets. Les plus importants sont :


\begin{savenotes}\sphinxattablestart
\sphinxthistablewithglobalstyle
\centering
\begin{tabular}[t]{\X{25}{100}\X{75}{100}}
\sphinxtoprule
\sphinxstyletheadfamily 
\sphinxAtStartPar
Onglet
&\sphinxstyletheadfamily 
\sphinxAtStartPar
Contenu
\\
\sphinxmidrule
\sphinxtableatstartofbodyhook
\sphinxAtStartPar
\sphinxstylestrong{Basic info}
&
\sphinxAtStartPar
Titre, résumé, langue, catégorie, mots\sphinxhyphen{}clés, point de contact.
\\
\sphinxhline
\sphinxAtStartPar
\sphinxstylestrong{Location and licences}
&
\sphinxAtStartPar
Emprise spatiale, SRC, champs de date, licence, attribution.
\\
\sphinxhline
\sphinxAtStartPar
\sphinxstylestrong{Optional metadata}
&
\sphinxAtStartPar
Champs ISO 19115 supplémentaires (édition, objectif, informations
complémentaires, …).
\\
\sphinxhline
\sphinxAtStartPar
\sphinxstylestrong{Settings}
&
\sphinxAtStartPar
Permissions, statut mis en avant, statut publié.
\\
\sphinxbottomrule
\end{tabular}
\sphinxtableafterendhook\par
\sphinxattableend\end{savenotes}


\subsection{Modification des Champs de Métadonnées Principaux}
\label{\detokenize{fr/03_gerer_metadonnees:modification-des-champs-de-metadonnees-principaux}}

\subsubsection{Titre}
\label{\detokenize{fr/03_gerer_metadonnees:titre}}
\sphinxAtStartPar
Le \sphinxstylestrong{titre} est le principal moyen pour les utilisateurs de trouver votre ressource.
Il doit être :
\begin{itemize}
\item {} 
\sphinxAtStartPar
Descriptif et spécifique (évitez des titres génériques comme “Données” ou “Couche
cartographique”).

\item {} 
\sphinxAtStartPar
Cohérent avec les conventions de nommage utilisées par le groupe Accord Pelagos.

\end{itemize}

\sphinxAtStartPar
Pour le modifier : cliquez sur le champ \sphinxstylestrong{Titre} dans l’onglet \sphinxstyleemphasis{Basic info} et tapez
vos modifications.


\subsubsection{Résumé (Abstract)}
\label{\detokenize{fr/03_gerer_metadonnees:resume-abstract}}
\sphinxAtStartPar
Le \sphinxstylestrong{résumé} est une description en langage courant de la ressource. Un bon résumé
répond aux questions :
\begin{itemize}
\item {} 
\sphinxAtStartPar
Que contient cette ressource ?

\item {} 
\sphinxAtStartPar
Quelle zone géographique couvre\sphinxhyphen{}t\sphinxhyphen{}elle ?

\item {} 
\sphinxAtStartPar
Quelle période temporelle représente\sphinxhyphen{}t\sphinxhyphen{}elle ?

\item {} 
\sphinxAtStartPar
À quoi sert\sphinxhyphen{}elle ?

\item {} 
\sphinxAtStartPar
Quelles sont ses limitations connues ?

\end{itemize}

\sphinxAtStartPar
Le champ résumé prend en charge la mise en forme de base. Visez au moins deux ou trois
phrases.


\subsubsection{Mots\sphinxhyphen{}clés}
\label{\detokenize{fr/03_gerer_metadonnees:mots-cles}}
\sphinxAtStartPar
Les mots\sphinxhyphen{}clés sont le principal mécanisme de filtrage et de découverte des ressources
dans KMaP.

\begin{sphinxadmonition}{important}{Important:}
\sphinxAtStartPar
Toutes les ressources liées à l’Accord Pelagos \sphinxstylestrong{doivent inclure} au moins l’un
des mots\sphinxhyphen{}clés suivants :
\begin{itemize}
\item {} 
\sphinxAtStartPar
\sphinxcode{\sphinxupquote{pelagos sanctuary}}

\item {} 
\sphinxAtStartPar
\sphinxcode{\sphinxupquote{pelagos agreement}}

\end{itemize}

\sphinxAtStartPar
Sans ces mots\sphinxhyphen{}clés, votre ressource n’apparaîtra pas lorsque d’autres membres du
groupe filtrent le catalogue avec ces termes.
\end{sphinxadmonition}

\sphinxAtStartPar
Pour ajouter un mot\sphinxhyphen{}clé :
\begin{enumerate}
\sphinxsetlistlabels{\arabic}{enumi}{enumii}{}{.}%
\item {} 
\sphinxAtStartPar
Dans le champ \sphinxstylestrong{Mots\sphinxhyphen{}clés}, commencez à taper le mot\sphinxhyphen{}clé.

\item {} 
\sphinxAtStartPar
S’il existe déjà dans le système, il apparaîtra dans un menu déroulant — sélectionnez\sphinxhyphen{}le.

\item {} 
\sphinxAtStartPar
S’il est nouveau, tapez le mot\sphinxhyphen{}clé complet et appuyez sur \sphinxstylestrong{Entrée} ou cliquez
sur \sphinxstylestrong{Ajouter}.

\item {} 
\sphinxAtStartPar
Pour supprimer un mot\sphinxhyphen{}clé, cliquez sur le \sphinxstylestrong{×} à côté de lui.

\end{enumerate}


\subsubsection{Catégorie}
\label{\detokenize{fr/03_gerer_metadonnees:categorie}}
\sphinxAtStartPar
La \sphinxstylestrong{catégorie} attribue la ressource à une classification thématique. Sélectionnez
la catégorie la plus appropriée dans le menu déroulant :
\begin{itemize}
\item {} 
\sphinxAtStartPar
Pêche et Aquaculture

\item {} 
\sphinxAtStartPar
Biodiversité Marine

\item {} 
\sphinxAtStartPar
Pollution

\item {} 
\sphinxAtStartPar
Changements Climatiques

\item {} 
\sphinxAtStartPar
Planification Spatiale Marine

\item {} 
\sphinxAtStartPar
Durabilité et Économie Bleue

\item {} 
\sphinxAtStartPar
Gouvernance

\end{itemize}

\sphinxAtStartPar
Une ressource ne peut appartenir qu’à une seule catégorie principale. Utilisez les
mots\sphinxhyphen{}clés pour l’étiquetage thématique supplémentaire.


\subsubsection{Emprise Spatiale}
\label{\detokenize{fr/03_gerer_metadonnees:emprise-spatiale}}
\sphinxAtStartPar
L’\sphinxstylestrong{emprise spatiale} définit le rectangle d’emprise géographique de la ressource.
Pour les datasets, elle est généralement définie automatiquement lors du téléversement.
Vous pouvez l’affiner manuellement :
\begin{enumerate}
\sphinxsetlistlabels{\arabic}{enumi}{enumii}{}{.}%
\item {} 
\sphinxAtStartPar
Allez dans l’onglet \sphinxstylestrong{Location and licences}.

\item {} 
\sphinxAtStartPar
Dans la section \sphinxstyleemphasis{Emprise spatiale}, utilisez le widget carte pour dessiner ou
ajuster le rectangle, ou saisissez les coordonnées (longitude et latitude min/max)
manuellement.

\end{enumerate}

\begin{sphinxadmonition}{tip}{Tip:}
\sphinxAtStartPar
Pour les ressources liées à Pelagos, l’emprise spatiale devrait couvrir au moins
la zone du Sanctuaire Pelagos (approximativement entre la France, l’Italie et Monaco
dans la mer Ligure et la Méditerranée occidentale adjacente).
\end{sphinxadmonition}


\subsubsection{Champs de Date}
\label{\detokenize{fr/03_gerer_metadonnees:champs-de-date}}
\sphinxAtStartPar
Trois champs de date sont disponibles :


\begin{savenotes}\sphinxattablestart
\sphinxthistablewithglobalstyle
\centering
\begin{tabular}[t]{\X{30}{100}\X{70}{100}}
\sphinxtoprule
\sphinxstyletheadfamily 
\sphinxAtStartPar
Champ
&\sphinxstyletheadfamily 
\sphinxAtStartPar
Description
\\
\sphinxmidrule
\sphinxtableatstartofbodyhook
\sphinxAtStartPar
\sphinxstylestrong{Date} (date de création)
&
\sphinxAtStartPar
Quand la ressource a été créée à l’origine ou quand les données ont été collectées.
\\
\sphinxhline
\sphinxAtStartPar
\sphinxstylestrong{Date de publication}
&
\sphinxAtStartPar
Quand la ressource a été publiée sur KMaP.
\\
\sphinxhline
\sphinxAtStartPar
\sphinxstylestrong{Date de révision}
&
\sphinxAtStartPar
Quand la ressource a été mise à jour pour la dernière fois.
\\
\sphinxbottomrule
\end{tabular}
\sphinxtableafterendhook\par
\sphinxattableend\end{savenotes}

\sphinxAtStartPar
Gardez ces dates précises — elles servent à filtrer et trier les ressources dans
le catalogue.


\subsubsection{Licence}
\label{\detokenize{fr/03_gerer_metadonnees:licence}}
\sphinxAtStartPar
La \sphinxstylestrong{licence} définit comment d’autres peuvent utiliser la ressource.
\begin{enumerate}
\sphinxsetlistlabels{\arabic}{enumi}{enumii}{}{.}%
\item {} 
\sphinxAtStartPar
Allez dans l’onglet \sphinxstylestrong{Location and licences}.

\item {} 
\sphinxAtStartPar
Sélectionnez une licence dans le menu déroulant. Les options courantes comprennent :
\begin{itemize}
\item {} 
\sphinxAtStartPar
\sphinxstyleemphasis{Creative Commons Attribution (CC BY)} \textendash{} réutilisation libre avec attribution.

\item {} 
\sphinxAtStartPar
\sphinxstyleemphasis{Creative Commons Attribution\sphinxhyphen{}Partage dans les Mêmes Conditions (CC BY\sphinxhyphen{}SA)} \textendash{}
réutilisation avec attribution et même licence.

\item {} 
\sphinxAtStartPar
\sphinxstyleemphasis{Open Database Licence (ODbL)} \textendash{} pour les datasets.

\item {} 
\sphinxAtStartPar
\sphinxstyleemphasis{Aucune restriction} \textendash{} domaine public.

\end{itemize}

\end{enumerate}

\begin{sphinxadmonition}{note}{Note:}
\sphinxAtStartPar
Si vous n’êtes pas sûr de la licence à appliquer, consultez votre administrateur
de groupe. Appliquer une licence incorrecte peut restreindre la réutilisation de
vos données de manière inappropriée.
\end{sphinxadmonition}


\subsubsection{Point de Contact}
\label{\detokenize{fr/03_gerer_metadonnees:point-de-contact}}
\sphinxAtStartPar
Le \sphinxstylestrong{point de contact} identifie qui contacter pour des questions sur la ressource.
Il peut s’agir d’une personne ou d’une organisation.
\begin{enumerate}
\sphinxsetlistlabels{\arabic}{enumi}{enumii}{}{.}%
\item {} 
\sphinxAtStartPar
Dans l’onglet \sphinxstyleemphasis{Basic info}, trouvez le champ \sphinxstylestrong{Point de contact}.

\item {} 
\sphinxAtStartPar
Commencez à taper un nom pour rechercher parmi les utilisateurs KMaP enregistrés,
ou saisissez le nom d’une organisation.

\end{enumerate}


\subsubsection{Miniature}
\label{\detokenize{fr/03_gerer_metadonnees:miniature}}
\sphinxAtStartPar
La \sphinxstylestrong{miniature} est l’image d’aperçu affichée dans les cartes du catalogue. Elle est
générée automatiquement pour les Datasets et les Cartes, mais peut être remplacée
manuellement :
\begin{enumerate}
\sphinxsetlistlabels{\arabic}{enumi}{enumii}{}{.}%
\item {} 
\sphinxAtStartPar
Cliquez sur \sphinxstylestrong{Edit} dans la page de détail.

\item {} 
\sphinxAtStartPar
Sélectionnez \sphinxstylestrong{Edit thumbnail}.

\item {} 
\sphinxAtStartPar
Téléversez une nouvelle image (recommandé : 600 × 400 px, JPEG ou PNG).

\end{enumerate}


\subsection{Enregistrer les Modifications}
\label{\detokenize{fr/03_gerer_metadonnees:enregistrer-les-modifications}}
\sphinxAtStartPar
Lorsque vous avez terminé de modifier, faites défiler jusqu’au bas de l’éditeur et
cliquez sur \sphinxstylestrong{Enregistrer} (ou \sphinxstylestrong{Mettre à jour} — l’étiquette peut varier). Un
message de confirmation apparaîtra.

\begin{sphinxadmonition}{warning}{Warning:}
\sphinxAtStartPar
Les modifications des métadonnées sont appliquées immédiatement et sont visibles
par tous les utilisateurs ayant accès à la ressource. Il n’existe pas de mode
brouillon ou de préproduction.
\end{sphinxadmonition}

\sphinxAtStartPar
Après l’enregistrement, vous êtes redirigé vers la page de détail de la ressource
où vous pouvez vérifier que tous les champs ont été correctement mis à jour.


\subsection{Liste de Contrôle de la Complétude des Métadonnées}
\label{\detokenize{fr/03_gerer_metadonnees:liste-de-controle-de-la-completude-des-metadonnees}}
\sphinxAtStartPar
Utilisez la liste de contrôle suivante avant de publier ou partager une ressource :


\begin{savenotes}\sphinxattablestart
\sphinxthistablewithglobalstyle
\centering
\begin{tabular}[t]{\X{10}{100}\X{90}{100}}
\sphinxtoprule
\sphinxstyletheadfamily 
\sphinxAtStartPar
\(\checkmark\)
&\sphinxstyletheadfamily 
\sphinxAtStartPar
Élément
\\
\sphinxmidrule
\sphinxtableatstartofbodyhook
\sphinxAtStartPar
☐
&
\sphinxAtStartPar
Le titre est descriptif et spécifique.
\\
\sphinxhline
\sphinxAtStartPar
☐
&
\sphinxAtStartPar
Le résumé explique quoi, où, quand et pourquoi.
\\
\sphinxhline
\sphinxAtStartPar
☐
&
\sphinxAtStartPar
Au moins un mot\sphinxhyphen{}clé Pelagos est présent (\sphinxcode{\sphinxupquote{pelagos sanctuary}} ou
\sphinxcode{\sphinxupquote{pelagos agreement}}).
\\
\sphinxhline
\sphinxAtStartPar
☐
&
\sphinxAtStartPar
La catégorie est définie.
\\
\sphinxhline
\sphinxAtStartPar
☐
&
\sphinxAtStartPar
L’emprise spatiale est définie et couvre la bonne zone.
\\
\sphinxhline
\sphinxAtStartPar
☐
&
\sphinxAtStartPar
Les champs de date sont remplis.
\\
\sphinxhline
\sphinxAtStartPar
☐
&
\sphinxAtStartPar
La licence est définie.
\\
\sphinxhline
\sphinxAtStartPar
☐
&
\sphinxAtStartPar
Le point de contact est identifié.
\\
\sphinxhline
\sphinxAtStartPar
☐
&
\sphinxAtStartPar
La miniature est représentative.
\\
\sphinxbottomrule
\end{tabular}
\sphinxtableafterendhook\par
\sphinxattableend\end{savenotes}

\sphinxstepscope


\section{Explorer les Cartes}
\label{\detokenize{fr/04_explorer_cartes:explorer-les-cartes}}\label{\detokenize{fr/04_explorer_cartes:maps-fr}}\label{\detokenize{fr/04_explorer_cartes::doc}}
\begin{sphinxcontents}
\sphinxstylecontentstitle{Sur cette page}
\begin{itemize}
\item {} 
\sphinxAtStartPar
\phantomsection\label{\detokenize{fr/04_explorer_cartes:id1}}{\hyperref[\detokenize{fr/04_explorer_cartes:trouver-une-carte}]{\sphinxcrossref{Trouver une Carte}}}

\item {} 
\sphinxAtStartPar
\phantomsection\label{\detokenize{fr/04_explorer_cartes:id2}}{\hyperref[\detokenize{fr/04_explorer_cartes:l-interface-du-visualiseur-de-carte}]{\sphinxcrossref{L’Interface du Visualiseur de Carte}}}

\item {} 
\sphinxAtStartPar
\phantomsection\label{\detokenize{fr/04_explorer_cartes:id3}}{\hyperref[\detokenize{fr/04_explorer_cartes:naviguer-dans-la-carte}]{\sphinxcrossref{Naviguer dans la Carte}}}
\begin{itemize}
\item {} 
\sphinxAtStartPar
\phantomsection\label{\detokenize{fr/04_explorer_cartes:id4}}{\hyperref[\detokenize{fr/04_explorer_cartes:navigation-de-base}]{\sphinxcrossref{Navigation de Base}}}

\item {} 
\sphinxAtStartPar
\phantomsection\label{\detokenize{fr/04_explorer_cartes:id5}}{\hyperref[\detokenize{fr/04_explorer_cartes:aller-a-un-emplacement-specifique}]{\sphinxcrossref{Aller à un Emplacement Spécifique}}}

\end{itemize}

\item {} 
\sphinxAtStartPar
\phantomsection\label{\detokenize{fr/04_explorer_cartes:id6}}{\hyperref[\detokenize{fr/04_explorer_cartes:travailler-avec-les-couches}]{\sphinxcrossref{Travailler avec les Couches}}}
\begin{itemize}
\item {} 
\sphinxAtStartPar
\phantomsection\label{\detokenize{fr/04_explorer_cartes:id7}}{\hyperref[\detokenize{fr/04_explorer_cartes:ouvrir-le-panneau-de-couches}]{\sphinxcrossref{Ouvrir le Panneau de Couches}}}

\item {} 
\sphinxAtStartPar
\phantomsection\label{\detokenize{fr/04_explorer_cartes:id8}}{\hyperref[\detokenize{fr/04_explorer_cartes:basculer-la-visibilite-des-couches}]{\sphinxcrossref{Basculer la Visibilité des Couches}}}

\item {} 
\sphinxAtStartPar
\phantomsection\label{\detokenize{fr/04_explorer_cartes:id9}}{\hyperref[\detokenize{fr/04_explorer_cartes:changer-l-ordre-des-couches}]{\sphinxcrossref{Changer l’Ordre des Couches}}}

\item {} 
\sphinxAtStartPar
\phantomsection\label{\detokenize{fr/04_explorer_cartes:id10}}{\hyperref[\detokenize{fr/04_explorer_cartes:ajuster-l-opacite-d-une-couche}]{\sphinxcrossref{Ajuster l’Opacité d’une Couche}}}

\item {} 
\sphinxAtStartPar
\phantomsection\label{\detokenize{fr/04_explorer_cartes:id11}}{\hyperref[\detokenize{fr/04_explorer_cartes:afficher-la-legende}]{\sphinxcrossref{Afficher la Légende}}}

\end{itemize}

\item {} 
\sphinxAtStartPar
\phantomsection\label{\detokenize{fr/04_explorer_cartes:id12}}{\hyperref[\detokenize{fr/04_explorer_cartes:interroger-les-elements}]{\sphinxcrossref{Interroger les Éléments}}}
\begin{itemize}
\item {} 
\sphinxAtStartPar
\phantomsection\label{\detokenize{fr/04_explorer_cartes:id13}}{\hyperref[\detokenize{fr/04_explorer_cartes:cliquer-sur-un-element}]{\sphinxcrossref{Cliquer sur un Élément}}}

\item {} 
\sphinxAtStartPar
\phantomsection\label{\detokenize{fr/04_explorer_cartes:id14}}{\hyperref[\detokenize{fr/04_explorer_cartes:outil-identifier}]{\sphinxcrossref{Outil Identifier}}}

\item {} 
\sphinxAtStartPar
\phantomsection\label{\detokenize{fr/04_explorer_cartes:id15}}{\hyperref[\detokenize{fr/04_explorer_cartes:selectionner-et-filtrer-les-elements}]{\sphinxcrossref{Sélectionner et Filtrer les Éléments}}}

\end{itemize}

\item {} 
\sphinxAtStartPar
\phantomsection\label{\detokenize{fr/04_explorer_cartes:id16}}{\hyperref[\detokenize{fr/04_explorer_cartes:outils-de-mesure}]{\sphinxcrossref{Outils de Mesure}}}
\begin{itemize}
\item {} 
\sphinxAtStartPar
\phantomsection\label{\detokenize{fr/04_explorer_cartes:id17}}{\hyperref[\detokenize{fr/04_explorer_cartes:mesure-de-distance}]{\sphinxcrossref{Mesure de Distance}}}

\item {} 
\sphinxAtStartPar
\phantomsection\label{\detokenize{fr/04_explorer_cartes:id18}}{\hyperref[\detokenize{fr/04_explorer_cartes:mesure-de-surface}]{\sphinxcrossref{Mesure de Surface}}}

\end{itemize}

\item {} 
\sphinxAtStartPar
\phantomsection\label{\detokenize{fr/04_explorer_cartes:id19}}{\hyperref[\detokenize{fr/04_explorer_cartes:imprimer-et-exporter-la-vue-de-la-carte}]{\sphinxcrossref{Imprimer et Exporter la Vue de la Carte}}}
\begin{itemize}
\item {} 
\sphinxAtStartPar
\phantomsection\label{\detokenize{fr/04_explorer_cartes:id20}}{\hyperref[\detokenize{fr/04_explorer_cartes:imprimer}]{\sphinxcrossref{Imprimer}}}

\item {} 
\sphinxAtStartPar
\phantomsection\label{\detokenize{fr/04_explorer_cartes:id21}}{\hyperref[\detokenize{fr/04_explorer_cartes:partager-permalien}]{\sphinxcrossref{Partager / Permalien}}}

\item {} 
\sphinxAtStartPar
\phantomsection\label{\detokenize{fr/04_explorer_cartes:id22}}{\hyperref[\detokenize{fr/04_explorer_cartes:integrer-embed}]{\sphinxcrossref{Intégrer (Embed)}}}

\end{itemize}

\item {} 
\sphinxAtStartPar
\phantomsection\label{\detokenize{fr/04_explorer_cartes:id23}}{\hyperref[\detokenize{fr/04_explorer_cartes:enregistrer-les-modifications-d-une-carte}]{\sphinxcrossref{Enregistrer les Modifications d’une Carte}}}

\item {} 
\sphinxAtStartPar
\phantomsection\label{\detokenize{fr/04_explorer_cartes:id24}}{\hyperref[\detokenize{fr/04_explorer_cartes:creer-une-nouvelle-carte}]{\sphinxcrossref{Créer une Nouvelle Carte}}}

\end{itemize}
\end{sphinxcontents}

\sphinxAtStartPar
Une \sphinxstylestrong{Carte} dans KMaP est une carte web interactive à plusieurs couches, construite
à partir d’un ou plusieurs Datasets. Cette section explique comment trouver, ouvrir et
naviguer dans les Cartes, comment travailler avec les couches, et comment extraire des
informations des données affichées.


\subsection{Trouver une Carte}
\label{\detokenize{fr/04_explorer_cartes:trouver-une-carte}}
\sphinxAtStartPar
Les Cartes sont répertoriées dans le catalogue aux côtés de tous les autres types de
ressources. Pour filtrer le catalogue afin d’afficher uniquement les Cartes :
\begin{enumerate}
\sphinxsetlistlabels{\arabic}{enumi}{enumii}{}{.}%
\item {} 
\sphinxAtStartPar
Ouvrez le catalogue : cliquez sur \sphinxstylestrong{Maps} dans la barre de navigation en haut.

\item {} 
\sphinxAtStartPar
Dans le panneau de filtres à gauche, sous \sphinxstylestrong{Type de ressource}, sélectionnez \sphinxstylestrong{Map}.

\item {} 
\sphinxAtStartPar
Pour trouver des cartes liées à Pelagos, appliquez également le filtre par mot\sphinxhyphen{}clé :
dans le champ \sphinxstylestrong{Mots\sphinxhyphen{}clés}, tapez \sphinxcode{\sphinxupquote{pelagos sanctuary}} ou \sphinxcode{\sphinxupquote{pelagos agreement}}.

\end{enumerate}

\sphinxAtStartPar
Alternativement, recherchez une carte par titre en utilisant la barre de recherche.

\sphinxAtStartPar
Cliquez sur une carte de ressource pour ouvrir sa \sphinxstylestrong{page de détail}. De là, cliquez
sur \sphinxstylestrong{View map} pour l’ouvrir dans le visualiseur de carte interactif.


\subsection{L’Interface du Visualiseur de Carte}
\label{\detokenize{fr/04_explorer_cartes:l-interface-du-visualiseur-de-carte}}
\sphinxAtStartPar
Le visualiseur de carte (basé sur \sphinxstylestrong{MapStore}) s’ouvre en plein écran. Ses principaux
éléments d’interface sont :


\begin{savenotes}\sphinxattablestart
\sphinxthistablewithglobalstyle
\centering
\begin{tabular}[t]{\X{30}{100}\X{70}{100}}
\sphinxtoprule
\sphinxstyletheadfamily 
\sphinxAtStartPar
Élément
&\sphinxstyletheadfamily 
\sphinxAtStartPar
Description
\\
\sphinxmidrule
\sphinxtableatstartofbodyhook
\sphinxAtStartPar
\sphinxstylestrong{Zone de carte principale}
&
\sphinxAtStartPar
La zone de carte interactive centrale.
\\
\sphinxhline
\sphinxAtStartPar
\sphinxstylestrong{Panneau de couches} (☰)
&
\sphinxAtStartPar
Bascule la visibilité des couches, les réordonne, accède aux paramètres.
\\
\sphinxhline
\sphinxAtStartPar
\sphinxstylestrong{Sélecteur de fond de carte}
&
\sphinxAtStartPar
Change le fond de carte (OpenStreetMap, imagerie satellitaire, …).
\\
\sphinxhline
\sphinxAtStartPar
\sphinxstylestrong{Contrôles de zoom} (+ / −)
&
\sphinxAtStartPar
Zoom avant et arrière. Prend également en charge la molette de la souris et le
geste de pincement.
\\
\sphinxhline
\sphinxAtStartPar
\sphinxstylestrong{Barre d’outils}
&
\sphinxAtStartPar
Rangée d’icônes d’outils en haut ou sur le côté du canvas.
\\
\sphinxhline
\sphinxAtStartPar
\sphinxstylestrong{Barre d’échelle}
&
\sphinxAtStartPar
Affiche l’échelle actuelle de la carte.
\\
\sphinxhline
\sphinxAtStartPar
\sphinxstylestrong{Panneau de coordonnées}
&
\sphinxAtStartPar
Affiche les coordonnées géographiques de la position du curseur.
\\
\sphinxbottomrule
\end{tabular}
\sphinxtableafterendhook\par
\sphinxattableend\end{savenotes}


\subsection{Naviguer dans la Carte}
\label{\detokenize{fr/04_explorer_cartes:naviguer-dans-la-carte}}

\subsubsection{Navigation de Base}
\label{\detokenize{fr/04_explorer_cartes:navigation-de-base}}

\begin{savenotes}\sphinxattablestart
\sphinxthistablewithglobalstyle
\centering
\begin{tabular}[t]{\X{35}{100}\X{65}{100}}
\sphinxtoprule
\sphinxstyletheadfamily 
\sphinxAtStartPar
Action
&\sphinxstyletheadfamily 
\sphinxAtStartPar
Comment l’effectuer
\\
\sphinxmidrule
\sphinxtableatstartofbodyhook
\sphinxAtStartPar
\sphinxstylestrong{Déplacer la carte (pan)}
&
\sphinxAtStartPar
Cliquez et faites glisser sur le canvas de la carte.
\\
\sphinxhline
\sphinxAtStartPar
\sphinxstylestrong{Zoom avant}
&
\sphinxAtStartPar
Faites défiler la molette de la souris vers l’avant, cliquez sur \sphinxstylestrong{+} ou
double\sphinxhyphen{}cliquez sur la carte.
\\
\sphinxhline
\sphinxAtStartPar
\sphinxstylestrong{Zoom arrière}
&
\sphinxAtStartPar
Faites défiler la molette de la souris vers l’arrière ou cliquez sur \sphinxstylestrong{−}.
\\
\sphinxhline
\sphinxAtStartPar
\sphinxstylestrong{Revenir à l’étendue complète}
&
\sphinxAtStartPar
Cliquez sur l’icône \sphinxstyleemphasis{Accueil} (🏠) dans la barre d’outils pour revenir à la
vue par défaut.
\\
\sphinxhline
\sphinxAtStartPar
\sphinxstylestrong{Zoom sur une couche}
&
\sphinxAtStartPar
Dans le panneau de couches, faites un clic droit sur une couche et sélectionnez
\sphinxstyleemphasis{Zoomer sur l’étendue de la couche}.
\\
\sphinxbottomrule
\end{tabular}
\sphinxtableafterendhook\par
\sphinxattableend\end{savenotes}


\subsubsection{Aller à un Emplacement Spécifique}
\label{\detokenize{fr/04_explorer_cartes:aller-a-un-emplacement-specifique}}
\sphinxAtStartPar
Utilisez l’outil \sphinxstylestrong{Recherche / Géocodeur} dans la barre d’outils (icône 🔍) pour
rechercher un nom de lieu. Tapez le nom et appuyez sur Entrée — la carte se déplacera
vers cet emplacement.


\subsection{Travailler avec les Couches}
\label{\detokenize{fr/04_explorer_cartes:travailler-avec-les-couches}}

\subsubsection{Ouvrir le Panneau de Couches}
\label{\detokenize{fr/04_explorer_cartes:ouvrir-le-panneau-de-couches}}
\sphinxAtStartPar
Cliquez sur l’icône \sphinxstylestrong{☰} (hamburger) pour ouvrir le panneau de couches. Il liste
toutes les couches incluses dans la carte, regroupées en :
\begin{itemize}
\item {} 
\sphinxAtStartPar
\sphinxstylestrong{Superpositions (Overlays)} \textendash{} les couches de données thématiques (ex. limite du
Sanctuaire Pelagos, points d’observation de cétacés, routes maritimes).

\item {} 
\sphinxAtStartPar
\sphinxstylestrong{Couches de fond (Background)} \textendash{} les options de fond de carte (une seule peut être
active à la fois).

\end{itemize}


\subsubsection{Basculer la Visibilité des Couches}
\label{\detokenize{fr/04_explorer_cartes:basculer-la-visibilite-des-couches}}
\sphinxAtStartPar
Chaque couche dans le panneau possède une \sphinxstylestrong{icône œil} (👁). Cliquez dessus pour
afficher ou masquer cette couche. Les couches masquées restent dans la composition de
la carte mais ne sont pas affichées sur le canvas.


\subsubsection{Changer l’Ordre des Couches}
\label{\detokenize{fr/04_explorer_cartes:changer-l-ordre-des-couches}}
\sphinxAtStartPar
Les couches sont dessinées dans l’ordre où elles apparaissent dans le panneau (les
couches du haut sont dessinées par\sphinxhyphen{}dessus). Pour réordonner :
\begin{enumerate}
\sphinxsetlistlabels{\arabic}{enumi}{enumii}{}{.}%
\item {} 
\sphinxAtStartPar
Cliquez et maintenez un nom de couche dans le panneau.

\item {} 
\sphinxAtStartPar
Faites\sphinxhyphen{}le glisser vers le haut ou le bas jusqu’à la position souhaitée.

\item {} 
\sphinxAtStartPar
Relâchez pour le déposer.

\end{enumerate}


\subsubsection{Ajuster l’Opacité d’une Couche}
\label{\detokenize{fr/04_explorer_cartes:ajuster-l-opacite-d-une-couche}}\begin{enumerate}
\sphinxsetlistlabels{\arabic}{enumi}{enumii}{}{.}%
\item {} 
\sphinxAtStartPar
Dans le panneau de couches, cliquez sur l’icône \sphinxstylestrong{⋮} (options) à côté d’une couche.

\item {} 
\sphinxAtStartPar
Sélectionnez \sphinxstylestrong{Paramètres} (ou l’icône de curseur).

\item {} 
\sphinxAtStartPar
Utilisez le curseur \sphinxstylestrong{Opacité} pour régler la transparence entre 0 \% (invisible)
et 100 \% (plein).

\end{enumerate}


\subsubsection{Afficher la Légende}
\label{\detokenize{fr/04_explorer_cartes:afficher-la-legende}}\begin{enumerate}
\sphinxsetlistlabels{\arabic}{enumi}{enumii}{}{.}%
\item {} 
\sphinxAtStartPar
Cliquez sur l’icône \sphinxstylestrong{⋮} (options) à côté d’une couche.

\item {} 
\sphinxAtStartPar
Sélectionnez \sphinxstylestrong{Légende} pour développer la clé des symboles de cette couche.

\end{enumerate}

\sphinxAtStartPar
Certaines configurations de carte affichent également un panneau de légende persistant
sur le côté du canvas.


\subsection{Interroger les Éléments}
\label{\detokenize{fr/04_explorer_cartes:interroger-les-elements}}

\subsubsection{Cliquer sur un Élément}
\label{\detokenize{fr/04_explorer_cartes:cliquer-sur-un-element}}
\sphinxAtStartPar
Cliquez sur tout élément visible (point, ligne ou polygone) sur le canvas de la carte
pour ouvrir une \sphinxstylestrong{fenêtre contextuelle d’attributs}. Celle\sphinxhyphen{}ci affiche les valeurs
d’attributs associées à cet élément — par exemple, le nom d’une aire marine protégée,
un nom d’espèce ou une valeur de surveillance.

\sphinxAtStartPar
Si plusieurs couches se superposent à l’emplacement cliqué, la fenêtre contextuelle
peut afficher des résultats de toutes les couches superposées. Utilisez les boutons
fléchés dans la fenêtre pour naviguer entre eux.


\subsubsection{Outil Identifier}
\label{\detokenize{fr/04_explorer_cartes:outil-identifier}}
\sphinxAtStartPar
L’outil \sphinxstylestrong{Identifier} (icône ℹ️ dans la barre d’outils) fonctionne de manière
similaire au clic direct sur la carte. Il interroge toutes les couches visibles au
point cliqué et renvoie leurs attributs dans un panneau.


\subsubsection{Sélectionner et Filtrer les Éléments}
\label{\detokenize{fr/04_explorer_cartes:selectionner-et-filtrer-les-elements}}
\sphinxAtStartPar
Pour les couches vectorielles, MapStore propose un outil \sphinxstylestrong{Filtrer} :
\begin{enumerate}
\sphinxsetlistlabels{\arabic}{enumi}{enumii}{}{.}%
\item {} 
\sphinxAtStartPar
Dans le panneau de couches, cliquez sur \sphinxstylestrong{⋮} à côté de la couche souhaitée.

\item {} 
\sphinxAtStartPar
Sélectionnez \sphinxstylestrong{Filtrer}.

\item {} 
\sphinxAtStartPar
Définissez une condition de filtre en utilisant les champs d’attributs
(ex. afficher uniquement les enregistrements où \sphinxstyleemphasis{pays = “France”}).

\item {} 
\sphinxAtStartPar
Cliquez sur \sphinxstylestrong{Appliquer}. Seuls les éléments correspondant au filtre sont affichés.

\end{enumerate}

\begin{sphinxadmonition}{tip}{Tip:}
\sphinxAtStartPar
Utilisez l’outil de filtre pour vous concentrer sur un sous\sphinxhyphen{}ensemble spécifique de
données Pelagos — par exemple, en affichant uniquement les observations de cétacés
d’une année spécifique, ou uniquement les routes traversant le sanctuaire.
\end{sphinxadmonition}


\subsection{Outils de Mesure}
\label{\detokenize{fr/04_explorer_cartes:outils-de-mesure}}

\subsubsection{Mesure de Distance}
\label{\detokenize{fr/04_explorer_cartes:mesure-de-distance}}\begin{enumerate}
\sphinxsetlistlabels{\arabic}{enumi}{enumii}{}{.}%
\item {} 
\sphinxAtStartPar
Cliquez sur l’outil \sphinxstylestrong{Mesure} (icône 📐) dans la barre d’outils.

\item {} 
\sphinxAtStartPar
Sélectionnez \sphinxstylestrong{Distance}.

\item {} 
\sphinxAtStartPar
Cliquez des points sur la carte pour définir une ligne. Double\sphinxhyphen{}cliquez pour terminer.

\item {} 
\sphinxAtStartPar
La distance totale est affichée dans l’unité de votre choix (km, miles, nm).

\end{enumerate}


\subsubsection{Mesure de Surface}
\label{\detokenize{fr/04_explorer_cartes:mesure-de-surface}}\begin{enumerate}
\sphinxsetlistlabels{\arabic}{enumi}{enumii}{}{.}%
\item {} 
\sphinxAtStartPar
Cliquez sur l’outil \sphinxstylestrong{Mesure}.

\item {} 
\sphinxAtStartPar
Sélectionnez \sphinxstylestrong{Surface}.

\item {} 
\sphinxAtStartPar
Cliquez pour définir les sommets d’un polygone. Double\sphinxhyphen{}cliquez pour fermer.

\item {} 
\sphinxAtStartPar
La surface calculée est affichée.

\end{enumerate}

\begin{sphinxadmonition}{note}{Note:}
\sphinxAtStartPar
Les mesures sont effectuées sur la projection de la carte et peuvent présenter de
légères imprécisions géométriques pour les grandes surfaces. Pour des mesures
précises, utilisez une application SIG de bureau.
\end{sphinxadmonition}


\subsection{Imprimer et Exporter la Vue de la Carte}
\label{\detokenize{fr/04_explorer_cartes:imprimer-et-exporter-la-vue-de-la-carte}}

\subsubsection{Imprimer}
\label{\detokenize{fr/04_explorer_cartes:imprimer}}\begin{enumerate}
\sphinxsetlistlabels{\arabic}{enumi}{enumii}{}{.}%
\item {} 
\sphinxAtStartPar
Cliquez sur l’outil \sphinxstylestrong{Imprimer} (icône 🖨) dans la barre d’outils.

\item {} 
\sphinxAtStartPar
Une boîte de dialogue d’impression apparaît. Configurez :
\begin{itemize}
\item {} 
\sphinxAtStartPar
\sphinxstylestrong{Titre} \textendash{} un titre à afficher sur la carte imprimée.

\item {} 
\sphinxAtStartPar
\sphinxstylestrong{Format de papier} \textendash{} A4, A3, Letter, etc.

\item {} 
\sphinxAtStartPar
\sphinxstylestrong{Résolution} \textendash{} 72, 150 ou 300 dpi.

\item {} 
\sphinxAtStartPar
\sphinxstylestrong{Mise en page} \textendash{} portrait ou paysage.

\end{itemize}

\item {} 
\sphinxAtStartPar
Cliquez sur \sphinxstylestrong{Imprimer} pour générer un PDF ou une image.

\end{enumerate}


\subsubsection{Partager / Permalien}
\label{\detokenize{fr/04_explorer_cartes:partager-permalien}}
\sphinxAtStartPar
Pour partager la vue actuelle de la carte (y compris le niveau de zoom, les couches
actives et le point central) :
\begin{enumerate}
\sphinxsetlistlabels{\arabic}{enumi}{enumii}{}{.}%
\item {} 
\sphinxAtStartPar
Cliquez sur le bouton \sphinxstylestrong{Partager} (icône de lien) dans la barre d’outils.

\item {} 
\sphinxAtStartPar
Un URL permanent est généré. Copiez et partagez cet URL — les destinataires
ouvriront la carte dans le même état.

\end{enumerate}


\subsubsection{Intégrer (Embed)}
\label{\detokenize{fr/04_explorer_cartes:integrer-embed}}\begin{enumerate}
\sphinxsetlistlabels{\arabic}{enumi}{enumii}{}{.}%
\item {} 
\sphinxAtStartPar
Cliquez sur le bouton \sphinxstylestrong{Partager}.

\item {} 
\sphinxAtStartPar
Passez à l’onglet \sphinxstylestrong{Intégrer}.

\item {} 
\sphinxAtStartPar
Copiez le code HTML \sphinxcode{\sphinxupquote{\textless{}iframe\textgreater{}}} et collez\sphinxhyphen{}le dans votre site web ou rapport.

\end{enumerate}


\subsection{Enregistrer les Modifications d’une Carte}
\label{\detokenize{fr/04_explorer_cartes:enregistrer-les-modifications-d-une-carte}}
\sphinxAtStartPar
Si vous avez la \sphinxstylestrong{permission de modification} sur une Carte, vous pouvez enregistrer
les modifications apportées aux couches, aux styles et à la vue initiale :
\begin{enumerate}
\sphinxsetlistlabels{\arabic}{enumi}{enumii}{}{.}%
\item {} 
\sphinxAtStartPar
Effectuez vos modifications dans le visualiseur de carte (réordonnez les couches,
modifiez les styles, définissez une nouvelle étendue initiale, etc.).

\item {} 
\sphinxAtStartPar
Cliquez sur le bouton \sphinxstylestrong{Enregistrer} (icône 💾) dans la barre d’outils.

\item {} 
\sphinxAtStartPar
Choisissez \sphinxstylestrong{Enregistrer} (écraser la carte existante) ou \sphinxstylestrong{Enregistrer sous}
(créer une nouvelle carte).

\end{enumerate}

\begin{sphinxadmonition}{warning}{Warning:}
\sphinxAtStartPar
L’enregistrement écrase la carte pour tous les utilisateurs qui y ont accès. Si vous
souhaitez expérimenter sans affecter les autres, utilisez d’abord \sphinxstylestrong{Enregistrer sous}
pour créer une copie personnelle.
\end{sphinxadmonition}


\subsection{Créer une Nouvelle Carte}
\label{\detokenize{fr/04_explorer_cartes:creer-une-nouvelle-carte}}
\sphinxAtStartPar
Les membres du groupe Accord Pelagos disposant des autorisations appropriées peuvent
créer de nouvelles Cartes :
\begin{enumerate}
\sphinxsetlistlabels{\arabic}{enumi}{enumii}{}{.}%
\item {} 
\sphinxAtStartPar
Allez dans \sphinxstylestrong{Maps} dans la barre de navigation.

\item {} 
\sphinxAtStartPar
Cliquez sur \sphinxstylestrong{+ Créer une carte} (ou le bouton équivalent en haut du catalogue).

\item {} 
\sphinxAtStartPar
Le visualiseur de carte s’ouvre avec un canvas vide.

\item {} 
\sphinxAtStartPar
Utilisez le bouton \sphinxstylestrong{Ajouter une couche} (icône + dans le panneau de couches) pour
rechercher et ajouter des Datasets depuis le catalogue KMaP.

\item {} 
\sphinxAtStartPar
Stylisez chaque couche selon vos besoins.

\item {} 
\sphinxAtStartPar
Cliquez sur \sphinxstylestrong{Enregistrer} et fournissez un titre, un résumé et des mots\sphinxhyphen{}clés
(n’oubliez pas d’inclure \sphinxcode{\sphinxupquote{pelagos sanctuary}} ou \sphinxcode{\sphinxupquote{pelagos agreement}}).

\end{enumerate}



\renewcommand{\indexname}{Index}
\printindex
\end{document}